% EES Latex-Vorlage M�rz 2012
\documentclass[
12pt,               																	% Schriftgr��e
a4paper,            																	% Layout f�r DINA4
german,             																	% deutsche Sprache, global
twoside,            																	% Layout f�r beidseitigen Druck
headinclude,        																	% Kopfzeile wird Seiten-Layouts mit ber�cksichtigt
headsepline,        																	% horizontale Linie unter Kolumnentitel
plainheadsepline,																			% horizontale Linie unter Kolumnentitel auch bei Chapter
BCOR20mm,           																	% Korrektur f�r die Bindung
DIV18,              																	% DIV-Wert f�r die Erstellung des Satzspiegels, siehe scrguide
parskip=half,       																	% Absatzabstand statt Absatzeinzug
openany,            																	% Kapitel k�nnen auf geraden und ungeraden Seiten beginnen
bibliography=totoc,version=first, 										% Literaturverz. wird ins Inhaltsverzeichnis eingetragen
numbers=noenddot,   																	% Kapitelnummern immer ohne Punkt
captions=tableheading,version=first, 									% korrekte Abst�nde bei Tabellen�berschriften
fleqn,             																		% fleqn: Glgen links (statt mittig)
listof=totoc,version=first														% Abbildungs- und Tabellenverzeichnis ins Inhaltsverzeichnis
]{scrbook}


%--------------- Packages ----------------
\usepackage[ngerman]{babel} 													% Neue deutsche Trennmuster
\usepackage[latin1]{inputenc} 												% direkte Eingabe von Umlauten & Co.
\usepackage[T1]{fontenc} 															% T1-Schriften
\usepackage{ae}               												% f�r PDF-Erstellung
\usepackage[format=hang,justification=raggedright,singlelinecheck=off,labelfont=bf,font=small,skip=4pt]{caption} % Captions ausrichten
\usepackage[centertags]{amsmath} 											% AMS-Mathematik, centertags zentriert Nummer bei split
\usepackage{tabularx}																	% erweiterte Tabellen
\usepackage{graphicx}            											% zum Einbinden von Grafiken
\usepackage{float}               											% u.a. genaue Plazierung von Gleitobjekten mit H
\usepackage{setspace}            											% Zeilenabstand einstellbar
\usepackage{scrpage2}           											% Kopf- und Fu�zeilen-Layout 
\usepackage{befehle}  																% eigene Befehle in der Datei definiert
\usepackage{listings}																	% Auflistungen
\usepackage{scrhack}																	% Damit kriegt man den Fehler mit den Kopf- und Fu�zeilen weg
\usepackage{makeidx}																	% Paket f�r das Stichwortverzeichnis
\usepackage{xcolor}																		% Paket f�r die Farben
\usepackage[colorlinks,pdfpagelabels,pdfstartview = FitH,bookmarksopen = true,bookmarksnumbered = true,linkcolor = black,plainpages = false,hypertexnames = false,citecolor = black] {hyperref}																			% Paket f�r das verlinkte PDF
\usepackage[absolute]{textpos}
\usepackage[clearempty]{titlesec}				

%%
\usepackage{amssymb}
\usepackage{rotating}
\usepackage{subcaption}
\usepackage{pdflscape}
\usepackage{longtable}
\usepackage{arydshln}
\usepackage{tikz}
\usepackage{pgfplots}
\usepackage{calc}
\usepackage{xparse}
\usetikzlibrary{intersections,decorations.pathreplacing,fit,calc,positioning,plothandlers,plotmarks,patterns}
\newlength{\hatchspread}
\newlength{\hatchthickness}
\newlength{\hatchshift}
\newcommand{\hatchcolor}{}
% declaring the keys in tikz
\tikzset{hatchspread/.code={\setlength{\hatchspread}{#1}},
         hatchthickness/.code={\setlength{\hatchthickness}{#1}},
         hatchshift/.code={\setlength{\hatchshift}{#1}},% must be >= 0
         hatchcolor/.code={\renewcommand{\hatchcolor}{#1}}}
% setting the default values
\tikzset{hatchspread=3pt,
         hatchthickness=0.4pt,
         hatchshift=0pt,% must be >= 0
         hatchcolor=black}
% declaring the pattern
\pgfdeclarepatternformonly[\hatchspread,\hatchthickness,\hatchshift,\hatchcolor]% variables
   {custom north west lines}% name
   {\pgfqpoint{\dimexpr-2\hatchthickness}{\dimexpr-2\hatchthickness}}% lower left corner
   {\pgfqpoint{\dimexpr\hatchspread+2\hatchthickness}{\dimexpr\hatchspread+2\hatchthickness}}% upper right corner
   {\pgfqpoint{\dimexpr\hatchspread}{\dimexpr\hatchspread}}% tile size
   {% shape description
    \pgfsetlinewidth{\hatchthickness}
    \pgfpathmoveto{\pgfqpoint{0pt}{\dimexpr\hatchspread+\hatchshift}}
    \pgfpathlineto{\pgfqpoint{\dimexpr\hatchspread+0.15pt+\hatchshift}{-0.15pt}}
    \ifdim \hatchshift > 0pt
      \pgfpathmoveto{\pgfqpoint{0pt}{\hatchshift}}
      \pgfpathlineto{\pgfqpoint{\dimexpr0.15pt+\hatchshift}{-0.15pt}}
    \fi
    \pgfsetstrokecolor{\hatchcolor}
%    \pgfsetdash{{1pt}{1pt}}{0pt}% dashing cannot work correctly in all situation this way
    \pgfusepath{stroke}
   }

%\usetikzlibrary{decorations.markings}
%\usetikzlibrary{decorations.pathreplacing,decorations.pathmorphing}

%--------------- Sonstiges ----------------

\pagestyle{scrheadings}																% Kopf- und Fu�zeile...
\renewcommand{\headfont}{\normalfont\sffamily}    		% Kolumnentitel serifenlos
\renewcommand{\pnumfont}{\normalfont\sffamily}    		% Seitennummern serifenlos
\ihead[]{\headmark}              											% Kopfzeile innen
\ohead[\pagemark]{\pagemark}     											% Kopfzeile au�en
\ifoot[]{} 																						% Fu�zeile innen
\ofoot[]{}    																				% Fu�zeile au�en	
\setlength{\headheight}{1.5\baselineskip}
\onehalfspacing																				% 1,5 Zeilenabstand
%\typearea[current]{current}        									% Neuberechnung des Satzspiegels mit alten Werten nach �nderung von Zeilenabstand,etc
\renewcommand{\bibname}{Literatur und Quellen} 				% Literaturverzeichnisbezeichnung
\renewcommand{\figurename}{Abb.}   										% Abbildungsbezeichnung
\renewcommand{\listfigurename}{Abbildungsverzeichnis} % Abbildungsverzeichnisbezeichnung
\renewcommand{\captionfont}{\small}										% Bildunterschriften klein kursiv
\graphicspath{{figs/}{bilder/}}    										% Bildverzeichnis


%--------------- Code-Darstellung ----------------
\lstdefinelanguage{Matlab2}
{morekeywords={for,if,else,end,function,disp,clc,input,subplot,plot,figure,axis,ylabel,xlabel,legend,set,gca,gcf,grid,switch,case,otherwise,eval,ss,sisotool,linmod,clear,load,addpath,pwd,},
sensitive=false,
morecomment=[l]{//},
morecomment=[l]{\%},
morecomment=[s]{/*}{*/},
morestring=[b]",
}
\definecolor{MyDarkBlue}{rgb}{0,0.08,0.45}
\definecolor{fett}{rgb}{0.965,0.984,0.961}
\definecolor{OliveGreen}{rgb}{0.43,0.66,0.45}
\lstset{language=Matlab2, basicstyle=\tiny, backgroundcolor=\color{fett}, keywordstyle=\color{MyDarkBlue}\bfseries ,numbers=left, numberstyle=\tiny, tabsize=2, commentstyle=\color{OliveGreen}}
\makeindex


%----------------------- Beginn des Dokuments -----------------------
\begin{document}
\pagenumbering{Roman}																	% R�mische Seitennummerierung
\newcommand{\dd}{\cdot}
\renewcommand{\tt}{(t)}
%\newcommand{\bs}{\boldsymbol}
\renewcommand{\d}{\mathrm{d}}
\renewcommand{\vec}[1]{\bs{\mathrm{#1}}}
\newcommand{\up}[1]{\mathrm{#1}}
\newcommand{\x}{\up{x}}
\renewcommand{\u}{\up{u}}
\newcommand{\C}{\up{C}}
\newcommand{\R}{\up{R}}
\newcommand{\T}{\up{T}}
\newcommand{\f}{\up{f}}
\newcommand{\eh}{\frac{1}{2}}
\renewcommand{\c}{\up{c}}
\renewcommand{\b}{\up{b}}
\renewcommand{\a}{\up{a}}
\newcommand{\s}{\up{s}}
\newcommand{\ul}[1]{\underline{#1}}
\newcommand{\com}[1]{\underline{\up{#1}}}
\newcommand{\comv}[1]{\underline{#1}}

\newcommand{\adj}{\mathrm{adj}}
%%%%%%%%%%%%%%%%%%%%%%%%%%
%TIKZ:
\tikzset{%
	>=latex,
  highlight/.style={rectangle,rounded corners,draw=orange,very thin,fill opacity=0.5,inner sep=-2pt},
  frame/.style={
	rectangle,minimum size=9mm,thick,draw=black
	},
 sum/.style={
	circle,minimum size =3mm,thick,draw=black, inner sep = 0pt
	},
 knot/.style={
	circle,minimum size =2mm,fill=black, inner sep=0pt
	},
 skip  loop/.style={to  path={--  ++(0,#1)  -|  (\tikztotarget)}},
 nonlin/.style = {frame, double,double distance=1pt},
 times/.style = {alias=sourcenode,frame, 
        append after command={
        ($(sourcenode.center)+(-0.2,0.2)$)--($(sourcenode.center)+(0.2,-0.2)$)
	  ($(sourcenode.center)+(-0.2,-0.2)$)--($(sourcenode.center)+(0.2,0.2)$)	
        }
	}
}



\newcommand{\tikzrightbrace}[4][0.5\textwidth]{%
%%% Ein rechter brace der die nodes 2 und 3 umschlie�t.
%%% rechts daneben wird der Text 4 mit der Breite 1 *optional angezeigt!
\tikz[overlay,remember picture]{
\node[fit=(#2.north east) (#3.south east),inner sep =0] (gesamt#2) {};
\draw[decorate, decoration=brace] (gesamt#2.north east)--(gesamt#2.south east);
\path (gesamt#2.east) node[right,xshift=2mm]{\parbox{#1}{#4}};
}
}


\newcommand{\tikzmark}[3][3pt]{\tikz[remember picture,baseline=(#2.base)] \node[inner sep=#1] (#2) {#3};}
%tikzmark: markiert inhalt als tikznode: 1(optional): inner sep der node
%									2: name der Node
%									3: Inhalt der Node
\newcommand{\tikzmarktwo}[3][2pt]{\tikz[remember picture,baseline=(#2.base)] 
\node[inner ysep=#1,inner xsep=0pt, outer xsep=0pt] (#2) {#3};}
%%tikzmark: markiert inhalt als tikznode: 1(optional): inner ysep der node
%%									2(optional): inner xsep der node
%%									3: name der Node
%%									4: Inhalt der Node

%% Befehl, um einfach Beschriftungen mit Pfeil hinzuzuf�gen: (h�lt Abstand nach unten/oben automatisch ein! (standardm��ig in footnotesize aber durch voranstellen einer gr��e �berschreibbar!)
% #1 (optional, def 0.5cm): Abstand der Beschriftung vom zu Beschriftenden
% #2 (optional, def 270): Winkel, in dem die Beschriftung angebracht wird.
% #3 zu beschriftender Text
% #4 Beschriftung 
\DeclareDocumentCommand \pin {O{0.5cm} O{270} m m}{
\tikz[baseline=(mynode.base), pin distance=#1,every pin edge/.style={<-}, inner sep=1pt]{
\node(mynode)[pin=#2:{\footnotesize  #4}]{#3};
\coordinate (pin) at (current bounding box.south);
\path let \p1=(pin), \p2=(mynode.center) in %
	node[inner sep=0pt](da) at (\x2,\y1){} node(gesamt)[fit=(mynode)(da),inner sep =0pt]{};
\pgfresetboundingbox
\useasboundingbox (gesamt.north west)rectangle(gesamt.south east);
%\draw[fill=red] (current bounding box.south west) circle(0.02cm);
%\draw[fill=red] (current bounding box.south east) circle(0.02cm);
}
}

\DeclareDocumentCommand \pintest{O{0.5cm} O{270} m m}{#1 #2 #3 #4}

\newcommand{\Highlight}[1][submatrix]{%
    \tikz[overlay,remember picture]{
    \node[highlight,fit=(left.north west) (right.south east)] (#1) {};}
}
\newcommand{\highlight}[2]{%
	\tikz[overlay,remember picture]{
	\node[highlight,fit=(#1.north west)(#1.south east)] (#2){};
	}
}
\newcommand{\underbrac}[2]{%
	\tikz[overlay,remember picture, very thin, orange]{
	\node[inner sep=-2pt,fit=(#1.north west)(#1.south east)] (#2){};
	\path (#1.south west)++(0.065,0.05) edge($(#1.south west)+(0.065,0)$)
		  ($(#1.south west)+(0.065,0)$) edge ($(#1.south east)-(0.065,0)$)
  	 	 ($(#1.south east)-(0.065,0)$)++(0,0.05) edge  ($(#1.south east)-(0.065,0)$);
	}
}
%%%%%%%%%%%%%%%%%%%%%%%%%%%%%%%%%%%%%%%
\newcommand{\tikzdoubleul}[1]{ % doppelte Unterstreichungen auch fuer groessere Objekte
\tikzmarktwo{tikzdoubleul}{#1}
\tikz[overlay, remember picture]{
\path (tikzdoubleul.south west) edge[double, thick, double distance = 1.5pt] (tikzdoubleul.south east);
}
}

\newcommand{\uldash}[1]{ % gestrichelte Unterstreichung
\tikzmarktwo[2pt]{uldash}{#1}
\tikz[overlay, remember picture]{
\path (uldash.south west) edge[draw, thick, dashed] (uldash.south east);
}
}

\newcommand{\tikzcancel}[1]{ % cancel gestrichelt
\tikzmarktwo{tikzcancel}{#1}
\tikz[overlay, remember picture]{
\path ($(tikzcancel.south west)-(1pt,2pt)$) edge[draw,dotted,thick,color =gray] ($(tikzcancel.north east)+(1pt,2pt)$);
}
}


%%%%%%%%%%%%%%%%%%%%%%%%%%%%%%
% fuer bessere Underbraces!!
% normal den underbrace Befehl verwenden aber dann den Text mathclapen! Bsp: \underbrace{a}{\mathclap{text}}
 \def\mathllap{\mathpalette\mathllapinternal}
 \def\mathllapinternal#1#2{%
 \llap{$\mathsurround=0pt#1{#2}$}%  $
 }
 \def\clap#1{\hbox  to  0pt{\hss#1\hss}}
 \def\mathclap{\mathpalette\mathclapinternal}
 \def\mathclapinternal#1#2{%
 \clap{$\mathsurround=0pt#1{#2}$}%
}
 \def\mathrlap{\mathpalette\mathrlapinternal}
 \def\mathrlapinternal#1#2{%
 \rlap{$\mathsurround=0pt#1{#2}$}%  $
 }
%%%%%%%%%%%%%%%%%%%%%%%%%%%%%%
%%%%%%%%%%%%%%%%%%%%%%%%%%%%
% underbrackets!
\makeatletter
\def\overbracket#1{\mathop{\vbox{\ialign{##\crcr\noalign{\kern3\p@}
\downbracketfill\crcr\noalign{\kern3\p@\nointerlineskip}
$\hfil\displaystyle{#1}\hfil$\crcr}}}\limits}
\def\underbracket#1{\mathop{\vtop{\ialign{##\crcr
$\hfil\displaystyle{#1}\hfil$\crcr\noalign{\kern3\p@\nointerlineskip}
\upbracketfill\crcr\noalign{\kern3\p@}}}}\limits}
\def\overparenthesis#1{\mathop{\vbox{\ialign{##\crcr\noalign{\kern3\p@}
\downparenthfill\crcr\noalign{\kern3\p@\nointerlineskip}
$\hfil\displaystyle{#1}\hfil$\crcr}}}\limits}
\def\underparenthesis#1{\mathop{\vtop{\ialign{##\crcr
$\hfil\displaystyle{#1}\hfil$\crcr\noalign{\kern3\p@\nointerlineskip}
\upparenthfill\crcr\noalign{\kern3\p@}}}}\limits}
\def\downparenthfill{$\m@th\braceld\leaders\vrule\hfill\bracerd$}
\def\upparenthfill{$\m@th\bracelu\leaders\vrule\hfill\braceru$}
\def\upbracketfill{$\m@th\makesm@sh{\llap{\vrule\@height3\p@\@width.7\p@}}%
\leaders\vrule\@height.7\p@\hfill
\makesm@sh{\rlap{\vrule\@height3\p@\@width.7\p@}}$}
\def\downbracketfill{$\m@th
\makesm@sh{\llap{\vrule\@height.7\p@\@depth2.3\p@\@width.7\p@}}%
\leaders\vrule\@height.7\p@\hfill
\makesm@sh{\rlap{\vrule\@height.7\p@\@depth2.3\p@\@width.7\p@}}$}
\makeatother
%%%%%%%%%%%%%%%%%%%%%%%%%%%%%%%%%%%%

% Markierung als Studenten rechnen lassen>
\newcommand{\markstudent}[3] %Argumente: 1: Abstand zwischen jetziger Zeile und Inhalt (Bei Equations: Baselineskip)
% 2.: L�nge des Inhalts /Markierung
%3. : Inhalt
{
\ifLeiter
\begin{minipage}[t]{0.05\textwidth}	
	\vspace{#1}
	\unitlength1cm
	{\color{orange}
	\makebox[0.025\textwidth]{\line(0,-1){#2}} 
	\makebox[0.025\textwidth]{\line(0,-1){#2}}}
\end{minipage}
\begin{minipage}[t]{0.95\textwidth}
#3
\end{minipage}
\else
#3
\fi
}

% Markierung f�r �bungsleiter>
\newcommand{\markteacher}[3] %Argumente: 1: Abstand zwischen jetziger Zeile und Inhalt (Bei Equations: Baselineskip)
% 2.: L�nge des Inhalts /Markierung
%3. : Inhalt
{
\ifLeiter
\begin{minipage}[t]{0.05\textwidth}	
	\vspace{#1}
	\unitlength1cm
	{\color{ForestGreen}
	\makebox[0.025\textwidth]{\begin{sideways}\uwave{\hspace{#2}}\end{sideways}}
	\makebox[0.025\textwidth]{\begin{sideways}\uwave{\hspace{#2}}\end{sideways}}  
	}
\end{minipage}
\begin{minipage}[t]{0.95\textwidth}
#3
\end{minipage}
\else
#3
\fi
}

% Markierung f�r �bungsleiter, einzeilig!
\newcommand{\markteachers}[1] 
{
\ifLeiter
\begin{minipage}[t]{0.05\textwidth}	
	\vspace{-0.75\baselineskip}
	\unitlength1cm
	{\color{OliveGreen}
	\makebox[0.025\textwidth]{\begin{sideways}\uwave{\hspace{0.85\baselineskip}}\end{sideways}}
	\makebox[0.025\textwidth]{\begin{sideways}\uwave{\hspace{0.85\baselineskip}}\end{sideways}}  
	}
\end{minipage}
\hspace{-0.5cm}
\begin{minipage}[t]{0.95\textwidth}
#1%
\end{minipage}
\else
#1
\fi
}

% Abs�tze machen:
\newcommand{\absatz}[2]
{
\ifdefined\wortlang 
\else
\newlength{\wortlang}
\fi
\settowidth{\wortlang}{#1}
\ifdim \wortlang<0.5\linewidth
\begin{tabular}{p{\wortlang}p{\linewidth-4\tabcolsep-\wortlang}}
\else
\begin{tabular}{p{0.5\linewidth-2\tabcolsep}p{0.5\linewidth-2\tabcolsep}}
\fi
#1 & #2
\end{tabular}
\let\wortlang\relax
}




\newcommand{\koords}{
\put(0,0){\makebox(0,0)[rb]{
\begin{picture}(0,0)
\put(4,0){\line(0,0){0,1}$4$}
\put(3,0){\line(0,0){0,1}$3$}
\put(2,0){\line(0,0){0,1}$2$}
\put(1,0){\line(0,0){0,1}$1$}
\put(0,0){\line(0,0){0,1}$0$}
\put(-1,0){\line(0,0){0,1}$1$}
\put(-2,0){\line(0,0){0,1}$2$}
\put(-3,0){\line(0,0){0,1}$3$}
\put(-4,0){\line(0,0){0,1}$4$}
\put(-5,0){\line(0,0){0,1}$5$}
\put(-6,0){\line(0,0){0,1}$6$}
\put(-7,0){\line(0,0){0,1}$7$}
\put(-8,0){\line(0,0){0,1}$8$}
\put(-9,0){\line(0,0){0,1}$9$}
\put(-10,0){\line(0,0){0,1}$10$}
\put(-11,0){\line(0,0){0,1}$11$}
\put(-12,0){\line(0,0){0,1}$12$}
\put(-13,0){\line(0,0){0,1}$13$}
\put(-14,0){\line(0,0){0,1}$14$}
\put(-15,0){\line(0,0){0,1}$15$}
\put(-16,0){\line(0,0){0,1}$16$}
\put(-17,0){\line(0,0){0,1}$17$}
\put(-18,0){\line(0,0){0,1}$18$}
\put(-19,0){\line(0,0){0,1}$19$}
\put(-20,0){\line(0,0){0,1}$20$}
\put(-21,0){\line(0,0){0,1}$21$}
\put(-22,0){\line(0,0){0,1}$22$}
\put(-23,0){\line(0,0){0,1}$23$}
\put(-24,0){\line(0,0){0,1}$24$}
\put(0,-5){\line(-1,0){0,1}$5$}
\put(0,-4){\line(-1,0){0,1}$4$}
\put(0,-3){\line(-1,0){0,1}$3$}
\put(0,-2){\line(-1,0){0,1}$2$}
\put(0,-1){\line(-1,0){0,1}$1$}
\put(0,0){\line(-1,0){0,1}}
\put(0,1){\line(-1,0){0,1}$1$}
\put(0,2){\line(-1,0){0,1}$2$}
\put(0,3){\line(-1,0){0,1}$3$}
\put(0,4){\line(-1,0){0,1}$4$}
\put(0,5){\line(-1,0){0,1}$5$}
\put(0,6){\line(-1,0){0,1}$6$}
\put(0,7){\line(-1,0){0,1}$7$}
\put(0,8){\line(-1,0){0,1}$8$}
\put(0,9){\line(-1,0){0,1}$9$}
\put(0,10){\line(-1,0){0,1}$10$}
\put(0,11){\line(-1,0){0,1}$11$}
\put(0,12){\line(-1,0){0,1}$12$}
\put(0,13){\line(-1,0){0,1}$13$}
\put(0,14){\line(-1,0){0,1}$14$}
\put(0,15){\line(-1,0){0,1}$15$}
\put(0,16){\line(-1,0){0,1}$16$}
\end{picture}
}
}
}

\renewcommand{\L}{\mathrsfs{L}}

\thispagestyle{empty}
\setcounter{page}{-1}

\begin{textblock*}{47mm}(30mm,33mm)
\includegraphics[width=47mm]{ees}
\end{textblock*}

\begin{textblock*}{93mm}(127mm,192mm)
\includegraphics[width=93mm]{fausiegel}
\end{textblock*}

{\raggedleft
\textsc{Lehrstuhl f�r Elektrische Energiesysteme}\\
Vorstand: Univ.-Prof. Dr.-Ing. Matthias Luther
\par}

\vspace{51mm}

{\centering

%##########################################################
% BEARBEITUNGSFELDER
\large{Art der Arbeit und Nr} \\	
\Large{Hier steht das Thema der schriftlichen Ausarbeitung}
%##########################################################

\par}

\vspace{105mm}

{\raggedright
\begin{tabbing}
XX \= XXXXXXXXXXX \= XXXXXXXXXXXXXXXXXXXXXX \kill
%##########################################################
% BEARBEITUNGSFELDER
		\> \textbf{Bearbeiter:} 	\> Vorname Nachname \\
		\>												\> Matrikelnummer		\\
 		\>												\>									\\
		\> \textbf{Betreuer:}			\> Vorname Nachname \\
		\>												\>									\\
		\> \textbf{Abgabedatum:}	\> TT.MM.JJJJ															
%##########################################################
\end{tabbing}
\par}

\thispagestyle{empty}
\cleardoublepage														% Deckblatt

Ich versichere, dass ich die vorliegende Arbeit ohne fremde Hilfe und ohne Benutzung anderer als der angegebenen Quellen angefertigt habe, und dass die Arbeit in gleicher oder �hnlicher Form noch keiner anderen Pr�fungsbeh�rde vorgelegen hat und von dieser als Teil einer Pr�fungsleistung angenommen wurde. Alle Ausf�hrungen, die w�rtlich oder sinngem�� �bernommen wurden, sind als solche gekennzeichnet.\\
\\

\begin{table}[!htb]
\centering
\begin{tabularx}{\textwidth}{lXl}

       \hspace{6cm} &  & \hspace{6cm} \\
\cline{1-1}\cline{3-3}
Ort, Datum  &  & Unterschrift
\end{tabularx}
\end{table}														% Erkl�rung der Eigenarbeit

\chapter*{Aufgabenstellung der Arbeit}
%\vspace*{4cm}

{\large \textbf{Thema:} \parbox[t]{0.8\textwidth}{Thema der Arbeit}}
\\

Hier wird die Aufgabenstellung beschrieben. Die Notwendigkeit dieser h�ngt vom Betreuer ab. 															% Aufgabenstellung (optional)

%\chapter*{Vorwort}

%\vspace*{4cm}

An dieser Stelle kann ein Vorwort eingef�gt werden.															% Vorwort (optional)

\cleardoublepage
\begin{spacing}{1.15}																	% evtl. kleinerer Zeilenabstand im IV, AV, TV
\pdfbookmark[1]{Inhaltsverzeichnis}{toc}							% Inhaltsverzeichnis bei den Lesezeichen rein
\tableofcontents 																			% Inhaltsverzeichnis erzeugen
\listoffigures   																			% Abbildungsverzeichnis
\listoftables   																			% Tabellenverzeichnis
\end{spacing}

%\addchap{Abstract}

Hierbei handelt es sich um einen, in englischer Sprache verfassten, kurzen Abriss der
Arbeit (Abstract). Der Umfang betr�gt eine halbe bis eine Seite. Alternativ zu dem
Abstract kann eine in englischer Sprache verfasste Zusammenfassung (Summary) nach
der deutschen Zusammenfassung stehen. Im Hauptseminar kann dieser Teil entfallen.															% Abstract

\mainmatter																						% Hauptteil beginnt

\chapter{Einleitung}
\label{kap:einleitung}
Die Analyse von komplexen Energieversorgungsnetzen ist ohne das Hilfsmittel Computer kaum denkbar. Die Gr��e der Netze, sowie die Anzahl deren Komponenten und Speicherelemente macht eine rein analytische Berechnung sehr schwierig. Der Einsatz des Computers hat es dem Ingenieur nicht nur in der Energietechnik erm�glicht, verschiedenste dynamische Systeme zu simulieren, um dadurch Systemeigenschaften zu beobachten oder mit unerprobten Konzepten zu experimentieren. 
Dynamische Systeme sind dabei im allgemeinen dadurch gekennzeichnet, dass sie Energiespeicher enthalten, die zu zeitver�nderlichen Vorg�ngen f�hren. Aus physikalischer Sicht sind diese Vorg�nge durch Differentialgleichungen charakterisiert, deren L�sungen sich nicht immer in analytischer Form, oder zumindest einfach angeben lassen.
Der Computer liefert die M�glichkeit, mit numerischen Methoden N�herungen dieser L�sungen zu berechnen.
Hierf�r gibt es eine Vielzahl verschiedener Methoden, welche unter gewissen Voraussetzungen zu unterschiedlichen Ergebnissen oder zu unterschiedlicher Genauigkeit f�hren k�nnen. 
Um Simulationssoftware anwenden zu k�nnen ist es deshalb notwendig, dass der Benutzer die Grundlagen der hinter der Software steckenden Numerik versteht.

Aus diesem Grund werden in dieser Arbeit die wichtigsten Verfahren zu numerischen Simulation von dynamischen Systemen vorgestellt. Dabei wird auf Einschrittverfahren, sowie auf Mehrschrittverfahren eingegangen. Dar�ber hinaus wird die Vorgehensweise der Schrittweitensteuerung erl�utert.
Im \ref{sec:stabi}. Kapitel werden die verschiedenen Verfahren auf Stabilit�t untersucht, welche eine essenzielle Eigenschaft der numerischen Verfahren darstellt.
Um die Bedeutung der Verfahren zu verdeutlichen, werden in Kapitel \ref{sec:anwend} einige Softwarepakete der zur Netzsimulation auf ihr verwendetes Verfahren untersucht.
														% Einleitung (Bezeichnung passend zur Arbeit!)

\chapter{Numerische Integrationsverfahren f�r Anfangswertprobleme}
\label{kap:hauptteil}
Die Modellbildung dynamischer Systeme f�hrt in vielen F�llen naturgem�� zu gew�hnlichen Differentialgleichungen (beziehungsweise Differentialgleichungssystemen) in Abh�ngigkeit der Zeit $t$, so genannten \textit{Anfangswertproblemen}:
\begin{align}
\dot{x}(t)=f\left(x(t),t\right), \quad x(0)=x_0.
\label{eq:AWA}
\end{align}
Dabei ist die zeitliche Ableitung der Zust�nde $\dot{x}(t)=\frac{\d}{\d t}x(t)$ zu einem bestimmten Zeitpunkt in Form der Funktion $f$ berechenbar. Dabei mag die Funktion $f$ auch nicht in analytische Form dargestellt werden k�nnen, sondern lediglich der Funktionswert berechenbar sein. Dar�ber hinaus ist der Anfangswert der L�sung $x(0)=x_0$ bekannt.\\
Um derartige Systeme simulieren zu k�nnen, muss eine L�sung $x(t)$ des Anfangswertproblems gefunden werden. Die exakte L�sung dieses Problems ergibt sich durch Integration der zeitlichen Ableitung:
\begin{align}
x\tt = x_0 + \int\limits_0^t \dot{x}(x,t) \d t.
\label{eq:integral}
\end{align}
Die Herausforderung der Simulation besteht also darin, eine numerische Approximation dieser Integration zu ermitteln, was das Problem auf die numerische Integration von Anfangswertproblemen reduziert. Diese Approximation soll zu bestimmten, diskreten Zeitpunkten $t=t_i=i \dd T, i=0,\ldots,n$ N�herungsl�sungen $x_i$ f�r den Tats�chlichen Funktionswert $x(t_i)$ liefern. Die Schrittweite $T$ legt dabei den Abstand von zwei benachbarten Abtastzeitpunkten $t_i$ fest. 

Es soll also allgemein eine Berechnungsvorschrift
\begin{align}
	x_{i+1} = \Phi(x, t)
\end{align}
mit der Verfahrensfunktion $\Phi(x,t)$ bestimmt werden, welche zum einen eine m�glichst genaue Approximation der tats�chlichen L�sung liefert. Zum anderen soll das Verfahren so gew�hlt werden, dass auch m�glichst komplexe Systeme approximiert werden k�nnen, das hei�t, es ist erstrebenswert, den dabei anfallenden Rechenaufwand gering zu halten.

Im Folgenden werden verschiedene Verfahren zur Integration vorgestellt. Diese werden in Einschrittverfahren und Mehrschrittverfahren unterteilt. Soweit nicht anders angegeben wird sich dabei auf \cite{gruene} bezogen.
%%%%%%%%%%%%%%%%%%%%%%%%%%%%%%%%%%%%%%%%%%%%%%%
\section{Einschrittverfahren}
\subsection{Explizites Eulerverfahren}
Der einfachste Ansatz, (\ref{eq:integral}) zu approximieren, ist, die Steigung der Funktion $x(t)$ im Zeitintervall zwischen $t_i$ und $t_{i+1}$ als konstant anzunehmen. Dies entspricht dem Taylor-Satz erster Ordnung und f�hrt auf die \textit{explizite} Berechnungsvorschrift
\begin{align}
		x_{i+1} = x_i + T f(x_i,t_i).
\end{align}
Dieses Verfahren wird als \textit{explizites Eulerverfahren} bezeichnet. 

\subsection{Impliztes Eulerverfahren}
Beim expliziten Eulerverfahren wird angenommen, dass die im Zeitpunkt $t_i$ vorliegende Steigung im gesamten Intervall $t_i \ldots t_{i+1}$ g�ltig ist. Genauso gut k�nnte angenommen werden, dass im betrachteten Intervall die Steigung des Zeitpunktes $t_{i+1}$ vorliegt:
\begin{align}
	x_{i+1}&= x_i + f(x_{i+1},t_{i+1})
\end{align}
In diesem Fall ergibt sich eine \textit{implizite Gleichung} f�r den approximierten Wert $x_{i+1}$, welche nur unter Kenntnis von $f$ aufgel�st werden kann. Deshalb geh�rt das Verfahren zur Klasse der \textit{impliziten Verfahren} und wird als \textit{implizites Eulerverfahren} bezeichnet.

\subsection{Trapezregel/Heun-Verfahren}
Einen Ansatz, die Verfahren zu verbessern, liefert die so genannte \textit{Trapezregel}. Bei dieser werden die beiden Euler-Verfahren gewisserma�en kombiniert.
Dabei wird f�r die Steigung im betrachteten Intervall der Mittelwert der Steigungen im $i$-ten und $i+1$-ten Schritt angenommen:
\begin{align}
	x_{i+1} = x_i + T \dd \eh \dd \Big(f\big(x_i,t_i\big) + f\big(x_{i+1},t_{i+1}\big)  \Big).
\end{align}
Dadurch ergibt sich allgemein ebenfalls ein implizites Verfahren. Um das Aufl�sen der impliziten Gleichung zu umgehen, kann der unbekannten Wert $f\big(x_{i+1},t_{i+1}\big)$ durch das explizite Euler-Verfahren abgesch�tzt werden. Dies f�hrt auf das \textit{Heun-Verfahren}:
\begin{align*}
	x_{i+1}	= x_i + T \dd \eh \dd \Big(f\big(x_i,t_i\big) + f\big( x_i +  T \dd f(x_i,t_i)  ,t_{i+1}\big) \Big).
	\label{eq:Heun}
\end{align*}

\subsection{Runge-Kutta-Verfahren}
Um die bisher betrachteten Methoden zu verbessern, gibt es mehrere Ans�tze. Ein naheliegender Gedanke ist, statt nur der rechten Seite der Differentialgleichung (\ref{eq:AWA}), Zeitableitungen h�herer Ordnung in die Absch�tzung einzubauen. Dies f�hrt zu den Taylor-Methoden, welche auf dem Taylor-Satz basieren.
Da die Differentation der Rechten Seite jedoch mit verh�ltnism��ig gro�em rechnerischen Aufwand verbunden ist, sind diese Verfahren in der Praxis eher ungebr�uchlich und sollen im Rahmen dieser Arbeit nicht n�her betrachtet werden.

Um einen anderen Ansatz zu zeigen, w�hlt man den \textit{globalen} Fehler $e(t_i):= \norm{x_i - x(t_i)}$, den ein Verfahren im Vergleich zur exakten L�sung macht, als ein Ma�, mit dem die Verfahren verglichen werden k�nnen. Es kann gezeigt werden, dass sich dieser f�r die betrachteten Integrationsverfahren mit der Landau-Notation durch
\begin{align}
	e(t) = \mathcal{O}(T^p)
\end{align}
absch�tzen l�sst. (Getroffene Annahmen �ber die rechte Seite der Differentialgleichung (\ref{eq:AWA}) sollen im Rahmen dieser Arbeit au�er Acht gelassen werden und k�nnen in \cite{gruene} nachgelesen werden.) Dabei legt die Konsistenzordnung $p$ fest, wie schnell der globale Fehler gegen Null geht, wenn die Schrittweite verkleinert wird und stellt damit ein Ma� f�r die \glqq G�te\grqq \ des Verfahrens dar. Es kann gezeigt werden, dass die Konsistenzordnung der Euler-Verfahren $p=1$ betr�gt. Das Heun-Verfahren ist mit $p=2$ bereits deutlich genauer.\\
Deshalb ist der erste Ansatz in der Praxis, die Idee der Trapezregel beziehungsweise des Heun-Verfahrens weiter zu f�hren und die Konvergenzordnung der Verfahren durch weitere St�tzstellen zu erh�hen.
Dies f�hrt zum allgemeinen Ansatz der $s$-\textit{stufigen expliziten Runge-Kutta-Verfahren}. Aus Gr�nden der �bersichtlichkeit wurde der Index $_i$ zum Kennzeichnen des aktuell betrachteten Abtastpunktes im Folgenden weggelassen.
\begin{align}
	\Phi(x,t,T)&= x + T\, \sum\limits_{i=1}^s b_i\,k_i, \quad \text{mit} \\
	k_i &= f\left(  x + T \,\sum\limits_{j=1}^{i-1} a_{ij} k_j ,t+c_i T    \right) \quad \text{f�r } i=1,\ldots,s.
\end{align}
Dabei wird die Rechte Seite der DGL an den St�tzstellen $t+c_i\,T$ berechnet, wobei die vorhergehenden St�tzstellen gewichtet mit den Koeffizienten $a_{ij}$ zur Berechnung des zur St�tzstelle geh�renden Funktionswertes herangezogen werden. Die Gewichtung der berechneten Steigungen bei der Ermittlung der resultierenden Verfahrensfunktion erfolgt durch die Koeffizienten $b_i$. Mit dieser Darstellung ergeben die Runge-Kutta-Verfahren eine unendliche Menge verschiedener Verfahren zur L�sung des Problems.\\
Um die Runge-Kutta-Verfahren zu klassifizieren, werden die Koeffizienten $a_{ij},b_{i},c_{i}$ h�ufig in Form des Butcher-Schemas, nach Tabelle \ref{tab:ButEx} dargestellt.
\begin{table}
\centering
$
	\begin{array}{c|ccccc}
		c_1\\
		c_2 & a_{21}\\
		c_3 & a_{31} & a_{32}\\
		\vdots & \vdots & \vdots & \ddots \\
		c_s & a_{s1} & a_{s2}&\cdots & a_{s \, s-1}\\ \hline
		& b_1 & b_2 & \cdots & b_{s-1} & b_s
	\end{array}
$\\[\baselineskip]
\caption{Allgemeines Butcher-Schema expliziter Runge-Kutta-Verfahren}
\label{tab:ButEx}
\end{table}
Sowohl das explizite Euler-Verfahren (s=p=1), als auch das Heun-Verfahren (s=p=2) lassen sich auf diese Weise beschreiben. Dar�ber hinaus ist das \textit{klassische} Runge-Kutta-Verfahren (s=p=4) von gro�er Bedeutung. F�r diese Verfahren ergeben sich die in Tabelle \ref{tab:Butcherexamples} dargestellten Butcher-Schemata.
\begin{table}
\centering
$
	\begin{array}[b]{c|c}
		0 &\\ \hline
		&1
	\end{array} \quad
	\begin{array}[b]{c|cc}
		0 &\\ 
		1&1 \\ \hline
		& \eh & \eh
	\end{array} \quad 
	\begin{array}[b]{c|cccc}
		0 &\\ 
		\eh&\eh \\ 
		\eh & 0 & \eh \\
		1&0&0&1 \\ \hline
		& \frac{1}{6}&\frac{2}{6}&\frac{2}{6}&\frac{1}{6}
	\end{array}
$\\[\baselineskip]
\caption{Butcher-Schemata f�r (von links nach rechts): Euler-Verfahren, Heun-Verfahren, \textit{klassisches} Runge-Kutta-Verfahren}
\label{tab:Butcherexamples}
\end{table}
\\Die selbe Vorgehensweise l�sst sich f�r implizite Verfahren anwenden, sodass sich als Berechnungsvorschrift f�r $s$-\textit{stufige implizite Runge-Kutta-Verfahren} ergibt:
\begin{align}
	\Phi(x,t,T) &= x + T\, \sum\limits_{i=1}^s b_i \, k_i \quad \text{mit}\\
	k_i &= f\left(  x + T\, \sum\limits_{j=1}^{s} a_{ij}k_j,  t+c_i T   \right) \quad \text{f�r } i=1,\ldots,s. 
\end{align}
Beispiele f�r implizite Runge-Kutta-Verfahren sind das bereits kennengelernte implizite Euler-Verfahren (s=p=1), die implizite Mittelpunktsregel (s=p=2), sowie die implizite Trapezregel (s=p=2), deren Butcher-Schemata in Tabelle \ref{tab:ButcherIm} dargestellt sind.
\begin{table}
	\centering
$
	\begin{array}[b]{c|c}
		1 & 1\\ \hline
		 &1
	\end{array} \quad
	\begin{array}[b]{c|c}
		\eh&\eh \\ \hline
		& 1
	\end{array} \quad 
	\begin{array}[b]{c|cc}
		0 & 0 & 0\\ 
		1&\eh & \eh \\  \hline
		& \eh & \eh
	\end{array}
$\\[\baselineskip]
\caption{Butcher-Schemata f�r implizite Verfahren (von links nach rechts): Euler, Mittelpunktsregel, Trapezregel}
\label{tab:ButcherIm}
\end{table}
Der Vorteil der impliziten Verfahren gegen�ber den expliziten Verfahren wird sich bei Betrachtung des Stabilit�tsverhaltens in Kapitel \ref{sec:stabi} zeigen. 

%%%%%%%%%%%%%%%%%%%%%%%%%%%%%%%%%%%%%%%%%%%%%
\section{Schrittweitensteuerung (Adaption)}
% NACH DER BESCHREIBUNG ALLER VERFAHREN!
Um die Verfahren weiter zu verbessern und gleichzeitig Rechenzeit einzusparen, liegt der Gedanke nahe, die bisher als konstant angenommene Schrittweite $T$ als steuerbar anzusehen, sodass $T_i$ an Stellen, bei welchen sich der Funktionswert st�rker �ndert, kleiner gew�hlt werden kann und an Stellen langsamer �nderungen hingegen gr��er. Aus Gr�nden der numerischen Effizienz ist es �blich, dass im $i$-ten Zeitschritt eine \glqq gute\grqq \ Schrittweite f�r den �bergang von $t_i$ nach $t_{i+1}$ bestimmt wird, jedoch keine Anpassung von vorhergehenden Schrittweiten mehr erfolgt.

Um die Schrittweite anzupassen wird gefordert, dass der \textit{lokale} Fehler $\varepsilon(t)$ unter einer gewissen Toleranzgrenze $tol$ liegen soll. Der \textit{lokale} Fehler ist dabei der Anteil am \textit{globalen} Fehler, der durch den Zeitschritt von $t_i$ bis $t_{i+1}$ hervorgerufen wird. Da die exakte L�sung des Funktionsverlaufs nicht bekannt ist, wird f�r das verwendete Verfahren $\Phi$ mit Konsistenzordnung $p$ ein anderes Verfahren $\hat{\Phi}$ zur Absch�tzung des Fehlers herangezogen. Die Konsistenzordnung des zweiten Verfahrens $\hat{p}$ wird kleiner als die des verwendeten Verfahrens gew�hlt, sodass $p \geq \hat{p}+1$ gilt. Der Schritt von $t_i$ nach $t_{i+1} = t_i+T_i$ wird mit beiden Verfahren berechnet und die Norm der Differenz der Ergebnisse als \textit{Fehlersch�tzer} $\overline{\varepsilon}$ bezeichnet.
\begin{equation}
	\overline{\varepsilon} := \norm{\hat{\Phi}(t_i,x_i,T_i) - \Phi(t_i,x_i,T_i)}.
\end{equation}
Mit Hilfe der Konvergenzeigenschaften l�sst sich zeigen, dass die gew�nschte Fehlertoleranz n�herungsweise eingehalten wird, falls die Schrittweite durch 
\begin{equation}
	T_{\mathrm{neu}} = \sqrt[\hat{p}+1]{\frac{tol}{\overline{\varepsilon}}}T_{\mathrm{alt}}
	\label{eq:Schrittweite}
\end{equation}
angepasst wird. Da dies nur approximativ gilt, wird das Argument der Wurzel in der Praxis jedoch h�ufig durch einen \glqq Sicherheitsfaktor\grqq \ $f_{ac} \in ]0,1[$ (typisch: $0.9$) skaliert.

Nach der Berechnung des Fehlersch�tzers $\overline{\varepsilon}$ im Zeitschritt $i$ wird dieser mit der zul�ssigen Toleranz verglichen. Ist der berechnete Fehler gr��er, als die zul�ssige Toleranz ($\overline{\varepsilon} > tol$), wird der selbe Zeitschritt mit der angepassten Schrittweite $T_{\mathrm{neu}}$ erneut berechnet.\\
Gilt $\overline{\varepsilon} \leq tol$, so ist die gew�nschte Genauigkeit erreicht und die der Berechnung zu Grunde liegende Schrittweite wird f�r diesen Zeitschritt angenommen. Nun wird erneut (\ref{eq:Schrittweite}) berechnet und $T_\mathrm{neu}$ als Schrittweite $T_{i+1}$ f�r den n�chsten Zeitschritt angenommen. Dies sorgt daf�r, dass die Schrittweite bei hoher Genauigkeit auch vergr��ert werden kann.

%%%%%%%%%%%%%%%%%%%%%%%%%%%%%%%%%%%%%%%%%
\section{Eingebettete Runge-Kutta-Verfahren}
Um den Rechenaufwand der Beschriebenen Schrittweitenanpassung zu reduzieren, k�nnen die Verfahrensfunktionen $\Phi$ und $\hat{\Phi}$ so geschickt gew�hlt werden, dass bei deren Auswertungen auf gleiche Zwischenergebnisse zur�ckgegriffen werden kann. Dies f�hrt auf die sehr h�ufig verwendeten eingebetteten Runge-Kutta verfahren. Die eingebetteten Verfahren werden mit $\mathrm{RK}p(\hat{p})$ bezeichnet.
Der Gedanke dahinter ist, dass sich die berechneten Werte der St�tzstellen $k_i$ und $\hat{k}_i$ nicht unterscheiden, sondern lediglich deren Linearkombination. Das hei�t die Koeffizienten $a_{ij}=\hat{a}_{ij}$ und $c_{i}=\hat{c}_{i}$ stimmen in beiden Verfahren �berein. Es gilt lediglich $b_{i}\neq \hat{b}_{j}$. 
F�r eingebettete Runge-Kutta-Verfahren lassen sich dehalb auch sehr einfach Butcher-Schemata nach Tabelle \ref{tab:ButcherEin} erstellen.
\begin{table}
	\centering
$
	\begin{array}{c|ccccc}
		c_1\\
		c_2 & a_{21}\\
		c_3 & a_{31} & a_{32}\\
		\vdots & \vdots & \vdots & \ddots \\
		c_s & a_{s1} & a_{s2}&\cdots & a_{s \, s-1}\\ \hline
		& b_1 & b_2 & \cdots & b_{s-1} & b_s \\ \hline
		& \hat{b}_1 & \hat{b}_2 & \cdots & \hat{b}_{s-1} & \hat{b}_{s}
	\end{array}
$\\[\baselineskip]
\caption{allgemeines Butcher-Schema f�r eingebettete Runge-Kutta-Verfahren}
\label{tab:ButcherEin}
\end{table}
Wie in \cite{gruene} gezeigt wird, muss in diesem Fall zur Konstruktion eines Verfahrens der Konsistenzordnung $\hat{p}=p-1$ ein Verfahren mit Stufenzahl $\hat{s}=s+1$ gew�hlt werden. F�r das klassische Runge-Kutta-Verfahren (s=p=4) ergibt sich dadurch das Butcher-Schema nach Tabelle \ref{tab:koeffbest}.
\begin{table}
		\centering
$
	\begin{array}{c|ccccc}
		0\\
		\eh & \eh \\[2pt]
		\eh & 0 & \eh \\[2pt]
		1 & 0 & 0 & 1 \\[2pt]%
		c_5 & a_{51} & a_{52}& a_{53} & a_{54}\\\hline
		& \tikzmark[2pt]{1}{$\frac{1}{6}$} & \frac{2}{6} & \frac{2}{6} & \frac{1}{6} & 0 \\[2pt] \hline
		& \tikzmark[2pt]{2}{$\hat{b}_1$} & \hat{b}_2 & \hat{b}_3 & \hat{b}_4 & \hat{b}_5
	\end{array}
$\\[\baselineskip]
\caption{Butcher-Schema zur Bestimmung der unbekannten Koeffizienten}
\label{tab:koeffbest}
\end{table}
Um das zugeh�rige Verfahren $\hat{\Phi}$ mit $\hat{s}=5$ und $\hat{p}=3$ zu bestimmen, m�ssen also die unbekannten Koeffizienten in \ref{tab:koeffbest} bestimmt werden. Dies kann in \cite{gruene} nachgelesen werden und soll an dieser Stelle nicht weiter vertieft werden.
Eines der am weitesten verbreiteten Verfahren dieser Klasse ist das Dormand-Prince-RK5(4)-Verfahren, welches in der Numerik-Software \Matlab \ unter dem Namen \textsc{ode45}  standardm��ig eingestellt ist. Die zugeh�rigen Koeffizienten k�nnen in Anhang \ref{sec:anheingeb} nachgelesen werden.

%%%%%%%%%%%%%%%%%%%%%%%%%%%%%%%%%%%%%%%%%%%%
\section{Mehrschrittverfahren}
Die bisher betrachteten Verfahren sind dadurch charakterisiert, dass die Verfahrensfunktion $\Phi(x,t)$ lediglich vom aktuellen Funktionswert $x_i$ abh�ngt. Deshalb werden diese Verfahren als Einschrittverfahren bezeichnet.\\
Im Gegensatz dazu h�ngt der Funktionswert $x_{i+1}$ bei Mehrschrittverfahren zus�tzlich dazu von einer beliebigen Anzahl an Vorg�ngerwerten $x_{i-k+1},\ldots,x_i$ ab. Damit wird ein $k$-stufiges lineares Mehrschrittverfahren folgenderma�en definiert:
\begin{align}
	&a_k x_{i+k}+a_{k-1} x_{i+k-1}+\ldots+a_0 x_i \notag\\
	&\quad = T\left( b_k f(x_{i+k},t_{i+k}) + b_{k-1}f(x_{i+k-1},t_{i+k-1})+\ldots+b_0f(x_{i},t_{i}) \right)
	\label{eq:mehrschritt}
\end{align}
mit $a_k \neq 0$. Durch $\mathcal{Z}$-Transformation (nachzulesen zum Beispiel in \cite{digit}) lassen sich Mehrschrittverfahren auch durch die beiden Polynome
\begin{align}
	P_a(z) &= a_0 + a_1 z + \ldots + a_k z^k \\
	P_b(z) &= b_0 + b_1 z + \ldots + b_k z^k	
\end{align}
charakterisieren.
Ein Beispiel f�r ein Mehrschrittverfahren ist das (implizite) Milne-Simpson-Verfahren:
\begin{align}	
	x_{i+1} = x_{i-1} + \frac{T}{3} \left(f(x_{i+1},t_{i+1})  +4f(x_i,t_i)  +f(x_{i-1},t_{i-1})  \right).
\end{align}	

%%%%%%%%%%%%%%%%%%%%%%%%%%%%%%%%%%%%%%%%%%%%
\subsection{Adams-Verfahren/Pr�diktor-Korrektor-Verfahren}
Wie auch bei den Einschrittverfahren k�nnen Mehrschrittverfahren in explizite und implizite Verfahren unterteilt werden. In \cite{faires} wird gezeigt, dass implizite Verfahren im allgemeinen erheblich bessere Ergebnisse erzielen, als explizite Verfahren, jedoch das Aufl�sen der impliziten Gleichung nach dem gesuchten Funktionswert nicht realisierbar sein kann oder zumindest mit erheblichem Rechenaufwand verbunden ist.\\
Dies f�hrt -- �hnlich wie bei dem Heun-Verfahren -- zu dem Gedanken, Approximationen, welche durch ein explizites Verfahren gewonnen wurden, durch ein implizites Verfahren zu verbessern. Diese Kombination aus explizitem und implizitem Mehrschrittverfahren wird als Pr�diktor-Korrektor-Verfahren bezeichnet. 

Die hierf�r am h�ufigsten verwendeten Mehrschrittverfahren sind die so genannten Adams-Verfahren. Dabei werden die expliziten Verfahren als \textit{Adams-Bashforth-Verfahren} und die impliziten Verfahren als \textit{Adams-Moulton-Verfahren} bezeichnet. Die Koeffizienten dieser Verfahren k�nnen in Anhang \ref{sec:anhmehr} nachgeschlagen werden.\\
Nach \cite{faires} liegen Rechenaufwand und Genauigkeit bei einem $m$-stufigen \textit{Adams-Bashforth-Verfahren} und einem $m$-1-stufigen \textit{Adams-Moulton-Verfahren} im gleichen Gr��enbereich, wobei der lokale Fehler in beiden F�llen von $T^{m+1}$ abh�ngt. Wie bereits eingangs erw�hnt, liefern die impliziten Verfahren jedoch deutlich genauere Approximationen.

Zur Durchf�hrung des Pr�diktor-Korrektor-Verfahrens wird zun�chst mit dem expliziten $m$-stufigen \textit{Adams-Bashforth-Verfahren} ein Approximationswert $x_{i+1}^{(0)}$ berechnet, also die Pr�diktion durchgef�hrt. Dieser wird dann in die rechte Seite der impliziten Gleichung des \textit{Adams-Moulton-Verfahren} eingesetzt, wodurch eine bessere Sch�tzung $x_{i+1}^{(1)}$ berechnet werden kann (Korrektur). Diese wird nun als Funktionswert $x(t_{i+1})$ angenommen.
In \cite{netomac} wird beschrieben, dass der Korrektor-Schritt in der Praxis iterativ wiederholt werden kann, um eine noch h�here Genauigkeit zu erzielen.

Da ein $k$-stufiges Mehrschrittverfahren zur Berechnung des Funktionswertes $x_i$ alle zur�ckliegenden Funktionswerte bis $x_{i-k}$ ben�tigt, ist der erste �berhaupt berechenbare Funktionswert $x(t=t_k)$. Aus diesem Grund m�ssen Mehrschrittverfahren durch geeignete Einschrittverfahren \glqq gestartet\grqq \ werden. Man bezeichnet Einschrittverfahren im Gegensatz zu Mehrschrittverfahren aus diesem Grund auch als \textit{selbststartend}.

\subsection{BDF-Verfahren}
Die zweite sehr bedeutende Klasse der Mehrschrittverfahren sind die so genannten \textit{Backwards Differentiation Formulas} (BDF). F�r diese Verfahren gilt $P_b(z) = z^k$. Damit lassen sich eine Reihe impliziter Verfahren erzeugen, die besonders gut f�r steife Differentialgleichungen geeignet sind. Die Bedeutung dessen wird im folgenden Kapitel erl�utert.

%
%\subsection{Schrittweitensteuerung}
%(WAHRSCHEINLICH WEGLASSEN!)
%Die Methoden zur Schrittweitensteuerung lassen sich weitgehend von Einschrittverfahren �bertragen. Bei Mehrschrittverfahren besteht lediglich das Problem, dass die Schrittweite in letzten $k$ Schritten �bereinstimmen muss, damit Gleichung (\ref{eq:mehrschritt}) sinnvoll ausgewertet werden kann. 
%\subsection{Extrapolationsverfahren}
%(wahrscheinlich weglassen, gehoert zu den Mehrschrittverfahren!!)
%TABELLE MIT VERFAHREN, KONSISTENZORDNUNG und STABILIT�T!!!!

% Nach Friedrich: Numerische Methoden:
%\section{Beschreibung der Integrationsverfahren}
%	\subsection{Explizite Einschrittverfahren}
%		\subsubsection{Explizites Eulerverfahren}
%		\subsubsection{Trapezregel}
%		\subsubsection{Runge-Kutta-Verfahren}
%	\subsection{Implizite Einschrittverfahren}
%	\subsection{Impliztes Eulerverfahren}
%	\subsection{explizite Mehrschrittverfahren}
%	\subsection{implizite Mehrschrittverfahren}
%	\subsection{Pr�diktor-Korrektor-Verfahren}

% Nach Faires: Numerische Methoden>
% \subsection{Eulersches Verfahren}
% \subsection{Runge-Kutta Verfahren}
% \subsubsection{Mittelpunktmethode}
% \subsubsection{modifiziertes Eulersches Verfahren $\hat{=}$ implizites Eulerverfahren}
% \subsubsection{Heunsches Verfahren}
% \subsubsection{Runge-Kutta Verfahren 4. Ordnung}
% \subsection{(Pr�diktor-Korrektor Verfahren )$\hat{=}$Mehrschrittverfahren}
% --> Unterscheidung von explizit und implizit (????)
% \subsubsection{Adam-Bashforth ...(explizit)}
% \subsubsection{Adams-Moulton ....(implizit)}
% \subsubsection{Pr�diktor-Korrektor Verfahren: Milnesche, Simpsonsche}
% Ausblick:(?) Extrapolationsverfahren und Adaptive Verfahren


\chapter{Stabilit�tsanalyse der wichtigsten Methoden}
\label{sec:stabi}
Neben der Konvergenz des Fehlers mit der Schrittweite $e(t) = \mathcal{O}(T^p)$ ist f�r die Anwendbarkeit eines Verfahrens entscheidend, welche Voraussetzungen an die Schrittweite gestellt werden m�ssen, damit das Verfahren bei der L�sung von Differentialgleichungen stabil ist. Unter der Stabilit�t soll in diesem Zusammenhang die exponentielle Stabilit�t verstanden werden, wie sie in \cite{gruene} beschrieben wird. Zur Einfachheit werden die Betrachtungen an dieser Stelle auf die Klasse \textit{linearer zeitinvarianter Systeme} beschr�nk, f�r die gilt:
\begin{align}
	\dot{x}=Ax.
\end{align}
Aus der Systemtheorie ist bekannt, dass ein solches System genau dann exponentiell stabil ist, wenn f�r alle Eigenwerte der Matrix $A$ gilt: $\Re \lambda_i < 0$, das hei�t, die Eigenwerte liegen in der linken Halbebene der komplexen Zahlenebene $\mathbb{C}^-$.

\section{Steife Differentialgleichungen}
Zur Einf�hrung soll das Anfangswertproblem
\begin{align}
	\dot{x}=\lambda x,\quad x(0)=x_0
	\label{eq:bsp}
\end{align}
betrachtet werden. F�r diese einfache Differentialgleichung lautet die geschlossene L�sung $x(t)= e^{\lambda t} x_0$. Das hei�t f�r $\lambda < 0$ konvergiert diese L�sung gegen $0$.
Eine Approximation des Systems mit dem expliziten Euler-Verfahren liefert
\begin{align}
	x_{i+1}=x_i + T \lambda x_i = (1+T\lambda)x_i.
\end{align}
Diese Differenzengleichung konvergiert genau dann ebenfalls gegen $0$, wenn 
\begin{align}
	|1+T\lambda| < 1
\end{align}
gilt. F�r die Schrittweite $T$ muss also die Bedingung
\begin{align}
	T < \frac{2}{|\lambda|}
	\label{eq:Tbed}
\end{align}
erf�llt sein, damit die Approximation eine Sinnvolle N�herung darstellt. Es ist augenscheinlich, dass dies bei sehr schnellen Systemen mit $\Re \lambda \ll 0$ zu sehr kleinen notwendigen Schrittweiten f�hrt.\\
Aus diesem Grund werden Differentialgleichungen mit sehr schnellen stabilen Eigenwerten $\Re \lambda \ll 0$ allgemein als \textit{steif} bezeichnet. F�r den Fall, dass das Anfangswertproblem ein Differentialgleichungssystem mit mehreren Eigenwerten darstellt spricht man von einem \textit{steifen} System, falls mindestens ein solcher Eigenwert vorliegt. Die Approximation von steifen Differentialgleichungen f�hrt vor allem bei expliziten Verfahren zu Problemen und stellt die wesentliche Motivation zur Verwendung von impliziten Verfahren dar.
In Abbildung \ref{fig:example} ist die Approximation von Gleichung (\ref{eq:bsp}) mit dem expliziten Euler-Verfahren f�r verschiedene Schrittweiten dargestellt. 
\begin{figure}
\centering
% This file was created by matlab2tikz v0.4.3.
% Copyright (c) 2008--2013, Nico Schlömer <nico.schloemer@gmail.com>
% All rights reserved.
% 
% The latest updates can be retrieved from
%   http://www.mathworks.com/matlabcentral/fileexchange/22022-matlab2tikz
% where you can also make suggestions and rate matlab2tikz.
% 
%
% defining custom colors
\definecolor{mycolor1}{rgb}{0,1,1}%
\definecolor{mycolor2}{rgb}{1,0,1}%
%
\begin{tikzpicture}

\begin{axis}[%
width=0.7\textwidth,
height=5cm,
scale only axis, % Angegebene Groesse bezieht sich auf die Achsen!!
xmin=0,
xmax=10,
xtick={0,2},
xticklabels={0,$\frac{2}{|\lambda|}$},
xmajorgrids,
ymin=-1.5,
ymax=1.5,
ytick={0,1},
yticklabels={0,$x_0$},
ymajorgrids,
axis x line*=bottom,
axis y line*=left,
legend style={draw=black,legend cell align=left,anchor=south east,at={(1,0)}},
xlabel ={$t, i$}
]
\addplot [
color=black,
dashed,
forget plot
]
table[row sep=crcr]{
0 1\\
0.2 0.818730773333333\\
0.4 0.670320079202998\\
0.6 0.548811676826732\\
0.8 0.449329008582714\\
1 0.367879486678025\\
1.2 0.301194256621369\\
1.4 0.246597006647172\\
1.6 0.201896557953924\\
1.8 0.165298925026955\\
2 0.135335316718487\\
2.2 0.110803188516239\\
2.4 0.0907179802216991\\
2.6 0.0742736021021498\\
2.8 0.0608100836873454\\
3 0.049787086843805\\
3.2 0.0407622201136423\\
3.4 0.0333732839964259\\
3.6 0.0273237346150667\\
3.8 0.0223707823717484\\
4 0.0183156479512932\\
4.2 0.0149955846112634\\
4.4 0.0122773465853651\\
4.6 0.0100518414643173\\
4.8 0.0082297519355046\\
5 0.00673795116649717\\
5.2 0.00551656796922847\\
5.4 0.00451658395959232\\
5.6 0.00369786627806195\\
5.8 0.00302755691752091\\
6 0.00247875401639258\\
6.2 0.0020294321927442\\
6.4 0.00166155858859302\\
6.6 0.00136036914817741\\
6.8 0.0011137760847061\\
7 0.000911882755151595\\
7.2 0.000746586473314596\\
7.4 0.000611253320657065\\
7.6 0.000500451903924127\\
7.8 0.000409735374315939\\
8 0.000335462959875712\\
8.2 0.000274653848563731\\
8.4 0.000224867557833559\\
8.6 0.000184105989522648\\
8.8 0.000150733239177176\\
9 0.000123409941478568\\
9.2 0.000101039516823769\\
9.4 8.27241617463509e-005\\
9.6 6.77288169199417e-005\\
9.8 5.54516666538156e-005\\
10 4.53999859221005e-005\\
}node(vergleich)[pos=0.15, inner sep =0pt]{};
%\addlegendentry{\scriptsize Vergleich};
\node(vergleichtext) at ($(vergleich)+(10,100)$){\scriptsize{Vergleichsl�sung}};
\draw[black](vergleich)--(vergleichtext);
\addplot [
color=mycolor1,
solid
]
table[row sep=crcr]{
0 1\\
0.1 0.9\\
0.2 0.81\\
0.3 0.729\\
0.4 0.6561\\
0.5 0.59049\\
0.6 0.531441\\
0.7 0.4782969\\
0.8 0.43046721\\
0.9 0.387420489\\
1 0.3486784401\\
1.1 0.31381059609\\
1.2 0.282429536481\\
1.3 0.2541865828329\\
1.4 0.22876792454961\\
1.5 0.205891132094649\\
1.6 0.185302018885184\\
1.7 0.166771816996666\\
1.8 0.150094635296999\\
1.9 0.135085171767299\\
2 0.121576654590569\\
2.1 0.109418989131512\\
2.2 0.0984770902183611\\
2.3 0.088629381196525\\
2.4 0.0797664430768725\\
2.5 0.0717897987691852\\
2.6 0.0646108188922667\\
2.7 0.05814973700304\\
2.8 0.052334763302736\\
2.9 0.0471012869724624\\
3 0.0423911582752162\\
3.1 0.0381520424476946\\
3.2 0.0343368382029251\\
3.3 0.0309031543826326\\
3.4 0.0278128389443693\\
3.5 0.0250315550499324\\
3.6 0.0225283995449392\\
3.7 0.0202755595904453\\
3.8 0.0182480036314007\\
3.9 0.0164232032682607\\
4 0.0147808829414346\\
4.1 0.0133027946472911\\
4.2 0.011972515182562\\
4.3 0.0107752636643058\\
4.4 0.00969773729787523\\
4.5 0.00872796356808771\\
4.6 0.00785516721127894\\
4.7 0.00706965049015105\\
4.8 0.00636268544113594\\
4.9 0.00572641689702235\\
5 0.00515377520732011\\
5.1 0.0046383976865881\\
5.2 0.00417455791792929\\
5.3 0.00375710212613636\\
5.4 0.00338139191352273\\
5.5 0.00304325272217045\\
5.6 0.00273892744995341\\
5.7 0.00246503470495807\\
5.8 0.00221853123446226\\
5.9 0.00199667811101603\\
6 0.00179701029991443\\
6.1 0.00161730926992299\\
6.2 0.00145557834293069\\
6.3 0.00131002050863762\\
6.4 0.00117901845777386\\
6.5 0.00106111661199647\\
6.6 0.000955004950796825\\
6.7 0.000859504455717143\\
6.8 0.000773554010145428\\
6.9 0.000696198609130885\\
7 0.000626578748217797\\
7.1 0.000563920873396017\\
7.2 0.000507528786056415\\
7.3 0.000456775907450774\\
7.4 0.000411098316705696\\
7.5 0.000369988485035127\\
7.6 0.000332989636531614\\
7.7 0.000299690672878453\\
7.8 0.000269721605590607\\
7.9 0.000242749445031547\\
8 0.000218474500528392\\
8.1 0.000196627050475553\\
8.2 0.000176964345427998\\
8.3 0.000159267910885198\\
8.4 0.000143341119796678\\
8.5 0.00012900700781701\\
8.6 0.000116106307035309\\
8.7 0.000104495676331778\\
8.8 9.40461086986004e-005\\
8.9 8.46414978287404e-005\\
9 7.61773480458663e-005\\
9.1 6.85596132412797e-005\\
9.2 6.17036519171517e-005\\
9.3 5.55332867254366e-005\\
9.4 4.99799580528929e-005\\
9.5 4.49819622476036e-005\\
9.6 4.04837660228433e-005\\
9.7 3.64353894205589e-005\\
9.8 3.2791850478503e-005\\
9.9 2.95126654306527e-005\\
10 2.65613988875875e-005\\
};
\addlegendentry{\scriptsize $T=0.1/|\lambda|$};

\addplot [
color=green,
solid
]
table[row sep=crcr]{
0 1\\
0.5 0.5\\
1 0.25\\
1.5 0.125\\
2 0.0625\\
2.5 0.03125\\
3 0.015625\\
3.5 0.0078125\\
4 0.00390625\\
4.5 0.001953125\\
5 0.0009765625\\
5.5 0.00048828125\\
6 0.000244140625\\
6.5 0.0001220703125\\
7 6.103515625e-005\\
7.5 3.0517578125e-005\\
8 1.52587890625e-005\\
8.5 7.62939453125e-006\\
9 3.814697265625e-006\\
9.5 1.9073486328125e-006\\
10 9.5367431640625e-007\\
};
\addlegendentry{\scriptsize $T=0.5/|\lambda|$};

\addplot [
color=blue,
solid, thick
]
table[row sep=crcr]{
0 1\\
1 0\\
2 0\\
3 0\\
4 0\\
5 0\\
6 0\\
7 0\\
8 0\\
9 0\\
10 0\\
};
\addlegendentry{\scriptsize $T=1/|\lambda|$};

\addplot [
color=red,
solid
]
table[row sep=crcr]{
0 1\\
2 -1\\
4 1\\
6 -1\\
8 1\\
10 -1\\
};
\addlegendentry{\scriptsize $T=2/|\lambda|$};

\addplot [
color=mycolor2,
solid
]
table[row sep=crcr]{
0 1\\
2.1 -1.1\\
4.2 1.21\\
6.3 -1.331\\
8.4 1.4641\\
};
\addlegendentry{\scriptsize $T=2.1/|\lambda|$};

\end{axis}
\end{tikzpicture}%
\caption{Simulation des Beispielsystems mit explizitem Euler-Verfahren}
\label{fig:example}
\end{figure}

\section{Einschrittverfahren}
Das Beispiel der steifen Differentialgleichungen hat gezeigt, dass durch die Wahl der Schrittweite im allgemeinen das Stabilit�tsverhalten des Verfahrens beeinflussen l�sst. Im Folgenden sollen Bedingungen aufgezeigt werden, die an die Schrittweite gestellt werden m�ssen, um die Stabilit�t eines Verfahrens allgemein zu gew�hrleisten.

Nach \cite{gruene} lassen sich alle Runge-Kutta-Verfahren mit konstanter Schrittweite, also alle bisher betrachteten Einschrittverfahren ohne Schrittweitensteuerung durch die Differenzengleichung
\begin{align}
	x_{i+1}= \tilde{A} x_i 
\end{align}
beschreiben.
Nach \cite{digit} ist eine solche Differenzengleichung genau dann exponentiell stabil, wenn f�r alle Eigenwerte der Matrix $\tilde{A}$ stabil sind, das hei�t $|\tilde{\lambda}_i|<1$ gilt. Die Eigenwerte m�ssen also im Einheitskreis $B_1(0)$ der komplexen Zahlenebene liegen. ($B_u(v)$ bezeichnet im Folgenden allgemein einen Kreis mit Radius $u$ um den Punkt $v$ in der komplexen Zahlenebene.)
Zur Untersuchung der Stabilit�t eines Verfahrens m�sste also die Matrix $\tilde{A}$ aus den Koeffizienten des Verfahrens und der Matrix $A$ der zu Grunde liegenden Differentialgleichung berechnet werden.
Die Matrix der Differenzengleichung $\tilde{A}$ berechnet sich also als Funktion der Art
\begin{align}
	\tilde{A}= R(TA).
\end{align}
Es kann gezeigt werden, dass diese Funktion f�r explizite Werte $z \in \mathbb{C}$ die Form
\begin{align}
	R(z) = 1 + zb^T \left( I-z\mathcal{A}  \right)^{-1} e
\end{align}
annimmt, wobei $I$ f�r die Einheitsmatrix steht, $b = (b_1, \ldots, b_s)^T$ und $\mathcal{A}=(a_{ij})_{i,j=1,\ldots,s}$ die Koeffizienten des Verfahrens beinhalten und $e=(1,\ldots,1)^T \in \mathbb{R}^s$ gilt. $R(z)$ wird als die Stabilit�tsfunktion des Verfahrens bezeichnet. Die Stabilit�tsfunktion an sich ist nun unabh�ngig von der Schrittweite und der Dynamik der betrachteten Differentialgleichung und kann damit noch keine endg�ltige Aussage �ber die Stabilit�t des Verfahrens treffen. 
Um diese zu erm�glichen, gilt es herauszufinden, welche Bedingungen $A$ und $T$ erf�llen m�ssen, damit die Eigenwerte der Matrix $\tilde{A}$ stabil sind. 
Hierf�r wird das Stabilit�tsgebiet $S \subset \mathbb{C}$ eines Runge-Kutta-Verfahrens mit Stabilit�tsfunktion $R$ definiert als die Menge von Eigenwerten $\lambda_i$, die $TA$ annehmen darf, damit $\tilde{A}=R(TA)$ exponentiell stabil ist, das hei�t, f�r den Betrag der Stabilit�tsfunktion gilt
\begin{align}
	S = \{z\in\mathbb{C} \big| |R(z)|<1 \}.
\end{align}
F�r die Stabilit�t des Verfahrens m�ssen also alle Eigenwerte $T \lambda_i$ der Matrix $\tilde{A}$ im Stabilit�tsgebiet liegen.
Dabei gilt zu beachten, dass f�r die Menge aller Eigenwerte $\Sigma$ zweier Matrizen gilt: \mbox{$\lambda_i \in \Sigma(A) \Rightarrow T\lambda_i \in \Sigma(TA)$}.\\
Damit l�sst sich die Schrittweite $T$ in Abh�ngigkeit der Eigenwerte von $A$ so berechnen, dass das verwendete Verfahren stabil ist. Allgemein gilt, je gr��er das Stabilit�tsgebiet $S$, desto gr��er kann die Schrittweite $T$ gew�hlt werden, um Stabilit�t zu gew�hrleisten. Ein Verfahren wird als $A$-stabil bezeichnet, wenn es f�r eine stabile Dynamik $A$ unabh�ngig von der Schrittweite stabil ist. In diesem Fall ist die gesamte linke Halbebene der komplexen Zahlenebene im Stabilit�tsgebiet enthalten: $\mathbb{C}^- \subseteq S$.

Zur Vergleichbarkeit sollen nun die betrachteten Verfahren auf das Stabilit�tsgebiet untersucht werden.
F�r das explizite Euler-Verfahren lautet die Stabilit�tsfunktion $R_\mathrm{Eu,ex}(z) = 1+z$. Damit ergibt sich die Ungleichung
\begin{align}
	|R_\mathrm{Eu,ex}(z)| <1  \ \Leftrightarrow \ |1+z|<1,
\end{align}
also $S = \{ z\in\mathbb{C} \big| |1+z|<1\} = B_1(-1)$. Die Schrittweite muss also so gew�hlt werden, dass $T\lambda_i \in B_1(-1)$ gilt. F�r das betrachtete Beispielsystem entspricht dies genau (\ref{eq:Tbed}).

F�r das implizite Euler-Verfahren erh�lt man $R_\mathrm{Eu,im}(z) = \frac{1}{1-z}$. F�r das Stabilit�tsgebiet bedeutet das
\begin{align}
	|R_\mathrm{Eu,im}(z)| <1  \ \Leftrightarrow \ |1-z|>1,
\end{align}
das hei�t $S$ ist $\mathbb{C}$ ohne den Kreis $B_1(1)$. Damit ist das implizite Euler-Verfahren -- wie viele implizite Verfahren -- $A$-stabil, was hilfreich bei der Simulation von stabilen Systemen ist. Allerdings enth�lt das Stabilit�tsgebiet auch Bereiche, in denen die Approximation stabil ist, obwohl das eigentliche System instabil ist, wie zum Beispiel $\Re z > 2$. In diesem Fall w�rde die Instabilit�t eines Systems in der Simulation m�glicherweise unentdeckt bleiben, was eines der Probleme des impliziten Euler-Verfahrens ist und eine Motivation f�r die Trapezregel liefert. 

F�r diese ergibt sich als Sabilt�tsfunktion $R_\mathrm{Trapez}(z)=\frac{z+2}{z-2}$. Das Stabilit�tsgebiet entspricht damit genau der linken Halbebene $S_\mathrm{Trapez}=\mathbb{C}^-$. Damit ist die Trapezregel ebenfalls $A$-stabil, enth�lt jedoch keine Bereiche, in denen ein instabiles System stabil approximiert werden w�rde. Man bezeichnet diese Tatsache mit 
\textit{Erhaltung der Isometrie}, das hei�t, die Approximation ist genau dann stabil, wenn auch das zu Grunde liegende System exponentiell stabil ist. Damit besitzt die Trapezregel quasi ideale Stabilit�tseigenschaften. \\
Aufgrund der schwierigen Realisierbarkeit der impliziten Verfahren wurde das Heun-Verfahren auf Basis der Trapezregel und dem expliziten Euler-Verfahren eingef�hrt. Hierf�r ergibt sich die Stabilit�tsfunktion $R_\mathrm{Heun}(z) = \frac{1}{2}z^2+z+1$. Das zugeh�rige Stabilit�tsgebiet ist in Abbildung \ref{fig:stabis} zu sehen. Vergleicht man die Stabilit�tsgebiete von explizitem Euler-Verfahren und Heun-Verfahren, so wird deutlich, dass das Stabilit�tsgebiet des Heun-Verfahrens nur wenig gr��er ist, als das des expliziten Euler-Verfahrens. Zwischen den Stabilit�tsgebieten von Heun-Verfahren und Trapezregel besteht hingegen ein deutlicher Gr��enunterschied.

Ein gr��eres Stabilit�tsgebiet kann zum Beispiel mit dem \textit{klassischen} Runge-Kutta-Verfahren mit der Stabilit�tsfunktion $R_{RK}(z)=\frac{1}{24}z^4+\frac{1}{6}z^3+\eh z^2+z+1$ erreicht werden. �hnlich, wie bei den impliziten Verfahren treten hier jedoch kleine Bereiche auf, in denen ein instabiles System zu einer stabilen Approximation f�hren w�rde.
Grunds�tzlich gilt f�r explizite Runge-Kutta-Verfahren: Je H�her die Ordnung des Verfahrens, desto gr��er das Stabilit�tsgebiet. Bei der Betrachtung der Stabilit�tsgebiete wird auch der gro�e Vorteil der impliziten Verfahren gegen�ber expliziten Verfahren deutlich. Zwar gibt es auch implizite Verfahren, die nur kleine Stabilit�tsgebiete besitzen und damit schlecht f�r die L�sung steifer Differentialgleichungen geeignet sind, jedoch sind die Stabilit�tsgebiete von expliziten Verfahren im Allgemeinen kleiner, als die der impliziten Verfahren.
\begin{figure}
\centering
% This file was created by matlab2tikz v0.4.3.
% Copyright (c) 2008--2013, Nico Schlömer <nico.schloemer@gmail.com>
% All rights reserved.
% 
% The latest updates can be retrieved from
%   http://www.mathworks.com/matlabcentral/fileexchange/22022-matlab2tikz
% where you can also make suggestions and rate matlab2tikz.
% 
%
% defining custom colors
\definecolor{mycolor1}{rgb}{1,0,1}%
%
\begin{tikzpicture}

\begin{axis}[%
width=4.52083333333333in,
height=6.02777777777778in,
area legend,
colormap/blackwhite,
unbounded coords=jump,
scale only axis,
xmin=-4,
xmax=2,
xmajorgrids,
ymin=-4,
ymax=4,
ymajorgrids,
axis x line*=bottom,
axis y line*=left
]

\addplot[line width=2.0pt,draw=black,forget plot,fill=black, fill opacity =0.2]
table[row sep=crcr]{
x y\\
-1.9 0.499999990000001 \\
-1.89999999 0.5 \\
-1.8 0.59999999 \\
-1.79999999 0.6 \\
-1.79999999 0.699999999999999 \\
-1.7 0.79999999 \\
-1.6 0.79999999 \\
-1.59999999 0.8 \\
-1.5 0.89999999 \\
-1.49999999 0.899999999999999 \\
-1.4 0.999999990000001 \\
-1.3 0.999999990000001 \\
-1.2 0.999999990000001 \\
-1.1 0.999999990000001 \\
-1 0.999999990000001 \\
-0.9 0.999999990000001 \\
-0.8 0.999999990000001 \\
-0.7 0.999999990000001 \\
-0.6 0.999999990000001 \\
-0.50000001 0.899999999999999 \\
-0.5 0.89999999 \\
-0.40000001 0.8 \\
-0.4 0.79999999 \\
-0.3 0.79999999 \\
-0.20000001 0.699999999999999 \\
-0.20000001 0.6 \\
-0.2 0.59999999 \\
-0.10000001 0.5 \\
-0.1 0.499999990000001 \\
-1.0000000116861e-008 0.399999999999999 \\
-1.0000000116861e-008 0.3 \\
-1.0000000116861e-008 0.199999999999999 \\
-1.0000000116861e-008 0.0999999999999996 \\
-1.0000000116861e-008 0 \\
-1.0000000116861e-008 -0.0999999999999996 \\
-1.0000000116861e-008 -0.199999999999999 \\
-1.0000000116861e-008 -0.3 \\
-1.0000000116861e-008 -0.399999999999999 \\
-0.1 -0.49999999 \\
-0.10000001 -0.5 \\
-0.2 -0.599999989999999 \\
-0.20000001 -0.6 \\
-0.20000001 -0.699999999999999 \\
-0.3 -0.79999999 \\
-0.30000001 -0.8 \\
-0.4 -0.899999989999999 \\
-0.5 -0.899999989999999 \\
-0.50000001 -0.899999999999999 \\
-0.6 -0.99999999 \\
-0.7 -0.99999999 \\
-0.8 -0.99999999 \\
-0.9 -0.99999999 \\
-1 -0.99999999 \\
-1.1 -0.99999999 \\
-1.2 -0.99999999 \\
-1.3 -0.99999999 \\
-1.4 -0.99999999 \\
-1.49999999 -0.899999999999999 \\
-1.5 -0.899999989999999 \\
-1.6 -0.899999989999999 \\
-1.69999999 -0.8 \\
-1.7 -0.79999999 \\
-1.79999999 -0.699999999999999 \\
-1.8 -0.699999989999999 \\
-1.89999999 -0.6 \\
-1.89999999 -0.5 \\
-1.9 -0.49999999 \\
-1.99999999 -0.399999999999999 \\
-1.99999999 -0.3 \\
-1.99999999 -0.199999999999999 \\
-1.99999999 -0.0999999999999996 \\
-1.99999999 0 \\
-1.99999999 0.0999999999999996 \\
-1.99999999 0.199999999999999 \\
-1.99999999 0.3 \\
-1.99999999 0.399999999999999 \\
-1.9 0.499999990000001 \\
};


%\addplot[dashed,line width=2.0pt,draw=black,forget plot]
%table[row sep=crcr]{
%x y\\
%-1.9 0.490909090909091 \\
%-1.89090909090909 0.5 \\
%-1.8 0.59090909090909 \\
%-1.79090909090909 0.6 \\
%-1.79090909090909 0.699999999999999 \\
%-1.7 0.79090909090909 \\
%-1.6 0.79090909090909 \\
%-1.59090909090909 0.8 \\
%-1.5 0.89090909090909 \\
%-1.49090909090909 0.899999999999999 \\
%-1.4 0.990909090909091 \\
%-1.3 0.990909090909091 \\
%-1.2 0.990909090909091 \\
%-1.1 0.990909090909091 \\
%-1 0.990909090909091 \\
%-0.9 0.990909090909091 \\
%-0.8 0.990909090909091 \\
%-0.7 0.990909090909091 \\
%-0.6 0.990909090909091 \\
%-0.509090909090909 0.899999999999999 \\
%-0.5 0.89090909090909 \\
%-0.40909090909091 0.8 \\
%-0.4 0.79090909090909 \\
%-0.3 0.79090909090909 \\
%-0.20909090909091 0.699999999999999 \\
%-0.20909090909091 0.6 \\
%-0.2 0.59090909090909 \\
%-0.109090909090909 0.5 \\
%-0.1 0.490909090909091 \\
%-0.00909090909090936 0.399999999999999 \\
%-0.00909090909090936 0.3 \\
%-0.00909090909090936 0.199999999999999 \\
%-0.00909090909090936 0.0999999999999996 \\
%-0.00909090909090936 0 \\
%-0.00909090909090936 -0.0999999999999996 \\
%-0.00909090909090936 -0.199999999999999 \\
%-0.00909090909090936 -0.3 \\
%-0.00909090909090936 -0.399999999999999 \\
%-0.1 -0.490909090909091 \\
%-0.109090909090909 -0.5 \\
%-0.2 -0.59090909090909 \\
%-0.20909090909091 -0.6 \\
%-0.20909090909091 -0.699999999999999 \\
%-0.3 -0.79090909090909 \\
%-0.30909090909091 -0.8 \\
%-0.4 -0.89090909090909 \\
%-0.5 -0.89090909090909 \\
%-0.509090909090909 -0.899999999999999 \\
%-0.6 -0.990909090909091 \\
%-0.7 -0.990909090909091 \\
%-0.8 -0.990909090909091 \\
%-0.9 -0.990909090909091 \\
%-1 -0.990909090909091 \\
%-1.1 -0.990909090909091 \\
%-1.2 -0.990909090909091 \\
%-1.3 -0.990909090909091 \\
%-1.4 -0.990909090909091 \\
%-1.49090909090909 -0.899999999999999 \\
%-1.5 -0.89090909090909 \\
%-1.6 -0.89090909090909 \\
%-1.69090909090909 -0.8 \\
%-1.7 -0.79090909090909 \\
%-1.79090909090909 -0.699999999999999 \\
%-1.8 -0.69090909090909 \\
%-1.89090909090909 -0.6 \\
%-1.89090909090909 -0.5 \\
%-1.9 -0.490909090909091 \\
%-1.99090909090909 -0.399999999999999 \\
%-1.99090909090909 -0.3 \\
%-1.99090909090909 -0.199999999999999 \\
%-1.99090909090909 -0.0999999999999996 \\
%-1.99090909090909 0 \\
%-1.99090909090909 0.0999999999999996 \\
%-1.99090909090909 0.199999999999999 \\
%-1.99090909090909 0.3 \\
%-1.99090909090909 0.399999999999999 \\
%-1.9 0.490909090909091 \\
%NaN NaN \\
%};


%\addplot[dashed,line width=2.0pt,draw=black,forget plot]
%table[row sep=crcr]{
%x y\\
%-1.9 0.481818181818181 \\
%-1.88181818181818 0.5 \\
%-1.8 0.581818181818181 \\
%-1.78181818181818 0.6 \\
%-1.78181818181818 0.699999999999999 \\
%-1.7 0.781818181818181 \\
%-1.6 0.781818181818181 \\
%-1.58181818181818 0.8 \\
%-1.5 0.881818181818181 \\
%-1.48181818181818 0.899999999999999 \\
%-1.4 0.981818181818181 \\
%-1.3 0.981818181818181 \\
%-1.2 0.981818181818181 \\
%-1.1 0.981818181818181 \\
%-1 0.981818181818181 \\
%-0.9 0.981818181818181 \\
%-0.8 0.981818181818181 \\
%-0.7 0.981818181818181 \\
%-0.6 0.981818181818181 \\
%-0.518181818181818 0.899999999999999 \\
%-0.5 0.881818181818181 \\
%-0.418181818181818 0.8 \\
%-0.4 0.781818181818181 \\
%-0.3 0.781818181818181 \\
%-0.218181818181818 0.699999999999999 \\
%-0.218181818181818 0.6 \\
%-0.2 0.581818181818181 \\
%-0.118181818181818 0.5 \\
%-0.1 0.481818181818181 \\
%-0.018181818181818 0.399999999999999 \\
%-0.018181818181818 0.3 \\
%-0.018181818181818 0.199999999999999 \\
%-0.018181818181818 0.0999999999999996 \\
%-0.018181818181818 0 \\
%-0.018181818181818 -0.0999999999999996 \\
%-0.018181818181818 -0.199999999999999 \\
%-0.018181818181818 -0.3 \\
%-0.018181818181818 -0.399999999999999 \\
%-0.1 -0.481818181818182 \\
%-0.118181818181818 -0.5 \\
%-0.2 -0.581818181818182 \\
%-0.218181818181818 -0.6 \\
%-0.218181818181818 -0.699999999999999 \\
%-0.3 -0.781818181818182 \\
%-0.318181818181818 -0.8 \\
%-0.4 -0.881818181818182 \\
%-0.5 -0.881818181818182 \\
%-0.518181818181818 -0.899999999999999 \\
%-0.6 -0.981818181818182 \\
%-0.7 -0.981818181818182 \\
%-0.8 -0.981818181818182 \\
%-0.9 -0.981818181818182 \\
%-1 -0.981818181818182 \\
%-1.1 -0.981818181818182 \\
%-1.2 -0.981818181818182 \\
%-1.3 -0.981818181818182 \\
%-1.4 -0.981818181818182 \\
%-1.48181818181818 -0.899999999999999 \\
%-1.5 -0.881818181818182 \\
%-1.6 -0.881818181818182 \\
%-1.68181818181818 -0.8 \\
%-1.7 -0.781818181818182 \\
%-1.78181818181818 -0.699999999999999 \\
%-1.8 -0.681818181818181 \\
%-1.88181818181818 -0.6 \\
%-1.88181818181818 -0.5 \\
%-1.9 -0.481818181818182 \\
%-1.98181818181818 -0.399999999999999 \\
%-1.98181818181818 -0.3 \\
%-1.98181818181818 -0.199999999999999 \\
%-1.98181818181818 -0.0999999999999996 \\
%-1.98181818181818 0 \\
%-1.98181818181818 0.0999999999999996 \\
%-1.98181818181818 0.199999999999999 \\
%-1.98181818181818 0.3 \\
%-1.98181818181818 0.399999999999999 \\
%-1.9 0.481818181818181 \\
%NaN NaN \\
%};
%
%
%\addplot[dashed,line width=2.0pt,draw=black,forget plot]
%table[row sep=crcr]{
%x y\\
%-1.9 0.472727272727273 \\
%-1.87272727272727 0.5 \\
%-1.8 0.572727272727273 \\
%-1.77272727272727 0.6 \\
%-1.77272727272727 0.699999999999999 \\
%-1.7 0.772727272727273 \\
%-1.6 0.772727272727273 \\
%-1.57272727272727 0.8 \\
%-1.5 0.872727272727273 \\
%-1.47272727272727 0.899999999999999 \\
%-1.4 0.972727272727273 \\
%-1.3 0.972727272727273 \\
%-1.2 0.972727272727273 \\
%-1.1 0.972727272727273 \\
%-1 0.972727272727273 \\
%-0.9 0.972727272727273 \\
%-0.8 0.972727272727273 \\
%-0.7 0.972727272727273 \\
%-0.6 0.972727272727273 \\
%-0.527272727272727 0.899999999999999 \\
%-0.5 0.872727272727273 \\
%-0.427272727272728 0.8 \\
%-0.4 0.772727272727273 \\
%-0.3 0.772727272727273 \\
%-0.227272727272728 0.699999999999999 \\
%-0.227272727272728 0.6 \\
%-0.2 0.572727272727273 \\
%-0.127272727272727 0.5 \\
%-0.1 0.472727272727273 \\
%-0.0272727272727274 0.399999999999999 \\
%-0.0272727272727274 0.3 \\
%-0.0272727272727274 0.199999999999999 \\
%-0.0272727272727274 0.0999999999999996 \\
%-0.0272727272727274 0 \\
%-0.0272727272727274 -0.0999999999999996 \\
%-0.0272727272727274 -0.199999999999999 \\
%-0.0272727272727274 -0.3 \\
%-0.0272727272727274 -0.399999999999999 \\
%-0.1 -0.472727272727273 \\
%-0.127272727272727 -0.5 \\
%-0.2 -0.572727272727272 \\
%-0.227272727272728 -0.6 \\
%-0.227272727272728 -0.699999999999999 \\
%-0.3 -0.772727272727272 \\
%-0.327272727272728 -0.8 \\
%-0.4 -0.872727272727272 \\
%-0.5 -0.872727272727272 \\
%-0.527272727272727 -0.899999999999999 \\
%-0.6 -0.972727272727273 \\
%-0.7 -0.972727272727273 \\
%-0.8 -0.972727272727273 \\
%-0.9 -0.972727272727273 \\
%-1 -0.972727272727273 \\
%-1.1 -0.972727272727273 \\
%-1.2 -0.972727272727273 \\
%-1.3 -0.972727272727273 \\
%-1.4 -0.972727272727273 \\
%-1.47272727272727 -0.899999999999999 \\
%-1.5 -0.872727272727272 \\
%-1.6 -0.872727272727272 \\
%-1.67272727272727 -0.8 \\
%-1.7 -0.772727272727272 \\
%-1.77272727272727 -0.699999999999999 \\
%-1.8 -0.672727272727272 \\
%-1.87272727272727 -0.6 \\
%-1.87272727272727 -0.5 \\
%-1.9 -0.472727272727273 \\
%-1.97272727272727 -0.399999999999999 \\
%-1.97272727272727 -0.3 \\
%-1.97272727272727 -0.199999999999999 \\
%-1.97272727272727 -0.0999999999999996 \\
%-1.97272727272727 0 \\
%-1.97272727272727 0.0999999999999996 \\
%-1.97272727272727 0.199999999999999 \\
%-1.97272727272727 0.3 \\
%-1.97272727272727 0.399999999999999 \\
%-1.9 0.472727272727273 \\
%NaN NaN \\
%};
%
%
%\addplot[dashed,line width=2.0pt,draw=black,forget plot]
%table[row sep=crcr]{
%x y\\
%-1.9 0.463636363636364 \\
%-1.86363636363636 0.5 \\
%-1.8 0.563636363636364 \\
%-1.76363636363636 0.6 \\
%-1.76363636363636 0.699999999999999 \\
%-1.7 0.763636363636364 \\
%-1.6 0.763636363636364 \\
%-1.56363636363636 0.8 \\
%-1.5 0.863636363636364 \\
%-1.46363636363636 0.899999999999999 \\
%-1.4 0.963636363636364 \\
%-1.3 0.963636363636364 \\
%-1.2 0.963636363636364 \\
%-1.1 0.963636363636364 \\
%-1 0.963636363636364 \\
%-0.9 0.963636363636364 \\
%-0.8 0.963636363636364 \\
%-0.7 0.963636363636364 \\
%-0.6 0.963636363636364 \\
%-0.536363636363637 0.899999999999999 \\
%-0.5 0.863636363636364 \\
%-0.436363636363637 0.8 \\
%-0.4 0.763636363636364 \\
%-0.3 0.763636363636364 \\
%-0.236363636363637 0.699999999999999 \\
%-0.236363636363637 0.6 \\
%-0.2 0.563636363636364 \\
%-0.136363636363637 0.5 \\
%-0.1 0.463636363636364 \\
%-0.0363636363636367 0.399999999999999 \\
%-0.0363636363636367 0.3 \\
%-0.0363636363636367 0.199999999999999 \\
%-0.0363636363636367 0.0999999999999996 \\
%-0.0363636363636367 0 \\
%-0.0363636363636367 -0.0999999999999996 \\
%-0.0363636363636367 -0.199999999999999 \\
%-0.0363636363636367 -0.3 \\
%-0.0363636363636367 -0.399999999999999 \\
%-0.1 -0.463636363636363 \\
%-0.136363636363637 -0.5 \\
%-0.2 -0.563636363636363 \\
%-0.236363636363637 -0.6 \\
%-0.236363636363637 -0.699999999999999 \\
%-0.3 -0.763636363636363 \\
%-0.336363636363637 -0.8 \\
%-0.4 -0.863636363636363 \\
%-0.5 -0.863636363636363 \\
%-0.536363636363637 -0.899999999999999 \\
%-0.6 -0.963636363636363 \\
%-0.7 -0.963636363636363 \\
%-0.8 -0.963636363636363 \\
%-0.9 -0.963636363636363 \\
%-1 -0.963636363636363 \\
%-1.1 -0.963636363636363 \\
%-1.2 -0.963636363636363 \\
%-1.3 -0.963636363636363 \\
%-1.4 -0.963636363636363 \\
%-1.46363636363636 -0.899999999999999 \\
%-1.5 -0.863636363636363 \\
%-1.6 -0.863636363636363 \\
%-1.66363636363636 -0.8 \\
%-1.7 -0.763636363636363 \\
%-1.76363636363636 -0.699999999999999 \\
%-1.8 -0.663636363636363 \\
%-1.86363636363636 -0.6 \\
%-1.86363636363636 -0.5 \\
%-1.9 -0.463636363636363 \\
%-1.96363636363636 -0.399999999999999 \\
%-1.96363636363636 -0.3 \\
%-1.96363636363636 -0.199999999999999 \\
%-1.96363636363636 -0.0999999999999996 \\
%-1.96363636363636 0 \\
%-1.96363636363636 0.0999999999999996 \\
%-1.96363636363636 0.199999999999999 \\
%-1.96363636363636 0.3 \\
%-1.96363636363636 0.399999999999999 \\
%-1.9 0.463636363636364 \\
%NaN NaN \\
%};
%
%
%\addplot[dashed,line width=2.0pt,draw=black,forget plot]
%table[row sep=crcr]{
%x y\\
%-1.9 0.454545454545454 \\
%-1.85454545454545 0.5 \\
%-1.8 0.554545454545454 \\
%-1.75454545454545 0.6 \\
%-1.75454545454545 0.699999999999999 \\
%-1.7 0.754545454545454 \\
%-1.6 0.754545454545454 \\
%-1.55454545454545 0.8 \\
%-1.5 0.854545454545454 \\
%-1.45454545454545 0.899999999999999 \\
%-1.4 0.954545454545454 \\
%-1.3 0.954545454545454 \\
%-1.2 0.954545454545454 \\
%-1.1 0.954545454545454 \\
%-1 0.954545454545454 \\
%-0.9 0.954545454545454 \\
%-0.8 0.954545454545454 \\
%-0.7 0.954545454545454 \\
%-0.6 0.954545454545454 \\
%-0.545454545454545 0.899999999999999 \\
%-0.5 0.854545454545454 \\
%-0.445454545454546 0.8 \\
%-0.4 0.754545454545454 \\
%-0.3 0.754545454545454 \\
%-0.245454545454546 0.699999999999999 \\
%-0.245454545454546 0.6 \\
%-0.2 0.554545454545454 \\
%-0.145454545454545 0.5 \\
%-0.1 0.454545454545454 \\
%-0.0454545454545454 0.399999999999999 \\
%-0.0454545454545454 0.3 \\
%-0.0454545454545454 0.199999999999999 \\
%-0.0454545454545454 0.0999999999999996 \\
%-0.0454545454545454 0 \\
%-0.0454545454545454 -0.0999999999999996 \\
%-0.0454545454545454 -0.199999999999999 \\
%-0.0454545454545454 -0.3 \\
%-0.0454545454545454 -0.399999999999999 \\
%-0.1 -0.454545454545454 \\
%-0.145454545454545 -0.5 \\
%-0.2 -0.554545454545454 \\
%-0.245454545454546 -0.6 \\
%-0.245454545454546 -0.699999999999999 \\
%-0.3 -0.754545454545454 \\
%-0.345454545454546 -0.8 \\
%-0.4 -0.854545454545454 \\
%-0.5 -0.854545454545454 \\
%-0.545454545454545 -0.899999999999999 \\
%-0.6 -0.954545454545454 \\
%-0.7 -0.954545454545454 \\
%-0.8 -0.954545454545454 \\
%-0.9 -0.954545454545454 \\
%-1 -0.954545454545454 \\
%-1.1 -0.954545454545454 \\
%-1.2 -0.954545454545454 \\
%-1.3 -0.954545454545454 \\
%-1.4 -0.954545454545454 \\
%-1.45454545454545 -0.899999999999999 \\
%-1.5 -0.854545454545454 \\
%-1.6 -0.854545454545454 \\
%-1.65454545454545 -0.8 \\
%-1.7 -0.754545454545454 \\
%-1.75454545454545 -0.699999999999999 \\
%-1.8 -0.654545454545454 \\
%-1.85454545454545 -0.6 \\
%-1.85454545454545 -0.5 \\
%-1.9 -0.454545454545454 \\
%-1.95454545454545 -0.399999999999999 \\
%-1.95454545454545 -0.3 \\
%-1.95454545454545 -0.199999999999999 \\
%-1.95454545454545 -0.0999999999999996 \\
%-1.95454545454545 0 \\
%-1.95454545454545 0.0999999999999996 \\
%-1.95454545454545 0.199999999999999 \\
%-1.95454545454545 0.3 \\
%-1.95454545454545 0.399999999999999 \\
%-1.9 0.454545454545454 \\
%NaN NaN \\
%};
%
%
%\addplot[dashed,line width=2.0pt,draw=black,forget plot]
%table[row sep=crcr]{
%x y\\
%-1.9 0.445454545454545 \\
%-1.84545454545455 0.5 \\
%-1.8 0.545454545454545 \\
%-1.74545454545455 0.6 \\
%-1.74545454545455 0.699999999999999 \\
%-1.7 0.745454545454545 \\
%-1.6 0.745454545454545 \\
%-1.54545454545455 0.8 \\
%-1.5 0.845454545454545 \\
%-1.44545454545455 0.899999999999999 \\
%-1.4 0.945454545454545 \\
%-1.3 0.945454545454545 \\
%-1.2 0.945454545454545 \\
%-1.1 0.945454545454545 \\
%-1 0.945454545454545 \\
%-0.9 0.945454545454545 \\
%-0.8 0.945454545454545 \\
%-0.7 0.945454545454545 \\
%-0.6 0.945454545454545 \\
%-0.554545454545455 0.899999999999999 \\
%-0.5 0.845454545454545 \\
%-0.454545454545455 0.8 \\
%-0.4 0.745454545454545 \\
%-0.3 0.745454545454545 \\
%-0.254545454545455 0.699999999999999 \\
%-0.254545454545455 0.6 \\
%-0.2 0.545454545454545 \\
%-0.154545454545455 0.5 \\
%-0.1 0.445454545454545 \\
%-0.0545454545454547 0.399999999999999 \\
%-0.0545454545454547 0.3 \\
%-0.0545454545454547 0.199999999999999 \\
%-0.0545454545454547 0.0999999999999996 \\
%-0.0545454545454547 0 \\
%-0.0545454545454547 -0.0999999999999996 \\
%-0.0545454545454547 -0.199999999999999 \\
%-0.0545454545454547 -0.3 \\
%-0.0545454545454547 -0.399999999999999 \\
%-0.1 -0.445454545454545 \\
%-0.154545454545455 -0.5 \\
%-0.2 -0.545454545454545 \\
%-0.254545454545455 -0.6 \\
%-0.254545454545455 -0.699999999999999 \\
%-0.3 -0.745454545454545 \\
%-0.354545454545455 -0.8 \\
%-0.4 -0.845454545454545 \\
%-0.5 -0.845454545454545 \\
%-0.554545454545455 -0.899999999999999 \\
%-0.6 -0.945454545454545 \\
%-0.7 -0.945454545454545 \\
%-0.8 -0.945454545454545 \\
%-0.9 -0.945454545454545 \\
%-1 -0.945454545454545 \\
%-1.1 -0.945454545454545 \\
%-1.2 -0.945454545454545 \\
%-1.3 -0.945454545454545 \\
%-1.4 -0.945454545454545 \\
%-1.44545454545455 -0.899999999999999 \\
%-1.5 -0.845454545454545 \\
%-1.6 -0.845454545454545 \\
%-1.64545454545455 -0.8 \\
%-1.7 -0.745454545454545 \\
%-1.74545454545455 -0.699999999999999 \\
%-1.8 -0.645454545454545 \\
%-1.84545454545455 -0.6 \\
%-1.84545454545455 -0.5 \\
%-1.9 -0.445454545454545 \\
%-1.94545454545455 -0.399999999999999 \\
%-1.94545454545455 -0.3 \\
%-1.94545454545455 -0.199999999999999 \\
%-1.94545454545455 -0.0999999999999996 \\
%-1.94545454545455 0 \\
%-1.94545454545455 0.0999999999999996 \\
%-1.94545454545455 0.199999999999999 \\
%-1.94545454545455 0.3 \\
%-1.94545454545455 0.399999999999999 \\
%-1.9 0.445454545454545 \\
%NaN NaN \\
%};
%
%
%\addplot[dashed,line width=2.0pt,draw=black,forget plot]
%table[row sep=crcr]{
%x y\\
%-1.9 0.436363636363636 \\
%-1.83636363636364 0.5 \\
%-1.8 0.536363636363636 \\
%-1.73636363636364 0.6 \\
%-1.73636363636364 0.699999999999999 \\
%-1.7 0.736363636363635 \\
%-1.6 0.736363636363635 \\
%-1.53636363636364 0.8 \\
%-1.5 0.836363636363636 \\
%-1.43636363636364 0.899999999999999 \\
%-1.4 0.936363636363636 \\
%-1.3 0.936363636363636 \\
%-1.2 0.936363636363636 \\
%-1.1 0.936363636363636 \\
%-1 0.936363636363636 \\
%-0.9 0.936363636363636 \\
%-0.8 0.936363636363636 \\
%-0.7 0.936363636363636 \\
%-0.6 0.936363636363636 \\
%-0.563636363636363 0.899999999999999 \\
%-0.5 0.836363636363636 \\
%-0.463636363636363 0.8 \\
%-0.4 0.736363636363635 \\
%-0.3 0.736363636363635 \\
%-0.263636363636364 0.699999999999999 \\
%-0.263636363636364 0.6 \\
%-0.2 0.536363636363636 \\
%-0.163636363636363 0.5 \\
%-0.1 0.436363636363636 \\
%-0.0636363636363634 0.399999999999999 \\
%-0.0636363636363634 0.3 \\
%-0.0636363636363634 0.199999999999999 \\
%-0.0636363636363634 0.0999999999999996 \\
%-0.0636363636363634 0 \\
%-0.0636363636363634 -0.0999999999999996 \\
%-0.0636363636363634 -0.199999999999999 \\
%-0.0636363636363634 -0.3 \\
%-0.0636363636363634 -0.399999999999999 \\
%-0.1 -0.436363636363636 \\
%-0.163636363636363 -0.5 \\
%-0.2 -0.536363636363637 \\
%-0.263636363636364 -0.6 \\
%-0.263636363636364 -0.699999999999999 \\
%-0.3 -0.736363636363636 \\
%-0.363636363636364 -0.8 \\
%-0.4 -0.836363636363636 \\
%-0.5 -0.836363636363636 \\
%-0.563636363636363 -0.899999999999999 \\
%-0.6 -0.936363636363636 \\
%-0.7 -0.936363636363636 \\
%-0.8 -0.936363636363636 \\
%-0.9 -0.936363636363636 \\
%-1 -0.936363636363636 \\
%-1.1 -0.936363636363636 \\
%-1.2 -0.936363636363636 \\
%-1.3 -0.936363636363636 \\
%-1.4 -0.936363636363636 \\
%-1.43636363636364 -0.899999999999999 \\
%-1.5 -0.836363636363636 \\
%-1.6 -0.836363636363636 \\
%-1.63636363636364 -0.8 \\
%-1.7 -0.736363636363636 \\
%-1.73636363636364 -0.699999999999999 \\
%-1.8 -0.636363636363636 \\
%-1.83636363636364 -0.6 \\
%-1.83636363636364 -0.5 \\
%-1.9 -0.436363636363636 \\
%-1.93636363636364 -0.399999999999999 \\
%-1.93636363636364 -0.3 \\
%-1.93636363636364 -0.199999999999999 \\
%-1.93636363636364 -0.0999999999999996 \\
%-1.93636363636364 0 \\
%-1.93636363636364 0.0999999999999996 \\
%-1.93636363636364 0.199999999999999 \\
%-1.93636363636364 0.3 \\
%-1.93636363636364 0.399999999999999 \\
%-1.9 0.436363636363636 \\
%NaN NaN \\
%};
%
%
%\addplot[dashed,line width=2.0pt,draw=black,forget plot]
%table[row sep=crcr]{
%x y\\
%-1.9 0.427272727272726 \\
%-1.82727272727273 0.5 \\
%-1.8 0.527272727272727 \\
%-1.72727272727273 0.6 \\
%-1.72727272727273 0.699999999999999 \\
%-1.7 0.727272727272726 \\
%-1.6 0.727272727272726 \\
%-1.52727272727273 0.8 \\
%-1.5 0.827272727272726 \\
%-1.42727272727273 0.899999999999999 \\
%-1.4 0.927272727272726 \\
%-1.3 0.927272727272726 \\
%-1.2 0.927272727272726 \\
%-1.1 0.927272727272726 \\
%-1 0.927272727272726 \\
%-0.9 0.927272727272726 \\
%-0.8 0.927272727272726 \\
%-0.7 0.927272727272726 \\
%-0.6 0.927272727272726 \\
%-0.572727272727273 0.899999999999999 \\
%-0.5 0.827272727272726 \\
%-0.472727272727273 0.8 \\
%-0.4 0.727272727272726 \\
%-0.3 0.727272727272726 \\
%-0.272727272727273 0.699999999999999 \\
%-0.272727272727273 0.6 \\
%-0.2 0.527272727272727 \\
%-0.172727272727273 0.5 \\
%-0.1 0.427272727272726 \\
%-0.0727272727272727 0.399999999999999 \\
%-0.0727272727272727 0.3 \\
%-0.0727272727272727 0.199999999999999 \\
%-0.0727272727272727 0.0999999999999996 \\
%-0.0727272727272727 0 \\
%-0.0727272727272727 -0.0999999999999996 \\
%-0.0727272727272727 -0.199999999999999 \\
%-0.0727272727272727 -0.3 \\
%-0.0727272727272727 -0.399999999999999 \\
%-0.1 -0.427272727272727 \\
%-0.172727272727273 -0.5 \\
%-0.2 -0.527272727272727 \\
%-0.272727272727273 -0.6 \\
%-0.272727272727273 -0.699999999999999 \\
%-0.3 -0.727272727272727 \\
%-0.372727272727273 -0.8 \\
%-0.4 -0.827272727272727 \\
%-0.5 -0.827272727272727 \\
%-0.572727272727273 -0.899999999999999 \\
%-0.6 -0.927272727272727 \\
%-0.7 -0.927272727272727 \\
%-0.8 -0.927272727272727 \\
%-0.9 -0.927272727272727 \\
%-1 -0.927272727272727 \\
%-1.1 -0.927272727272727 \\
%-1.2 -0.927272727272727 \\
%-1.3 -0.927272727272727 \\
%-1.4 -0.927272727272727 \\
%-1.42727272727273 -0.899999999999999 \\
%-1.5 -0.827272727272727 \\
%-1.6 -0.827272727272727 \\
%-1.62727272727273 -0.8 \\
%-1.7 -0.727272727272727 \\
%-1.72727272727273 -0.699999999999999 \\
%-1.8 -0.627272727272727 \\
%-1.82727272727273 -0.6 \\
%-1.82727272727273 -0.5 \\
%-1.9 -0.427272727272727 \\
%-1.92727272727273 -0.399999999999999 \\
%-1.92727272727273 -0.3 \\
%-1.92727272727273 -0.199999999999999 \\
%-1.92727272727273 -0.0999999999999996 \\
%-1.92727272727273 0 \\
%-1.92727272727273 0.0999999999999996 \\
%-1.92727272727273 0.199999999999999 \\
%-1.92727272727273 0.3 \\
%-1.92727272727273 0.399999999999999 \\
%-1.9 0.427272727272726 \\
%NaN NaN \\
%};
%
%
%\addplot[dashed,line width=2.0pt,draw=black,forget plot]
%table[row sep=crcr]{
%x y\\
%-1.9 0.418181818181818 \\
%-1.81818181818182 0.5 \\
%-1.8 0.518181818181819 \\
%-1.71818181818182 0.6 \\
%-1.71818181818182 0.699999999999999 \\
%-1.7 0.718181818181818 \\
%-1.6 0.718181818181818 \\
%-1.51818181818182 0.8 \\
%-1.5 0.818181818181818 \\
%-1.41818181818182 0.899999999999999 \\
%-1.4 0.918181818181818 \\
%-1.3 0.918181818181818 \\
%-1.2 0.918181818181818 \\
%-1.1 0.918181818181818 \\
%-1 0.918181818181818 \\
%-0.9 0.918181818181818 \\
%-0.8 0.918181818181818 \\
%-0.7 0.918181818181818 \\
%-0.6 0.918181818181818 \\
%-0.581818181818182 0.899999999999999 \\
%-0.5 0.818181818181818 \\
%-0.481818181818182 0.8 \\
%-0.4 0.718181818181818 \\
%-0.3 0.718181818181818 \\
%-0.281818181818182 0.699999999999999 \\
%-0.281818181818182 0.6 \\
%-0.2 0.518181818181819 \\
%-0.181818181818182 0.5 \\
%-0.1 0.418181818181818 \\
%-0.0818181818181821 0.399999999999999 \\
%-0.0818181818181821 0.3 \\
%-0.0818181818181821 0.199999999999999 \\
%-0.0818181818181821 0.0999999999999996 \\
%-0.0818181818181821 0 \\
%-0.0818181818181821 -0.0999999999999996 \\
%-0.0818181818181821 -0.199999999999999 \\
%-0.0818181818181821 -0.3 \\
%-0.0818181818181821 -0.399999999999999 \\
%-0.1 -0.418181818181818 \\
%-0.181818181818182 -0.5 \\
%-0.2 -0.518181818181818 \\
%-0.281818181818182 -0.6 \\
%-0.281818181818182 -0.699999999999999 \\
%-0.3 -0.718181818181817 \\
%-0.381818181818182 -0.8 \\
%-0.4 -0.818181818181818 \\
%-0.5 -0.818181818181818 \\
%-0.581818181818182 -0.899999999999999 \\
%-0.6 -0.918181818181818 \\
%-0.7 -0.918181818181818 \\
%-0.8 -0.918181818181818 \\
%-0.9 -0.918181818181818 \\
%-1 -0.918181818181818 \\
%-1.1 -0.918181818181818 \\
%-1.2 -0.918181818181818 \\
%-1.3 -0.918181818181818 \\
%-1.4 -0.918181818181818 \\
%-1.41818181818182 -0.899999999999999 \\
%-1.5 -0.818181818181818 \\
%-1.6 -0.818181818181818 \\
%-1.61818181818182 -0.8 \\
%-1.7 -0.718181818181817 \\
%-1.71818181818182 -0.699999999999999 \\
%-1.8 -0.618181818181818 \\
%-1.81818181818182 -0.6 \\
%-1.81818181818182 -0.5 \\
%-1.9 -0.418181818181818 \\
%-1.91818181818182 -0.399999999999999 \\
%-1.91818181818182 -0.3 \\
%-1.91818181818182 -0.199999999999999 \\
%-1.91818181818182 -0.0999999999999996 \\
%-1.91818181818182 0 \\
%-1.91818181818182 0.0999999999999996 \\
%-1.91818181818182 0.199999999999999 \\
%-1.91818181818182 0.3 \\
%-1.91818181818182 0.399999999999999 \\
%-1.9 0.418181818181818 \\
%NaN NaN \\
%};
%
%
%\addplot[dashed,line width=2.0pt,draw=black,forget plot]
%table[row sep=crcr]{
%x y\\
%-1.9 0.409090909090909 \\
%-1.80909090909091 0.5 \\
%-1.8 0.509090909090909 \\
%-1.70909090909091 0.6 \\
%-1.70909090909091 0.699999999999999 \\
%-1.7 0.709090909090909 \\
%-1.6 0.709090909090909 \\
%-1.50909090909091 0.8 \\
%-1.5 0.809090909090909 \\
%-1.40909090909091 0.899999999999999 \\
%-1.4 0.909090909090909 \\
%-1.3 0.909090909090909 \\
%-1.2 0.909090909090909 \\
%-1.1 0.909090909090909 \\
%-1 0.909090909090909 \\
%-0.9 0.909090909090909 \\
%-0.8 0.909090909090909 \\
%-0.7 0.909090909090909 \\
%-0.6 0.909090909090909 \\
%-0.590909090909091 0.899999999999999 \\
%-0.5 0.809090909090909 \\
%-0.490909090909091 0.8 \\
%-0.4 0.709090909090909 \\
%-0.3 0.709090909090909 \\
%-0.290909090909091 0.699999999999999 \\
%-0.290909090909091 0.6 \\
%-0.2 0.509090909090909 \\
%-0.190909090909091 0.5 \\
%-0.1 0.409090909090909 \\
%-0.0909090909090907 0.399999999999999 \\
%-0.0909090909090907 0.3 \\
%-0.0909090909090907 0.199999999999999 \\
%-0.0909090909090907 0.0999999999999996 \\
%-0.0909090909090907 0 \\
%-0.0909090909090907 -0.0999999999999996 \\
%-0.0909090909090907 -0.199999999999999 \\
%-0.0909090909090907 -0.3 \\
%-0.0909090909090907 -0.399999999999999 \\
%-0.1 -0.409090909090909 \\
%-0.190909090909091 -0.5 \\
%-0.2 -0.509090909090909 \\
%-0.290909090909091 -0.6 \\
%-0.290909090909091 -0.699999999999999 \\
%-0.3 -0.709090909090909 \\
%-0.390909090909091 -0.8 \\
%-0.4 -0.809090909090909 \\
%-0.5 -0.809090909090909 \\
%-0.590909090909091 -0.899999999999999 \\
%-0.6 -0.909090909090909 \\
%-0.7 -0.909090909090909 \\
%-0.8 -0.909090909090909 \\
%-0.9 -0.909090909090909 \\
%-1 -0.909090909090909 \\
%-1.1 -0.909090909090909 \\
%-1.2 -0.909090909090909 \\
%-1.3 -0.909090909090909 \\
%-1.4 -0.909090909090909 \\
%-1.40909090909091 -0.899999999999999 \\
%-1.5 -0.809090909090909 \\
%-1.6 -0.809090909090909 \\
%-1.60909090909091 -0.8 \\
%-1.7 -0.709090909090909 \\
%-1.70909090909091 -0.699999999999999 \\
%-1.8 -0.609090909090909 \\
%-1.80909090909091 -0.6 \\
%-1.80909090909091 -0.5 \\
%-1.9 -0.409090909090909 \\
%-1.90909090909091 -0.399999999999999 \\
%-1.90909090909091 -0.3 \\
%-1.90909090909091 -0.199999999999999 \\
%-1.90909090909091 -0.0999999999999996 \\
%-1.90909090909091 0 \\
%-1.90909090909091 0.0999999999999996 \\
%-1.90909090909091 0.199999999999999 \\
%-1.90909090909091 0.3 \\
%-1.90909090909091 0.399999999999999 \\
%-1.9 0.409090909090909 \\
%NaN NaN \\
%};
%
%
%\addplot[dashed,line width=2.0pt,draw=black,forget plot]
%table[row sep=crcr]{
%x y\\
%-1.9 0.399999999999999 \\
%-1.8 0.5 \\
%-1.8 0.5 \\
%-1.7 0.6 \\
%-1.7 0.699999999999999 \\
%-1.7 0.699999999999999 \\
%-1.6 0.699999999999999 \\
%-1.5 0.8 \\
%-1.5 0.8 \\
%-1.4 0.899999999999999 \\
%-1.4 0.899999999999999 \\
%-1.3 0.899999999999999 \\
%-1.2 0.899999999999999 \\
%-1.1 0.899999999999999 \\
%-1 0.899999999999999 \\
%-0.9 0.899999999999999 \\
%-0.8 0.899999999999999 \\
%-0.7 0.899999999999999 \\
%-0.6 0.899999999999999 \\
%-0.6 0.899999999999999 \\
%-0.5 0.8 \\
%-0.5 0.8 \\
%-0.4 0.699999999999999 \\
%-0.3 0.699999999999999 \\
%-0.3 0.699999999999999 \\
%-0.3 0.6 \\
%-0.2 0.5 \\
%-0.2 0.5 \\
%-0.1 0.399999999999999 \\
%-0.1 0.399999999999999 \\
%-0.1 0.3 \\
%-0.1 0.199999999999999 \\
%-0.1 0.0999999999999996 \\
%-0.1 0 \\
%-0.1 -0.0999999999999996 \\
%-0.1 -0.199999999999999 \\
%-0.1 -0.3 \\
%-0.1 -0.399999999999999 \\
%-0.1 -0.399999999999999 \\
%-0.2 -0.5 \\
%-0.2 -0.5 \\
%-0.3 -0.6 \\
%-0.3 -0.699999999999999 \\
%-0.3 -0.699999999999999 \\
%-0.4 -0.8 \\
%-0.4 -0.8 \\
%-0.5 -0.8 \\
%-0.6 -0.899999999999999 \\
%-0.6 -0.899999999999999 \\
%-0.7 -0.899999999999999 \\
%-0.8 -0.899999999999999 \\
%-0.9 -0.899999999999999 \\
%-1 -0.899999999999999 \\
%-1.1 -0.899999999999999 \\
%-1.2 -0.899999999999999 \\
%-1.3 -0.899999999999999 \\
%-1.4 -0.899999999999999 \\
%-1.4 -0.899999999999999 \\
%-1.5 -0.8 \\
%-1.6 -0.8 \\
%-1.6 -0.8 \\
%-1.7 -0.699999999999999 \\
%-1.7 -0.699999999999999 \\
%-1.8 -0.6 \\
%-1.8 -0.6 \\
%-1.8 -0.5 \\
%-1.9 -0.399999999999999 \\
%-1.9 -0.399999999999999 \\
%-1.9 -0.3 \\
%-1.9 -0.199999999999999 \\
%-1.9 -0.0999999999999996 \\
%-1.9 0 \\
%-1.9 0.0999999999999996 \\
%-1.9 0.199999999999999 \\
%-1.9 0.3 \\
%-1.9 0.399999999999999 \\
%-1.9 0.399999999999999 \\
%NaN NaN \\
%};


\addplot[dashed,line width=2.0pt,draw=red,forget plot]
table[row sep=crcr]{
x y\\
2 1.00000001168609e-008 \\
1.90000001 0.0999999999999996 \\
1.90000001 0.199999999999999 \\
1.90000001 0.3 \\
1.90000001 0.399999999999999 \\
1.9 0.400000009999999 \\
1.80000001 0.5 \\
1.8 0.500000009999999 \\
1.70000001 0.6 \\
1.70000001 0.699999999999999 \\
1.7 0.700000009999999 \\
1.6 0.700000009999999 \\
1.50000001 0.8 \\
1.5 0.800000009999999 \\
1.40000001 0.899999999999999 \\
1.4 0.900000009999999 \\
1.3 0.900000009999999 \\
1.2 0.900000009999999 \\
1.1 0.900000009999999 \\
1.00000001 1 \\
1 1.00000001 \\
0.99999999 1 \\
0.9 0.900000009999999 \\
0.8 0.900000009999999 \\
0.7 0.900000009999999 \\
0.6 0.900000009999999 \\
0.59999999 0.899999999999999 \\
0.5 0.800000009999999 \\
0.49999999 0.8 \\
0.4 0.700000009999999 \\
0.3 0.700000009999999 \\
0.29999999 0.699999999999999 \\
0.29999999 0.6 \\
0.2 0.500000009999999 \\
0.19999999 0.5 \\
0.0999999999999999 0.400000009999999 \\
0.0999999899999997 0.399999999999999 \\
0.0999999899999997 0.3 \\
0.0999999899999997 0.199999999999999 \\
0.0999999899999997 0.0999999999999996 \\
0 1.00000001168609e-008 \\
-1.0000000116861e-008 0 \\
0 -1.00000001168609e-008 \\
0.0999999899999997 -0.0999999999999996 \\
0.0999999899999997 -0.199999999999999 \\
0.0999999899999997 -0.3 \\
0.0999999899999997 -0.399999999999999 \\
0.0999999999999999 -0.40000001 \\
0.19999999 -0.5 \\
0.19999999 -0.6 \\
0.2 -0.60000001 \\
0.29999999 -0.699999999999999 \\
0.3 -0.700000009999999 \\
0.39999999 -0.8 \\
0.4 -0.80000001 \\
0.5 -0.80000001 \\
0.59999999 -0.899999999999999 \\
0.6 -0.90000001 \\
0.7 -0.90000001 \\
0.8 -0.90000001 \\
0.9 -0.90000001 \\
0.99999999 -1 \\
1 -1.00000001 \\
1.00000001 -1 \\
1.1 -0.90000001 \\
1.2 -0.90000001 \\
1.3 -0.90000001 \\
1.4 -0.90000001 \\
1.40000001 -0.899999999999999 \\
1.5 -0.80000001 \\
1.50000001 -0.8 \\
1.6 -0.700000009999999 \\
1.7 -0.700000009999999 \\
1.70000001 -0.699999999999999 \\
1.70000001 -0.6 \\
1.8 -0.50000001 \\
1.80000001 -0.5 \\
1.9 -0.40000001 \\
1.90000001 -0.399999999999999 \\
1.90000001 -0.3 \\
1.90000001 -0.199999999999999 \\
1.90000001 -0.0999999999999996 \\
2 -1.00000001168609e-008 \\
NaN NaN \\
};


\addplot[dashed,line width=2.0pt,draw=red,forget plot]
table[row sep=crcr]{
x y\\
2 0.00909090909090932 \\
1.90909090909091 0.0999999999999996 \\
1.90909090909091 0.199999999999999 \\
1.90909090909091 0.3 \\
1.90909090909091 0.399999999999999 \\
1.9 0.409090909090909 \\
1.80909090909091 0.5 \\
1.8 0.509090909090909 \\
1.70909090909091 0.6 \\
1.70909090909091 0.699999999999999 \\
1.7 0.709090909090909 \\
1.6 0.709090909090909 \\
1.50909090909091 0.8 \\
1.5 0.809090909090909 \\
1.40909090909091 0.899999999999999 \\
1.4 0.909090909090909 \\
1.3 0.909090909090909 \\
1.2 0.909090909090909 \\
1.1 0.909090909090909 \\
1.00909090909091 1 \\
1 1.00909090909091 \\
0.990909090909091 1 \\
0.9 0.909090909090909 \\
0.8 0.909090909090909 \\
0.7 0.909090909090909 \\
0.6 0.909090909090909 \\
0.590909090909091 0.899999999999999 \\
0.5 0.809090909090909 \\
0.490909090909091 0.8 \\
0.4 0.709090909090909 \\
0.3 0.709090909090909 \\
0.29090909090909 0.699999999999999 \\
0.29090909090909 0.6 \\
0.2 0.509090909090909 \\
0.190909090909091 0.5 \\
0.0999999999999999 0.409090909090909 \\
0.0909090909090905 0.399999999999999 \\
0.0909090909090905 0.3 \\
0.0909090909090905 0.199999999999999 \\
0.0909090909090905 0.0999999999999996 \\
0 0.00909090909090932 \\
-0.00909090909090936 0 \\
0 -0.00909090909090932 \\
0.0909090909090905 -0.0999999999999996 \\
0.0909090909090905 -0.199999999999999 \\
0.0909090909090905 -0.3 \\
0.0909090909090905 -0.399999999999999 \\
0.0999999999999999 -0.409090909090909 \\
0.190909090909091 -0.5 \\
0.190909090909091 -0.6 \\
0.2 -0.609090909090909 \\
0.29090909090909 -0.699999999999999 \\
0.3 -0.709090909090909 \\
0.390909090909091 -0.8 \\
0.4 -0.809090909090909 \\
0.5 -0.809090909090909 \\
0.590909090909091 -0.899999999999999 \\
0.6 -0.909090909090909 \\
0.7 -0.909090909090909 \\
0.8 -0.909090909090909 \\
0.9 -0.909090909090909 \\
0.990909090909091 -1 \\
1 -1.00909090909091 \\
1.00909090909091 -1 \\
1.1 -0.909090909090909 \\
1.2 -0.909090909090909 \\
1.3 -0.909090909090909 \\
1.4 -0.909090909090909 \\
1.40909090909091 -0.899999999999999 \\
1.5 -0.809090909090909 \\
1.50909090909091 -0.8 \\
1.6 -0.709090909090909 \\
1.7 -0.709090909090909 \\
1.70909090909091 -0.699999999999999 \\
1.70909090909091 -0.6 \\
1.8 -0.509090909090909 \\
1.80909090909091 -0.5 \\
1.9 -0.409090909090909 \\
1.90909090909091 -0.399999999999999 \\
1.90909090909091 -0.3 \\
1.90909090909091 -0.199999999999999 \\
1.90909090909091 -0.0999999999999996 \\
2 -0.00909090909090932 \\
NaN NaN \\
};


\addplot[dashed,line width=2.0pt,draw=red,forget plot]
table[row sep=crcr]{
x y\\
2 0.0181818181818179 \\
1.91818181818182 0.0999999999999996 \\
1.91818181818182 0.199999999999999 \\
1.91818181818182 0.3 \\
1.91818181818182 0.399999999999999 \\
1.9 0.418181818181818 \\
1.81818181818182 0.5 \\
1.8 0.518181818181819 \\
1.71818181818182 0.6 \\
1.71818181818182 0.699999999999999 \\
1.7 0.718181818181818 \\
1.6 0.718181818181818 \\
1.51818181818182 0.8 \\
1.5 0.818181818181818 \\
1.41818181818182 0.899999999999999 \\
1.4 0.918181818181818 \\
1.3 0.918181818181818 \\
1.2 0.918181818181818 \\
1.1 0.918181818181818 \\
1.01818181818182 1 \\
1 1.01818181818182 \\
0.981818181818182 1 \\
0.9 0.918181818181818 \\
0.8 0.918181818181818 \\
0.7 0.918181818181818 \\
0.6 0.918181818181818 \\
0.581818181818182 0.899999999999999 \\
0.5 0.818181818181818 \\
0.481818181818182 0.8 \\
0.4 0.718181818181818 \\
0.3 0.718181818181818 \\
0.281818181818182 0.699999999999999 \\
0.281818181818182 0.6 \\
0.2 0.518181818181819 \\
0.181818181818182 0.5 \\
0.0999999999999999 0.418181818181818 \\
0.0818181818181819 0.399999999999999 \\
0.0818181818181819 0.3 \\
0.0818181818181819 0.199999999999999 \\
0.0818181818181819 0.0999999999999996 \\
0 0.0181818181818179 \\
-0.018181818181818 0 \\
0 -0.0181818181818179 \\
0.0818181818181819 -0.0999999999999996 \\
0.0818181818181819 -0.199999999999999 \\
0.0818181818181819 -0.3 \\
0.0818181818181819 -0.399999999999999 \\
0.0999999999999999 -0.418181818181818 \\
0.181818181818182 -0.5 \\
0.181818181818182 -0.6 \\
0.2 -0.618181818181818 \\
0.281818181818182 -0.699999999999999 \\
0.3 -0.718181818181817 \\
0.381818181818182 -0.8 \\
0.4 -0.818181818181818 \\
0.5 -0.818181818181818 \\
0.581818181818182 -0.899999999999999 \\
0.6 -0.918181818181818 \\
0.7 -0.918181818181818 \\
0.8 -0.918181818181818 \\
0.9 -0.918181818181818 \\
0.981818181818182 -1 \\
1 -1.01818181818182 \\
1.01818181818182 -1 \\
1.1 -0.918181818181818 \\
1.2 -0.918181818181818 \\
1.3 -0.918181818181818 \\
1.4 -0.918181818181818 \\
1.41818181818182 -0.899999999999999 \\
1.5 -0.818181818181818 \\
1.51818181818182 -0.8 \\
1.6 -0.718181818181817 \\
1.7 -0.718181818181817 \\
1.71818181818182 -0.699999999999999 \\
1.71818181818182 -0.6 \\
1.8 -0.518181818181818 \\
1.81818181818182 -0.5 \\
1.9 -0.418181818181818 \\
1.91818181818182 -0.399999999999999 \\
1.91818181818182 -0.3 \\
1.91818181818182 -0.199999999999999 \\
1.91818181818182 -0.0999999999999996 \\
2 -0.0181818181818179 \\
NaN NaN \\
};


\addplot[dashed,line width=2.0pt,draw=red,forget plot]
table[row sep=crcr]{
x y\\
2 0.0272727272727272 \\
1.92727272727273 0.0999999999999996 \\
1.92727272727273 0.199999999999999 \\
1.92727272727273 0.3 \\
1.92727272727273 0.399999999999999 \\
1.9 0.427272727272726 \\
1.82727272727273 0.5 \\
1.8 0.527272727272727 \\
1.72727272727273 0.6 \\
1.72727272727273 0.699999999999999 \\
1.7 0.727272727272726 \\
1.6 0.727272727272726 \\
1.52727272727273 0.8 \\
1.5 0.827272727272726 \\
1.42727272727273 0.899999999999999 \\
1.4 0.927272727272726 \\
1.3 0.927272727272726 \\
1.2 0.927272727272726 \\
1.1 0.927272727272726 \\
1.02727272727273 1 \\
1 1.02727272727273 \\
0.972727272727273 1 \\
0.9 0.927272727272726 \\
0.8 0.927272727272726 \\
0.7 0.927272727272726 \\
0.6 0.927272727272726 \\
0.572727272727273 0.899999999999999 \\
0.5 0.827272727272726 \\
0.472727272727273 0.8 \\
0.4 0.727272727272726 \\
0.3 0.727272727272726 \\
0.272727272727273 0.699999999999999 \\
0.272727272727273 0.6 \\
0.2 0.527272727272727 \\
0.172727272727273 0.5 \\
0.0999999999999999 0.427272727272726 \\
0.0727272727272726 0.399999999999999 \\
0.0727272727272726 0.3 \\
0.0727272727272726 0.199999999999999 \\
0.0727272727272726 0.0999999999999996 \\
0 0.0272727272727272 \\
-0.0272727272727274 0 \\
0 -0.0272727272727272 \\
0.0727272727272726 -0.0999999999999996 \\
0.0727272727272726 -0.199999999999999 \\
0.0727272727272726 -0.3 \\
0.0727272727272726 -0.399999999999999 \\
0.0999999999999999 -0.427272727272727 \\
0.172727272727273 -0.5 \\
0.172727272727273 -0.6 \\
0.2 -0.627272727272727 \\
0.272727272727273 -0.699999999999999 \\
0.3 -0.727272727272727 \\
0.372727272727273 -0.8 \\
0.4 -0.827272727272727 \\
0.5 -0.827272727272727 \\
0.572727272727273 -0.899999999999999 \\
0.6 -0.927272727272727 \\
0.7 -0.927272727272727 \\
0.8 -0.927272727272727 \\
0.9 -0.927272727272727 \\
0.972727272727273 -1 \\
1 -1.02727272727273 \\
1.02727272727273 -1 \\
1.1 -0.927272727272727 \\
1.2 -0.927272727272727 \\
1.3 -0.927272727272727 \\
1.4 -0.927272727272727 \\
1.42727272727273 -0.899999999999999 \\
1.5 -0.827272727272727 \\
1.52727272727273 -0.8 \\
1.6 -0.727272727272727 \\
1.7 -0.727272727272727 \\
1.72727272727273 -0.699999999999999 \\
1.72727272727273 -0.6 \\
1.8 -0.527272727272727 \\
1.82727272727273 -0.5 \\
1.9 -0.427272727272727 \\
1.92727272727273 -0.399999999999999 \\
1.92727272727273 -0.3 \\
1.92727272727273 -0.199999999999999 \\
1.92727272727273 -0.0999999999999996 \\
2 -0.0272727272727272 \\
NaN NaN \\
};


\addplot[dashed,line width=2.0pt,draw=red,forget plot]
table[row sep=crcr]{
x y\\
2 0.0363636363636366 \\
1.93636363636364 0.0999999999999996 \\
1.93636363636364 0.199999999999999 \\
1.93636363636364 0.3 \\
1.93636363636364 0.399999999999999 \\
1.9 0.436363636363636 \\
1.83636363636364 0.5 \\
1.8 0.536363636363636 \\
1.73636363636364 0.6 \\
1.73636363636364 0.699999999999999 \\
1.7 0.736363636363635 \\
1.6 0.736363636363635 \\
1.53636363636364 0.8 \\
1.5 0.836363636363636 \\
1.43636363636364 0.899999999999999 \\
1.4 0.936363636363636 \\
1.3 0.936363636363636 \\
1.2 0.936363636363636 \\
1.1 0.936363636363636 \\
1.03636363636364 1 \\
1 1.03636363636364 \\
0.963636363636363 1 \\
0.9 0.936363636363636 \\
0.8 0.936363636363636 \\
0.7 0.936363636363636 \\
0.6 0.936363636363636 \\
0.563636363636363 0.899999999999999 \\
0.5 0.836363636363636 \\
0.463636363636363 0.8 \\
0.4 0.736363636363635 \\
0.3 0.736363636363635 \\
0.263636363636363 0.699999999999999 \\
0.263636363636363 0.6 \\
0.2 0.536363636363636 \\
0.163636363636363 0.5 \\
0.0999999999999999 0.436363636363636 \\
0.0636363636363632 0.399999999999999 \\
0.0636363636363632 0.3 \\
0.0636363636363632 0.199999999999999 \\
0.0636363636363632 0.0999999999999996 \\
0 0.0363636363636366 \\
-0.0363636363636367 0 \\
0 -0.0363636363636366 \\
0.0636363636363632 -0.0999999999999996 \\
0.0636363636363632 -0.199999999999999 \\
0.0636363636363632 -0.3 \\
0.0636363636363632 -0.399999999999999 \\
0.0999999999999999 -0.436363636363636 \\
0.163636363636363 -0.5 \\
0.163636363636363 -0.6 \\
0.2 -0.636363636363636 \\
0.263636363636363 -0.699999999999999 \\
0.3 -0.736363636363636 \\
0.363636363636363 -0.8 \\
0.4 -0.836363636363636 \\
0.5 -0.836363636363636 \\
0.563636363636363 -0.899999999999999 \\
0.6 -0.936363636363636 \\
0.7 -0.936363636363636 \\
0.8 -0.936363636363636 \\
0.9 -0.936363636363636 \\
0.963636363636363 -1 \\
1 -1.03636363636364 \\
1.03636363636364 -1 \\
1.1 -0.936363636363636 \\
1.2 -0.936363636363636 \\
1.3 -0.936363636363636 \\
1.4 -0.936363636363636 \\
1.43636363636364 -0.899999999999999 \\
1.5 -0.836363636363636 \\
1.53636363636364 -0.8 \\
1.6 -0.736363636363636 \\
1.7 -0.736363636363636 \\
1.73636363636364 -0.699999999999999 \\
1.73636363636364 -0.6 \\
1.8 -0.536363636363637 \\
1.83636363636364 -0.5 \\
1.9 -0.436363636363636 \\
1.93636363636364 -0.399999999999999 \\
1.93636363636364 -0.3 \\
1.93636363636364 -0.199999999999999 \\
1.93636363636364 -0.0999999999999996 \\
2 -0.0363636363636366 \\
NaN NaN \\
};


\addplot[dashed,line width=2.0pt,draw=red,forget plot]
table[row sep=crcr]{
x y\\
2 0.0454545454545452 \\
1.94545454545455 0.0999999999999996 \\
1.94545454545455 0.199999999999999 \\
1.94545454545455 0.3 \\
1.94545454545455 0.399999999999999 \\
1.9 0.445454545454545 \\
1.84545454545455 0.5 \\
1.8 0.545454545454545 \\
1.74545454545455 0.6 \\
1.74545454545455 0.699999999999999 \\
1.7 0.745454545454545 \\
1.6 0.745454545454545 \\
1.54545454545455 0.8 \\
1.5 0.845454545454545 \\
1.44545454545455 0.899999999999999 \\
1.4 0.945454545454545 \\
1.3 0.945454545454545 \\
1.2 0.945454545454545 \\
1.1 0.945454545454545 \\
1.04545454545455 1 \\
1 1.04545454545455 \\
0.954545454545455 1 \\
0.9 0.945454545454545 \\
0.8 0.945454545454545 \\
0.7 0.945454545454545 \\
0.6 0.945454545454545 \\
0.554545454545455 0.899999999999999 \\
0.5 0.845454545454545 \\
0.454545454545455 0.8 \\
0.4 0.745454545454545 \\
0.3 0.745454545454545 \\
0.254545454545455 0.699999999999999 \\
0.254545454545455 0.6 \\
0.2 0.545454545454545 \\
0.154545454545455 0.5 \\
0.0999999999999999 0.445454545454545 \\
0.0545454545454546 0.399999999999999 \\
0.0545454545454546 0.3 \\
0.0545454545454546 0.199999999999999 \\
0.0545454545454546 0.0999999999999996 \\
0 0.0454545454545452 \\
-0.0454545454545454 0 \\
0 -0.0454545454545452 \\
0.0545454545454546 -0.0999999999999996 \\
0.0545454545454546 -0.199999999999999 \\
0.0545454545454546 -0.3 \\
0.0545454545454546 -0.399999999999999 \\
0.0999999999999999 -0.445454545454545 \\
0.154545454545455 -0.5 \\
0.154545454545455 -0.6 \\
0.2 -0.645454545454545 \\
0.254545454545455 -0.699999999999999 \\
0.3 -0.745454545454545 \\
0.354545454545455 -0.8 \\
0.4 -0.845454545454545 \\
0.5 -0.845454545454545 \\
0.554545454545455 -0.899999999999999 \\
0.6 -0.945454545454545 \\
0.7 -0.945454545454545 \\
0.8 -0.945454545454545 \\
0.9 -0.945454545454545 \\
0.954545454545455 -1 \\
1 -1.04545454545455 \\
1.04545454545455 -1 \\
1.1 -0.945454545454545 \\
1.2 -0.945454545454545 \\
1.3 -0.945454545454545 \\
1.4 -0.945454545454545 \\
1.44545454545455 -0.899999999999999 \\
1.5 -0.845454545454545 \\
1.54545454545455 -0.8 \\
1.6 -0.745454545454545 \\
1.7 -0.745454545454545 \\
1.74545454545455 -0.699999999999999 \\
1.74545454545455 -0.6 \\
1.8 -0.545454545454545 \\
1.84545454545455 -0.5 \\
1.9 -0.445454545454545 \\
1.94545454545455 -0.399999999999999 \\
1.94545454545455 -0.3 \\
1.94545454545455 -0.199999999999999 \\
1.94545454545455 -0.0999999999999996 \\
2 -0.0454545454545452 \\
NaN NaN \\
};


\addplot[dashed,line width=2.0pt,draw=red,forget plot]
table[row sep=crcr]{
x y\\
2 0.0545454545454545 \\
1.95454545454545 0.0999999999999996 \\
1.95454545454545 0.199999999999999 \\
1.95454545454545 0.3 \\
1.95454545454545 0.399999999999999 \\
1.9 0.454545454545454 \\
1.85454545454545 0.5 \\
1.8 0.554545454545454 \\
1.75454545454545 0.6 \\
1.75454545454545 0.699999999999999 \\
1.7 0.754545454545454 \\
1.6 0.754545454545454 \\
1.55454545454545 0.8 \\
1.5 0.854545454545454 \\
1.45454545454545 0.899999999999999 \\
1.4 0.954545454545454 \\
1.3 0.954545454545454 \\
1.2 0.954545454545454 \\
1.1 0.954545454545454 \\
1.05454545454545 1 \\
1 1.05454545454545 \\
0.945454545454545 1 \\
0.9 0.954545454545454 \\
0.8 0.954545454545454 \\
0.7 0.954545454545454 \\
0.6 0.954545454545454 \\
0.545454545454545 0.899999999999999 \\
0.5 0.854545454545454 \\
0.445454545454545 0.8 \\
0.4 0.754545454545454 \\
0.3 0.754545454545454 \\
0.245454545454545 0.699999999999999 \\
0.245454545454545 0.6 \\
0.2 0.554545454545454 \\
0.145454545454545 0.5 \\
0.0999999999999999 0.454545454545454 \\
0.0454545454545453 0.399999999999999 \\
0.0454545454545453 0.3 \\
0.0454545454545453 0.199999999999999 \\
0.0454545454545453 0.0999999999999996 \\
0 0.0545454545454545 \\
-0.0545454545454547 0 \\
0 -0.0545454545454545 \\
0.0454545454545453 -0.0999999999999996 \\
0.0454545454545453 -0.199999999999999 \\
0.0454545454545453 -0.3 \\
0.0454545454545453 -0.399999999999999 \\
0.0999999999999999 -0.454545454545454 \\
0.145454545454545 -0.5 \\
0.145454545454545 -0.6 \\
0.2 -0.654545454545454 \\
0.245454545454545 -0.699999999999999 \\
0.3 -0.754545454545454 \\
0.345454545454545 -0.8 \\
0.4 -0.854545454545454 \\
0.5 -0.854545454545454 \\
0.545454545454545 -0.899999999999999 \\
0.6 -0.954545454545454 \\
0.7 -0.954545454545454 \\
0.8 -0.954545454545454 \\
0.9 -0.954545454545454 \\
0.945454545454545 -1 \\
1 -1.05454545454545 \\
1.05454545454545 -1 \\
1.1 -0.954545454545454 \\
1.2 -0.954545454545454 \\
1.3 -0.954545454545454 \\
1.4 -0.954545454545454 \\
1.45454545454545 -0.899999999999999 \\
1.5 -0.854545454545454 \\
1.55454545454545 -0.8 \\
1.6 -0.754545454545454 \\
1.7 -0.754545454545454 \\
1.75454545454545 -0.699999999999999 \\
1.75454545454545 -0.6 \\
1.8 -0.554545454545454 \\
1.85454545454545 -0.5 \\
1.9 -0.454545454545454 \\
1.95454545454545 -0.399999999999999 \\
1.95454545454545 -0.3 \\
1.95454545454545 -0.199999999999999 \\
1.95454545454545 -0.0999999999999996 \\
2 -0.0545454545454545 \\
NaN NaN \\
};


\addplot[dashed,line width=2.0pt,draw=red,forget plot]
table[row sep=crcr]{
x y\\
2 0.0636363636363631 \\
1.96363636363636 0.0999999999999996 \\
1.96363636363636 0.199999999999999 \\
1.96363636363636 0.3 \\
1.96363636363636 0.399999999999999 \\
1.9 0.463636363636364 \\
1.86363636363636 0.5 \\
1.8 0.563636363636364 \\
1.76363636363636 0.6 \\
1.76363636363636 0.699999999999999 \\
1.7 0.763636363636364 \\
1.6 0.763636363636364 \\
1.56363636363636 0.8 \\
1.5 0.863636363636364 \\
1.46363636363636 0.899999999999999 \\
1.4 0.963636363636364 \\
1.3 0.963636363636364 \\
1.2 0.963636363636364 \\
1.1 0.963636363636364 \\
1.06363636363636 1 \\
1 1.06363636363636 \\
0.936363636363637 1 \\
0.9 0.963636363636364 \\
0.8 0.963636363636364 \\
0.7 0.963636363636364 \\
0.6 0.963636363636364 \\
0.536363636363637 0.899999999999999 \\
0.5 0.863636363636364 \\
0.436363636363637 0.8 \\
0.4 0.763636363636364 \\
0.3 0.763636363636364 \\
0.236363636363637 0.699999999999999 \\
0.236363636363637 0.6 \\
0.2 0.563636363636364 \\
0.136363636363637 0.5 \\
0.0999999999999999 0.463636363636364 \\
0.0363636363636366 0.399999999999999 \\
0.0363636363636366 0.3 \\
0.0363636363636366 0.199999999999999 \\
0.0363636363636366 0.0999999999999996 \\
0 0.0636363636363631 \\
-0.0636363636363634 0 \\
0 -0.0636363636363631 \\
0.0363636363636366 -0.0999999999999996 \\
0.0363636363636366 -0.199999999999999 \\
0.0363636363636366 -0.3 \\
0.0363636363636366 -0.399999999999999 \\
0.0999999999999999 -0.463636363636363 \\
0.136363636363637 -0.5 \\
0.136363636363637 -0.6 \\
0.2 -0.663636363636363 \\
0.236363636363637 -0.699999999999999 \\
0.3 -0.763636363636363 \\
0.336363636363637 -0.8 \\
0.4 -0.863636363636363 \\
0.5 -0.863636363636363 \\
0.536363636363637 -0.899999999999999 \\
0.6 -0.963636363636363 \\
0.7 -0.963636363636363 \\
0.8 -0.963636363636363 \\
0.9 -0.963636363636363 \\
0.936363636363637 -1 \\
1 -1.06363636363636 \\
1.06363636363636 -1 \\
1.1 -0.963636363636363 \\
1.2 -0.963636363636363 \\
1.3 -0.963636363636363 \\
1.4 -0.963636363636363 \\
1.46363636363636 -0.899999999999999 \\
1.5 -0.863636363636363 \\
1.56363636363636 -0.8 \\
1.6 -0.763636363636363 \\
1.7 -0.763636363636363 \\
1.76363636363636 -0.699999999999999 \\
1.76363636363636 -0.6 \\
1.8 -0.563636363636363 \\
1.86363636363636 -0.5 \\
1.9 -0.463636363636363 \\
1.96363636363636 -0.399999999999999 \\
1.96363636363636 -0.3 \\
1.96363636363636 -0.199999999999999 \\
1.96363636363636 -0.0999999999999996 \\
2 -0.0636363636363631 \\
NaN NaN \\
};


\addplot[dashed,line width=2.0pt,draw=red,forget plot]
table[row sep=crcr]{
x y\\
2 0.0727272727272724 \\
1.97272727272727 0.0999999999999996 \\
1.97272727272727 0.199999999999999 \\
1.97272727272727 0.3 \\
1.97272727272727 0.399999999999999 \\
1.9 0.472727272727273 \\
1.87272727272727 0.5 \\
1.8 0.572727272727273 \\
1.77272727272727 0.6 \\
1.77272727272727 0.699999999999999 \\
1.7 0.772727272727273 \\
1.6 0.772727272727273 \\
1.57272727272727 0.8 \\
1.5 0.872727272727273 \\
1.47272727272727 0.899999999999999 \\
1.4 0.972727272727273 \\
1.3 0.972727272727273 \\
1.2 0.972727272727273 \\
1.1 0.972727272727273 \\
1.07272727272727 1 \\
1 1.07272727272727 \\
0.927272727272727 1 \\
0.9 0.972727272727273 \\
0.8 0.972727272727273 \\
0.7 0.972727272727273 \\
0.6 0.972727272727273 \\
0.527272727272727 0.899999999999999 \\
0.5 0.872727272727273 \\
0.427272727272727 0.8 \\
0.4 0.772727272727273 \\
0.3 0.772727272727273 \\
0.227272727272727 0.699999999999999 \\
0.227272727272727 0.6 \\
0.2 0.572727272727273 \\
0.127272727272727 0.5 \\
0.0999999999999999 0.472727272727273 \\
0.0272727272727273 0.399999999999999 \\
0.0272727272727273 0.3 \\
0.0272727272727273 0.199999999999999 \\
0.0272727272727273 0.0999999999999996 \\
0 0.0727272727272724 \\
-0.0727272727272727 0 \\
0 -0.0727272727272724 \\
0.0272727272727273 -0.0999999999999996 \\
0.0272727272727273 -0.199999999999999 \\
0.0272727272727273 -0.3 \\
0.0272727272727273 -0.399999999999999 \\
0.0999999999999999 -0.472727272727273 \\
0.127272727272727 -0.5 \\
0.127272727272727 -0.6 \\
0.2 -0.672727272727272 \\
0.227272727272727 -0.699999999999999 \\
0.3 -0.772727272727272 \\
0.327272727272727 -0.8 \\
0.4 -0.872727272727272 \\
0.5 -0.872727272727272 \\
0.527272727272727 -0.899999999999999 \\
0.6 -0.972727272727273 \\
0.7 -0.972727272727273 \\
0.8 -0.972727272727273 \\
0.9 -0.972727272727273 \\
0.927272727272727 -1 \\
1 -1.07272727272727 \\
1.07272727272727 -1 \\
1.1 -0.972727272727273 \\
1.2 -0.972727272727273 \\
1.3 -0.972727272727273 \\
1.4 -0.972727272727273 \\
1.47272727272727 -0.899999999999999 \\
1.5 -0.872727272727272 \\
1.57272727272727 -0.8 \\
1.6 -0.772727272727272 \\
1.7 -0.772727272727272 \\
1.77272727272727 -0.699999999999999 \\
1.77272727272727 -0.6 \\
1.8 -0.572727272727272 \\
1.87272727272727 -0.5 \\
1.9 -0.472727272727273 \\
1.97272727272727 -0.399999999999999 \\
1.97272727272727 -0.3 \\
1.97272727272727 -0.199999999999999 \\
1.97272727272727 -0.0999999999999996 \\
2 -0.0727272727272724 \\
NaN NaN \\
};


\addplot[dashed,line width=2.0pt,draw=red,forget plot]
table[row sep=crcr]{
x y\\
2 0.0818181818181817 \\
1.98181818181818 0.0999999999999996 \\
1.98181818181818 0.199999999999999 \\
1.98181818181818 0.3 \\
1.98181818181818 0.399999999999999 \\
1.9 0.481818181818181 \\
1.88181818181818 0.5 \\
1.8 0.581818181818181 \\
1.78181818181818 0.6 \\
1.78181818181818 0.699999999999999 \\
1.7 0.781818181818181 \\
1.6 0.781818181818181 \\
1.58181818181818 0.8 \\
1.5 0.881818181818181 \\
1.48181818181818 0.899999999999999 \\
1.4 0.981818181818181 \\
1.3 0.981818181818181 \\
1.2 0.981818181818181 \\
1.1 0.981818181818181 \\
1.08181818181818 1 \\
1 1.08181818181818 \\
0.918181818181818 1 \\
0.9 0.981818181818181 \\
0.8 0.981818181818181 \\
0.7 0.981818181818181 \\
0.6 0.981818181818181 \\
0.518181818181818 0.899999999999999 \\
0.5 0.881818181818181 \\
0.418181818181818 0.8 \\
0.4 0.781818181818181 \\
0.3 0.781818181818181 \\
0.218181818181818 0.699999999999999 \\
0.218181818181818 0.6 \\
0.2 0.581818181818181 \\
0.118181818181818 0.5 \\
0.0999999999999999 0.481818181818181 \\
0.018181818181818 0.399999999999999 \\
0.018181818181818 0.3 \\
0.018181818181818 0.199999999999999 \\
0.018181818181818 0.0999999999999996 \\
0 0.0818181818181817 \\
-0.0818181818181821 0 \\
0 -0.0818181818181817 \\
0.018181818181818 -0.0999999999999996 \\
0.018181818181818 -0.199999999999999 \\
0.018181818181818 -0.3 \\
0.018181818181818 -0.399999999999999 \\
0.0999999999999999 -0.481818181818182 \\
0.118181818181818 -0.5 \\
0.118181818181818 -0.6 \\
0.2 -0.681818181818181 \\
0.218181818181818 -0.699999999999999 \\
0.3 -0.781818181818182 \\
0.318181818181818 -0.8 \\
0.4 -0.881818181818182 \\
0.5 -0.881818181818182 \\
0.518181818181818 -0.899999999999999 \\
0.6 -0.981818181818182 \\
0.7 -0.981818181818182 \\
0.8 -0.981818181818182 \\
0.9 -0.981818181818182 \\
0.918181818181818 -1 \\
1 -1.08181818181818 \\
1.08181818181818 -1 \\
1.1 -0.981818181818182 \\
1.2 -0.981818181818182 \\
1.3 -0.981818181818182 \\
1.4 -0.981818181818182 \\
1.48181818181818 -0.899999999999999 \\
1.5 -0.881818181818182 \\
1.58181818181818 -0.8 \\
1.6 -0.781818181818182 \\
1.7 -0.781818181818182 \\
1.78181818181818 -0.699999999999999 \\
1.78181818181818 -0.6 \\
1.8 -0.581818181818182 \\
1.88181818181818 -0.5 \\
1.9 -0.481818181818182 \\
1.98181818181818 -0.399999999999999 \\
1.98181818181818 -0.3 \\
1.98181818181818 -0.199999999999999 \\
1.98181818181818 -0.0999999999999996 \\
2 -0.0818181818181817 \\
NaN NaN \\
};


\addplot[dashed,line width=2.0pt,draw=red,forget plot]
table[row sep=crcr]{
x y\\
2 0.0909090909090903 \\
1.99090909090909 0.0999999999999996 \\
1.99090909090909 0.199999999999999 \\
1.99090909090909 0.3 \\
1.99090909090909 0.399999999999999 \\
1.9 0.490909090909091 \\
1.89090909090909 0.5 \\
1.8 0.59090909090909 \\
1.79090909090909 0.6 \\
1.79090909090909 0.699999999999999 \\
1.7 0.79090909090909 \\
1.6 0.79090909090909 \\
1.59090909090909 0.8 \\
1.5 0.89090909090909 \\
1.49090909090909 0.899999999999999 \\
1.4 0.990909090909091 \\
1.3 0.990909090909091 \\
1.2 0.990909090909091 \\
1.1 0.990909090909091 \\
1.09090909090909 1 \\
1 1.09090909090909 \\
0.909090909090909 1 \\
0.9 0.990909090909091 \\
0.8 0.990909090909091 \\
0.7 0.990909090909091 \\
0.6 0.990909090909091 \\
0.509090909090909 0.899999999999999 \\
0.5 0.89090909090909 \\
0.409090909090909 0.8 \\
0.4 0.79090909090909 \\
0.3 0.79090909090909 \\
0.209090909090909 0.699999999999999 \\
0.209090909090909 0.6 \\
0.2 0.59090909090909 \\
0.109090909090909 0.5 \\
0.0999999999999999 0.490909090909091 \\
0.00909090909090934 0.399999999999999 \\
0.00909090909090934 0.3 \\
0.00909090909090934 0.199999999999999 \\
0.00909090909090934 0.0999999999999996 \\
0 0.0909090909090903 \\
-0.0909090909090907 0 \\
0 -0.0909090909090903 \\
0.00909090909090934 -0.0999999999999996 \\
0.00909090909090934 -0.199999999999999 \\
0.00909090909090934 -0.3 \\
0.00909090909090934 -0.399999999999999 \\
0.0999999999999999 -0.490909090909091 \\
0.109090909090909 -0.5 \\
0.109090909090909 -0.6 \\
0.2 -0.69090909090909 \\
0.209090909090909 -0.699999999999999 \\
0.3 -0.79090909090909 \\
0.309090909090909 -0.8 \\
0.4 -0.89090909090909 \\
0.5 -0.89090909090909 \\
0.509090909090909 -0.899999999999999 \\
0.6 -0.990909090909091 \\
0.7 -0.990909090909091 \\
0.8 -0.990909090909091 \\
0.9 -0.990909090909091 \\
0.909090909090909 -1 \\
1 -1.09090909090909 \\
1.09090909090909 -1 \\
1.1 -0.990909090909091 \\
1.2 -0.990909090909091 \\
1.3 -0.990909090909091 \\
1.4 -0.990909090909091 \\
1.49090909090909 -0.899999999999999 \\
1.5 -0.89090909090909 \\
1.59090909090909 -0.8 \\
1.6 -0.79090909090909 \\
1.7 -0.79090909090909 \\
1.79090909090909 -0.699999999999999 \\
1.79090909090909 -0.6 \\
1.8 -0.59090909090909 \\
1.89090909090909 -0.5 \\
1.9 -0.490909090909091 \\
1.99090909090909 -0.399999999999999 \\
1.99090909090909 -0.3 \\
1.99090909090909 -0.199999999999999 \\
1.99090909090909 -0.0999999999999996 \\
2 -0.0909090909090903 \\
NaN NaN \\
};


\addplot[dashed,line width=2.0pt,draw=red,forget plot]
table[row sep=crcr]{
x y\\
2 0.0999999999999996 \\
2 0.0999999999999996 \\
2 0.199999999999999 \\
2 0.3 \\
2 0.399999999999999 \\
1.9 0.5 \\
1.9 0.5 \\
1.8 0.6 \\
1.8 0.6 \\
1.8 0.699999999999999 \\
1.7 0.8 \\
1.6 0.8 \\
1.6 0.8 \\
1.5 0.899999999999999 \\
1.5 0.899999999999999 \\
1.4 1 \\
1.3 1 \\
1.2 1 \\
1.1 1 \\
1.1 1 \\
1 1.1 \\
0.9 1 \\
0.9 1 \\
0.8 1 \\
0.7 1 \\
0.6 1 \\
0.5 0.899999999999999 \\
0.5 0.899999999999999 \\
0.4 0.8 \\
0.4 0.8 \\
0.3 0.8 \\
0.2 0.699999999999999 \\
0.2 0.6 \\
0.2 0.6 \\
0.0999999999999999 0.5 \\
0.0999999999999999 0.5 \\
0 0.399999999999999 \\
0 0.3 \\
0 0.199999999999999 \\
0 0.0999999999999996 \\
0 0.0999999999999996 \\
-0.1 0 \\
0 -0.0999999999999996 \\
0 -0.0999999999999996 \\
0 -0.199999999999999 \\
0 -0.3 \\
0 -0.399999999999999 \\
0.0999999999999999 -0.5 \\
0.0999999999999999 -0.5 \\
0.0999999999999999 -0.6 \\
0.2 -0.699999999999999 \\
0.2 -0.699999999999999 \\
0.3 -0.8 \\
0.3 -0.8 \\
0.4 -0.899999999999999 \\
0.5 -0.899999999999999 \\
0.5 -0.899999999999999 \\
0.6 -1 \\
0.7 -1 \\
0.8 -1 \\
0.9 -1 \\
0.9 -1 \\
1 -1.1 \\
1.1 -1 \\
1.1 -1 \\
1.2 -1 \\
1.3 -1 \\
1.4 -1 \\
1.5 -0.899999999999999 \\
1.5 -0.899999999999999 \\
1.6 -0.8 \\
1.6 -0.8 \\
1.7 -0.8 \\
1.8 -0.699999999999999 \\
1.8 -0.6 \\
1.8 -0.6 \\
1.9 -0.5 \\
1.9 -0.5 \\
2 -0.399999999999999 \\
2 -0.3 \\
2 -0.199999999999999 \\
2 -0.0999999999999996 \\
2 -0.0999999999999996 \\
NaN NaN \\
};


\addplot[dashed,line width=2.0pt,draw=green,forget plot]
table[row sep=crcr]{
x y\\
0.0999999899999997 6 \\
0 5.90000001 \\
-1.0000000116861e-008 5.9 \\
-1.0000000116861e-008 5.8 \\
-1.0000000116861e-008 5.7 \\
0 5.69999999 \\
0.0999999899999997 5.6 \\
0.0999999899999997 5.5 \\
0 5.40000001 \\
-1.0000000116861e-008 5.4 \\
-1.0000000116861e-008 5.3 \\
0 5.29999999 \\
0.0999999899999997 5.2 \\
0 5.10000001 \\
-1.0000000116861e-008 5.1 \\
-1.0000000116861e-008 5 \\
0 4.99999999 \\
0.0999999899999997 4.9 \\
0 4.80000001 \\
-1.0000000116861e-008 4.8 \\
-1.0000000116861e-008 4.7 \\
-1.0000000116861e-008 4.6 \\
0 4.59999999 \\
0.0999999899999997 4.5 \\
0.0999999899999997 4.4 \\
0.0999999899999997 4.3 \\
0.0999999899999997 4.2 \\
0.0999999899999997 4.1 \\
0 4.00000001 \\
-1.0000000116861e-008 4 \\
-1.0000000116861e-008 3.9 \\
-1.0000000116861e-008 3.8 \\
-1.0000000116861e-008 3.7 \\
0 3.69999999 \\
0.0999999899999997 3.6 \\
0 3.50000001 \\
-1.0000000116861e-008 3.5 \\
-1.0000000116861e-008 3.4 \\
0 3.39999999 \\
0.0999999899999997 3.3 \\
0.0999999899999997 3.2 \\
0.0999999899999997 3.1 \\
0 3.00000001 \\
-1.0000000116861e-008 3 \\
0 2.99999999 \\
0.0999999899999997 2.9 \\
0 2.80000001 \\
-1.0000000116861e-008 2.8 \\
-1.0000000116861e-008 2.7 \\
-1.0000000116861e-008 2.6 \\
-1.0000000116861e-008 2.5 \\
-1.0000000116861e-008 2.4 \\
-1.0000000116861e-008 2.3 \\
0 2.29999999 \\
0.0999999899999997 2.2 \\
0.0999999899999997 2.1 \\
0 2.00000001 \\
-1.0000000116861e-008 2 \\
-1.0000000116861e-008 1.9 \\
-1.0000000116861e-008 1.8 \\
0 1.79999999 \\
0.0999999899999997 1.7 \\
0.0999999899999997 1.6 \\
0 1.50000001 \\
-1.0000000116861e-008 1.5 \\
-1.0000000116861e-008 1.4 \\
0 1.39999999 \\
0.0999999899999997 1.3 \\
0 1.20000001 \\
-1.0000000116861e-008 1.2 \\
-1.0000000116861e-008 1.1 \\
-1.0000000116861e-008 1 \\
-1.0000000116861e-008 0.899999999999999 \\
0 0.89999999 \\
0.0999999899999997 0.8 \\
0 0.700000009999999 \\
-1.0000000116861e-008 0.699999999999999 \\
-1.0000000116861e-008 0.6 \\
-1.0000000116861e-008 0.5 \\
-1.0000000116861e-008 0.399999999999999 \\
-1.0000000116861e-008 0.3 \\
-1.0000000116861e-008 0.199999999999999 \\
-1.0000000116861e-008 0.0999999999999996 \\
-1.0000000116861e-008 0 \\
-1.0000000116861e-008 -0.0999999999999996 \\
0 -0.10000001 \\
0.0999999899999997 -0.199999999999999 \\
0 -0.29999999 \\
-1.0000000116861e-008 -0.3 \\
-1.0000000116861e-008 -0.399999999999999 \\
-1.0000000116861e-008 -0.5 \\
-1.0000000116861e-008 -0.6 \\
0 -0.60000001 \\
0.0999999899999997 -0.699999999999999 \\
0 -0.79999999 \\
-1.0000000116861e-008 -0.8 \\
-1.0000000116861e-008 -0.899999999999999 \\
-1.0000000116861e-008 -1 \\
-1.0000000116861e-008 -1.1 \\
0 -1.10000001 \\
0.0999999899999997 -1.2 \\
0 -1.29999999 \\
-1.0000000116861e-008 -1.3 \\
-1.0000000116861e-008 -1.4 \\
-1.0000000116861e-008 -1.5 \\
-1.0000000116861e-008 -1.6 \\
0 -1.60000001 \\
0.0999999899999997 -1.7 \\
0 -1.79999999 \\
-1.0000000116861e-008 -1.8 \\
-1.0000000116861e-008 -1.9 \\
-1.0000000116861e-008 -2 \\
0 -2.00000001 \\
0.0999999899999997 -2.1 \\
0 -2.19999999 \\
-1.0000000116861e-008 -2.2 \\
-1.0000000116861e-008 -2.3 \\
-1.0000000116861e-008 -2.4 \\
-1.0000000116861e-008 -2.5 \\
-1.0000000116861e-008 -2.6 \\
-1.0000000116861e-008 -2.7 \\
-1.0000000116861e-008 -2.8 \\
-1.0000000116861e-008 -2.9 \\
-1.0000000116861e-008 -3 \\
0 -3.00000001 \\
0.0999999899999997 -3.1 \\
0 -3.19999999 \\
-1.0000000116861e-008 -3.2 \\
-1.0000000116861e-008 -3.3 \\
-1.0000000116861e-008 -3.4 \\
-1.0000000116861e-008 -3.5 \\
-1.0000000116861e-008 -3.6 \\
-1.0000000116861e-008 -3.7 \\
-1.0000000116861e-008 -3.8 \\
0 -3.80000001 \\
0.0999999899999997 -3.9 \\
0 -3.99999999 \\
-1.0000000116861e-008 -4 \\
-1.0000000116861e-008 -4.1 \\
0 -4.10000001 \\
0.0999999899999997 -4.2 \\
0 -4.29999999 \\
-1.0000000116861e-008 -4.3 \\
0 -4.30000001 \\
0.0999999899999997 -4.4 \\
0.0999999899999997 -4.5 \\
0.0999999899999997 -4.6 \\
0.0999999899999997 -4.7 \\
0 -4.79999999 \\
-1.0000000116861e-008 -4.8 \\
0 -4.80000001 \\
0.0999999899999997 -4.9 \\
0 -4.99999999 \\
-1.0000000116861e-008 -5 \\
-1.0000000116861e-008 -5.1 \\
-1.0000000116861e-008 -5.2 \\
-1.0000000116861e-008 -5.3 \\
-1.0000000116861e-008 -5.4 \\
0 -5.40000001 \\
0.0999999899999997 -5.5 \\
0.0999999899999997 -5.6 \\
0 -5.69999999 \\
-1.0000000116861e-008 -5.7 \\
-1.0000000116861e-008 -5.8 \\
-1.0000000116861e-008 -5.9 \\
0 -5.90000001 \\
0.0999999899999997 -6 \\
NaN NaN \\
};


\addplot[dashed,line width=2.0pt,draw=green,forget plot]
table[row sep=crcr]{
x y\\
0.0909090909090905 6 \\
0 5.90909090909091 \\
-0.00909090909090936 5.9 \\
-0.00909090909090936 5.8 \\
-0.00909090909090936 5.7 \\
0 5.69090909090909 \\
0.0909090909090905 5.6 \\
0.0909090909090905 5.5 \\
0 5.40909090909091 \\
-0.00909090909090936 5.4 \\
-0.00909090909090936 5.3 \\
0 5.29090909090909 \\
0.0909090909090905 5.2 \\
0 5.10909090909091 \\
-0.00909090909090936 5.1 \\
-0.00909090909090936 5 \\
0 4.99090909090909 \\
0.0909090909090905 4.9 \\
0 4.80909090909091 \\
-0.00909090909090936 4.8 \\
-0.00909090909090936 4.7 \\
-0.00909090909090936 4.6 \\
0 4.59090909090909 \\
0.0909090909090905 4.5 \\
0.0909090909090905 4.4 \\
0.0909090909090905 4.3 \\
0.0909090909090905 4.2 \\
0.0909090909090905 4.1 \\
0 4.00909090909091 \\
-0.00909090909090936 4 \\
-0.00909090909090936 3.9 \\
-0.00909090909090936 3.8 \\
-0.00909090909090936 3.7 \\
0 3.69090909090909 \\
0.0909090909090905 3.6 \\
0 3.50909090909091 \\
-0.00909090909090936 3.5 \\
-0.00909090909090936 3.4 \\
0 3.39090909090909 \\
0.0909090909090905 3.3 \\
0.0909090909090905 3.2 \\
0.0909090909090905 3.1 \\
0 3.00909090909091 \\
-0.00909090909090936 3 \\
0 2.99090909090909 \\
0.0909090909090905 2.9 \\
0 2.80909090909091 \\
-0.00909090909090936 2.8 \\
-0.00909090909090936 2.7 \\
-0.00909090909090936 2.6 \\
-0.00909090909090936 2.5 \\
-0.00909090909090936 2.4 \\
-0.00909090909090936 2.3 \\
0 2.29090909090909 \\
0.0909090909090905 2.2 \\
0.0909090909090905 2.1 \\
0 2.00909090909091 \\
-0.00909090909090936 2 \\
-0.00909090909090936 1.9 \\
-0.00909090909090936 1.8 \\
0 1.79090909090909 \\
0.0909090909090905 1.7 \\
0.0909090909090905 1.6 \\
0 1.50909090909091 \\
-0.00909090909090936 1.5 \\
-0.00909090909090936 1.4 \\
0 1.39090909090909 \\
0.0909090909090905 1.3 \\
0 1.20909090909091 \\
-0.00909090909090936 1.2 \\
-0.00909090909090936 1.1 \\
-0.00909090909090936 1 \\
-0.00909090909090936 0.899999999999999 \\
0 0.89090909090909 \\
0.0909090909090905 0.8 \\
0 0.709090909090909 \\
-0.00909090909090936 0.699999999999999 \\
-0.00909090909090936 0.6 \\
-0.00909090909090936 0.5 \\
-0.00909090909090936 0.399999999999999 \\
-0.00909090909090936 0.3 \\
-0.00909090909090936 0.199999999999999 \\
-0.00909090909090936 0.0999999999999996 \\
-0.00909090909090936 0 \\
-0.00909090909090936 -0.0999999999999996 \\
0 -0.109090909090909 \\
0.0909090909090905 -0.199999999999999 \\
0 -0.29090909090909 \\
-0.00909090909090936 -0.3 \\
-0.00909090909090936 -0.399999999999999 \\
-0.00909090909090936 -0.5 \\
-0.00909090909090936 -0.6 \\
0 -0.609090909090909 \\
0.0909090909090905 -0.699999999999999 \\
0 -0.79090909090909 \\
-0.00909090909090936 -0.8 \\
-0.00909090909090936 -0.899999999999999 \\
-0.00909090909090936 -1 \\
-0.00909090909090936 -1.1 \\
0 -1.10909090909091 \\
0.0909090909090905 -1.2 \\
0 -1.29090909090909 \\
-0.00909090909090936 -1.3 \\
-0.00909090909090936 -1.4 \\
-0.00909090909090936 -1.5 \\
-0.00909090909090936 -1.6 \\
0 -1.60909090909091 \\
0.0909090909090905 -1.7 \\
0 -1.79090909090909 \\
-0.00909090909090936 -1.8 \\
-0.00909090909090936 -1.9 \\
-0.00909090909090936 -2 \\
0 -2.00909090909091 \\
0.0909090909090905 -2.1 \\
0 -2.19090909090909 \\
-0.00909090909090936 -2.2 \\
-0.00909090909090936 -2.3 \\
-0.00909090909090936 -2.4 \\
-0.00909090909090936 -2.5 \\
-0.00909090909090936 -2.6 \\
-0.00909090909090936 -2.7 \\
-0.00909090909090936 -2.8 \\
-0.00909090909090936 -2.9 \\
-0.00909090909090936 -3 \\
0 -3.00909090909091 \\
0.0909090909090905 -3.1 \\
0 -3.19090909090909 \\
-0.00909090909090936 -3.2 \\
-0.00909090909090936 -3.3 \\
-0.00909090909090936 -3.4 \\
-0.00909090909090936 -3.5 \\
-0.00909090909090936 -3.6 \\
-0.00909090909090936 -3.7 \\
-0.00909090909090936 -3.8 \\
0 -3.80909090909091 \\
0.0909090909090905 -3.9 \\
0 -3.99090909090909 \\
-0.00909090909090936 -4 \\
-0.00909090909090936 -4.1 \\
0 -4.10909090909091 \\
0.0909090909090905 -4.2 \\
0 -4.29090909090909 \\
-0.00909090909090936 -4.3 \\
0 -4.30909090909091 \\
0.0909090909090905 -4.4 \\
0.0909090909090905 -4.5 \\
0.0909090909090905 -4.6 \\
0.0909090909090905 -4.7 \\
0 -4.79090909090909 \\
-0.00909090909090936 -4.8 \\
0 -4.80909090909091 \\
0.0909090909090905 -4.9 \\
0 -4.99090909090909 \\
-0.00909090909090936 -5 \\
-0.00909090909090936 -5.1 \\
-0.00909090909090936 -5.2 \\
-0.00909090909090936 -5.3 \\
-0.00909090909090936 -5.4 \\
0 -5.40909090909091 \\
0.0909090909090905 -5.5 \\
0.0909090909090905 -5.6 \\
0 -5.69090909090909 \\
-0.00909090909090936 -5.7 \\
-0.00909090909090936 -5.8 \\
-0.00909090909090936 -5.9 \\
0 -5.90909090909091 \\
0.0909090909090905 -6 \\
NaN NaN \\
};


\addplot[dashed,line width=2.0pt,draw=green,forget plot]
table[row sep=crcr]{
x y\\
0.0818181818181819 6 \\
0 5.91818181818182 \\
-0.018181818181818 5.9 \\
-0.018181818181818 5.8 \\
-0.018181818181818 5.7 \\
0 5.68181818181818 \\
0.0818181818181819 5.6 \\
0.0818181818181819 5.5 \\
0 5.41818181818182 \\
-0.018181818181818 5.4 \\
-0.018181818181818 5.3 \\
0 5.28181818181818 \\
0.0818181818181819 5.2 \\
0 5.11818181818182 \\
-0.018181818181818 5.1 \\
-0.018181818181818 5 \\
0 4.98181818181818 \\
0.0818181818181819 4.9 \\
0 4.81818181818182 \\
-0.018181818181818 4.8 \\
-0.018181818181818 4.7 \\
-0.018181818181818 4.6 \\
0 4.58181818181818 \\
0.0818181818181819 4.5 \\
0.0818181818181819 4.4 \\
0.0818181818181819 4.3 \\
0.0818181818181819 4.2 \\
0.0818181818181819 4.1 \\
0 4.01818181818182 \\
-0.018181818181818 4 \\
-0.018181818181818 3.9 \\
-0.018181818181818 3.8 \\
-0.018181818181818 3.7 \\
0 3.68181818181818 \\
0.0818181818181819 3.6 \\
0 3.51818181818182 \\
-0.018181818181818 3.5 \\
-0.018181818181818 3.4 \\
0 3.38181818181818 \\
0.0818181818181819 3.3 \\
0.0818181818181819 3.2 \\
0.0818181818181819 3.1 \\
0 3.01818181818182 \\
-0.018181818181818 3 \\
0 2.98181818181818 \\
0.0818181818181819 2.9 \\
0 2.81818181818182 \\
-0.018181818181818 2.8 \\
-0.018181818181818 2.7 \\
-0.018181818181818 2.6 \\
-0.018181818181818 2.5 \\
-0.018181818181818 2.4 \\
-0.018181818181818 2.3 \\
0 2.28181818181818 \\
0.0818181818181819 2.2 \\
0.0818181818181819 2.1 \\
0 2.01818181818182 \\
-0.018181818181818 2 \\
-0.018181818181818 1.9 \\
-0.018181818181818 1.8 \\
0 1.78181818181818 \\
0.0818181818181819 1.7 \\
0.0818181818181819 1.6 \\
0 1.51818181818182 \\
-0.018181818181818 1.5 \\
-0.018181818181818 1.4 \\
0 1.38181818181818 \\
0.0818181818181819 1.3 \\
0 1.21818181818182 \\
-0.018181818181818 1.2 \\
-0.018181818181818 1.1 \\
-0.018181818181818 1 \\
-0.018181818181818 0.899999999999999 \\
0 0.881818181818181 \\
0.0818181818181819 0.8 \\
0 0.718181818181818 \\
-0.018181818181818 0.699999999999999 \\
-0.018181818181818 0.6 \\
-0.018181818181818 0.5 \\
-0.018181818181818 0.399999999999999 \\
-0.018181818181818 0.3 \\
-0.018181818181818 0.199999999999999 \\
-0.018181818181818 0.0999999999999996 \\
-0.018181818181818 0 \\
-0.018181818181818 -0.0999999999999996 \\
0 -0.118181818181818 \\
0.0818181818181819 -0.199999999999999 \\
0 -0.281818181818182 \\
-0.018181818181818 -0.3 \\
-0.018181818181818 -0.399999999999999 \\
-0.018181818181818 -0.5 \\
-0.018181818181818 -0.6 \\
0 -0.618181818181818 \\
0.0818181818181819 -0.699999999999999 \\
0 -0.781818181818182 \\
-0.018181818181818 -0.8 \\
-0.018181818181818 -0.899999999999999 \\
-0.018181818181818 -1 \\
-0.018181818181818 -1.1 \\
0 -1.11818181818182 \\
0.0818181818181819 -1.2 \\
0 -1.28181818181818 \\
-0.018181818181818 -1.3 \\
-0.018181818181818 -1.4 \\
-0.018181818181818 -1.5 \\
-0.018181818181818 -1.6 \\
0 -1.61818181818182 \\
0.0818181818181819 -1.7 \\
0 -1.78181818181818 \\
-0.018181818181818 -1.8 \\
-0.018181818181818 -1.9 \\
-0.018181818181818 -2 \\
0 -2.01818181818182 \\
0.0818181818181819 -2.1 \\
0 -2.18181818181818 \\
-0.018181818181818 -2.2 \\
-0.018181818181818 -2.3 \\
-0.018181818181818 -2.4 \\
-0.018181818181818 -2.5 \\
-0.018181818181818 -2.6 \\
-0.018181818181818 -2.7 \\
-0.018181818181818 -2.8 \\
-0.018181818181818 -2.9 \\
-0.018181818181818 -3 \\
0 -3.01818181818182 \\
0.0818181818181819 -3.1 \\
0 -3.18181818181818 \\
-0.018181818181818 -3.2 \\
-0.018181818181818 -3.3 \\
-0.018181818181818 -3.4 \\
-0.018181818181818 -3.5 \\
-0.018181818181818 -3.6 \\
-0.018181818181818 -3.7 \\
-0.018181818181818 -3.8 \\
0 -3.81818181818182 \\
0.0818181818181819 -3.9 \\
0 -3.98181818181818 \\
-0.018181818181818 -4 \\
-0.018181818181818 -4.1 \\
0 -4.11818181818182 \\
0.0818181818181819 -4.2 \\
0 -4.28181818181818 \\
-0.018181818181818 -4.3 \\
0 -4.31818181818182 \\
0.0818181818181819 -4.4 \\
0.0818181818181819 -4.5 \\
0.0818181818181819 -4.6 \\
0.0818181818181819 -4.7 \\
0 -4.78181818181818 \\
-0.018181818181818 -4.8 \\
0 -4.81818181818182 \\
0.0818181818181819 -4.9 \\
0 -4.98181818181818 \\
-0.018181818181818 -5 \\
-0.018181818181818 -5.1 \\
-0.018181818181818 -5.2 \\
-0.018181818181818 -5.3 \\
-0.018181818181818 -5.4 \\
0 -5.41818181818182 \\
0.0818181818181819 -5.5 \\
0.0818181818181819 -5.6 \\
0 -5.68181818181818 \\
-0.018181818181818 -5.7 \\
-0.018181818181818 -5.8 \\
-0.018181818181818 -5.9 \\
0 -5.91818181818182 \\
0.0818181818181819 -6 \\
NaN NaN \\
};


\addplot[dashed,line width=2.0pt,draw=green,forget plot]
table[row sep=crcr]{
x y\\
0.0727272727272726 6 \\
0 5.92727272727273 \\
-0.0272727272727274 5.9 \\
-0.0272727272727274 5.8 \\
-0.0272727272727274 5.7 \\
0 5.67272727272727 \\
0.0727272727272726 5.6 \\
0.0727272727272726 5.5 \\
0 5.42727272727273 \\
-0.0272727272727274 5.4 \\
-0.0272727272727274 5.3 \\
0 5.27272727272727 \\
0.0727272727272726 5.2 \\
0 5.12727272727273 \\
-0.0272727272727274 5.1 \\
-0.0272727272727274 5 \\
0 4.97272727272727 \\
0.0727272727272726 4.9 \\
0 4.82727272727273 \\
-0.0272727272727274 4.8 \\
-0.0272727272727274 4.7 \\
-0.0272727272727274 4.6 \\
0 4.57272727272727 \\
0.0727272727272726 4.5 \\
0.0727272727272726 4.4 \\
0.0727272727272726 4.3 \\
0.0727272727272726 4.2 \\
0.0727272727272726 4.1 \\
0 4.02727272727273 \\
-0.0272727272727274 4 \\
-0.0272727272727274 3.9 \\
-0.0272727272727274 3.8 \\
-0.0272727272727274 3.7 \\
0 3.67272727272727 \\
0.0727272727272726 3.6 \\
0 3.52727272727273 \\
-0.0272727272727274 3.5 \\
-0.0272727272727274 3.4 \\
0 3.37272727272727 \\
0.0727272727272726 3.3 \\
0.0727272727272726 3.2 \\
0.0727272727272726 3.1 \\
0 3.02727272727273 \\
-0.0272727272727274 3 \\
0 2.97272727272727 \\
0.0727272727272726 2.9 \\
0 2.82727272727273 \\
-0.0272727272727274 2.8 \\
-0.0272727272727274 2.7 \\
-0.0272727272727274 2.6 \\
-0.0272727272727274 2.5 \\
-0.0272727272727274 2.4 \\
-0.0272727272727274 2.3 \\
0 2.27272727272727 \\
0.0727272727272726 2.2 \\
0.0727272727272726 2.1 \\
0 2.02727272727273 \\
-0.0272727272727274 2 \\
-0.0272727272727274 1.9 \\
-0.0272727272727274 1.8 \\
0 1.77272727272727 \\
0.0727272727272726 1.7 \\
0.0727272727272726 1.6 \\
0 1.52727272727273 \\
-0.0272727272727274 1.5 \\
-0.0272727272727274 1.4 \\
0 1.37272727272727 \\
0.0727272727272726 1.3 \\
0 1.22727272727273 \\
-0.0272727272727274 1.2 \\
-0.0272727272727274 1.1 \\
-0.0272727272727274 1 \\
-0.0272727272727274 0.899999999999999 \\
0 0.872727272727273 \\
0.0727272727272726 0.8 \\
0 0.727272727272726 \\
-0.0272727272727274 0.699999999999999 \\
-0.0272727272727274 0.6 \\
-0.0272727272727274 0.5 \\
-0.0272727272727274 0.399999999999999 \\
-0.0272727272727274 0.3 \\
-0.0272727272727274 0.199999999999999 \\
-0.0272727272727274 0.0999999999999996 \\
-0.0272727272727274 0 \\
-0.0272727272727274 -0.0999999999999996 \\
0 -0.127272727272727 \\
0.0727272727272726 -0.199999999999999 \\
0 -0.272727272727272 \\
-0.0272727272727274 -0.3 \\
-0.0272727272727274 -0.399999999999999 \\
-0.0272727272727274 -0.5 \\
-0.0272727272727274 -0.6 \\
0 -0.627272727272727 \\
0.0727272727272726 -0.699999999999999 \\
0 -0.772727272727272 \\
-0.0272727272727274 -0.8 \\
-0.0272727272727274 -0.899999999999999 \\
-0.0272727272727274 -1 \\
-0.0272727272727274 -1.1 \\
0 -1.12727272727273 \\
0.0727272727272726 -1.2 \\
0 -1.27272727272727 \\
-0.0272727272727274 -1.3 \\
-0.0272727272727274 -1.4 \\
-0.0272727272727274 -1.5 \\
-0.0272727272727274 -1.6 \\
0 -1.62727272727273 \\
0.0727272727272726 -1.7 \\
0 -1.77272727272727 \\
-0.0272727272727274 -1.8 \\
-0.0272727272727274 -1.9 \\
-0.0272727272727274 -2 \\
0 -2.02727272727273 \\
0.0727272727272726 -2.1 \\
0 -2.17272727272727 \\
-0.0272727272727274 -2.2 \\
-0.0272727272727274 -2.3 \\
-0.0272727272727274 -2.4 \\
-0.0272727272727274 -2.5 \\
-0.0272727272727274 -2.6 \\
-0.0272727272727274 -2.7 \\
-0.0272727272727274 -2.8 \\
-0.0272727272727274 -2.9 \\
-0.0272727272727274 -3 \\
0 -3.02727272727273 \\
0.0727272727272726 -3.1 \\
0 -3.17272727272727 \\
-0.0272727272727274 -3.2 \\
-0.0272727272727274 -3.3 \\
-0.0272727272727274 -3.4 \\
-0.0272727272727274 -3.5 \\
-0.0272727272727274 -3.6 \\
-0.0272727272727274 -3.7 \\
-0.0272727272727274 -3.8 \\
0 -3.82727272727273 \\
0.0727272727272726 -3.9 \\
0 -3.97272727272727 \\
-0.0272727272727274 -4 \\
-0.0272727272727274 -4.1 \\
0 -4.12727272727273 \\
0.0727272727272726 -4.2 \\
0 -4.27272727272727 \\
-0.0272727272727274 -4.3 \\
0 -4.32727272727273 \\
0.0727272727272726 -4.4 \\
0.0727272727272726 -4.5 \\
0.0727272727272726 -4.6 \\
0.0727272727272726 -4.7 \\
0 -4.77272727272727 \\
-0.0272727272727274 -4.8 \\
0 -4.82727272727273 \\
0.0727272727272726 -4.9 \\
0 -4.97272727272727 \\
-0.0272727272727274 -5 \\
-0.0272727272727274 -5.1 \\
-0.0272727272727274 -5.2 \\
-0.0272727272727274 -5.3 \\
-0.0272727272727274 -5.4 \\
0 -5.42727272727273 \\
0.0727272727272726 -5.5 \\
0.0727272727272726 -5.6 \\
0 -5.67272727272727 \\
-0.0272727272727274 -5.7 \\
-0.0272727272727274 -5.8 \\
-0.0272727272727274 -5.9 \\
0 -5.92727272727273 \\
0.0727272727272726 -6 \\
NaN NaN \\
};


\addplot[dashed,line width=2.0pt,draw=green,forget plot]
table[row sep=crcr]{
x y\\
0.0636363636363632 6 \\
0 5.93636363636364 \\
-0.0363636363636367 5.9 \\
-0.0363636363636367 5.8 \\
-0.0363636363636367 5.7 \\
0 5.66363636363636 \\
0.0636363636363632 5.6 \\
0.0636363636363632 5.5 \\
0 5.43636363636364 \\
-0.0363636363636367 5.4 \\
-0.0363636363636367 5.3 \\
0 5.26363636363636 \\
0.0636363636363632 5.2 \\
0 5.13636363636364 \\
-0.0363636363636367 5.1 \\
-0.0363636363636367 5 \\
0 4.96363636363636 \\
0.0636363636363632 4.9 \\
0 4.83636363636364 \\
-0.0363636363636367 4.8 \\
-0.0363636363636367 4.7 \\
-0.0363636363636367 4.6 \\
0 4.56363636363636 \\
0.0636363636363632 4.5 \\
0.0636363636363632 4.4 \\
0.0636363636363632 4.3 \\
0.0636363636363632 4.2 \\
0.0636363636363632 4.1 \\
0 4.03636363636364 \\
-0.0363636363636367 4 \\
-0.0363636363636367 3.9 \\
-0.0363636363636367 3.8 \\
-0.0363636363636367 3.7 \\
0 3.66363636363636 \\
0.0636363636363632 3.6 \\
0 3.53636363636364 \\
-0.0363636363636367 3.5 \\
-0.0363636363636367 3.4 \\
0 3.36363636363636 \\
0.0636363636363632 3.3 \\
0.0636363636363632 3.2 \\
0.0636363636363632 3.1 \\
0 3.03636363636364 \\
-0.0363636363636367 3 \\
0 2.96363636363636 \\
0.0636363636363632 2.9 \\
0 2.83636363636364 \\
-0.0363636363636367 2.8 \\
-0.0363636363636367 2.7 \\
-0.0363636363636367 2.6 \\
-0.0363636363636367 2.5 \\
-0.0363636363636367 2.4 \\
-0.0363636363636367 2.3 \\
0 2.26363636363636 \\
0.0636363636363632 2.2 \\
0.0636363636363632 2.1 \\
0 2.03636363636364 \\
-0.0363636363636367 2 \\
-0.0363636363636367 1.9 \\
-0.0363636363636367 1.8 \\
0 1.76363636363636 \\
0.0636363636363632 1.7 \\
0.0636363636363632 1.6 \\
0 1.53636363636364 \\
-0.0363636363636367 1.5 \\
-0.0363636363636367 1.4 \\
0 1.36363636363636 \\
0.0636363636363632 1.3 \\
0 1.23636363636364 \\
-0.0363636363636367 1.2 \\
-0.0363636363636367 1.1 \\
-0.0363636363636367 1 \\
-0.0363636363636367 0.899999999999999 \\
0 0.863636363636364 \\
0.0636363636363632 0.8 \\
0 0.736363636363635 \\
-0.0363636363636367 0.699999999999999 \\
-0.0363636363636367 0.6 \\
-0.0363636363636367 0.5 \\
-0.0363636363636367 0.399999999999999 \\
-0.0363636363636367 0.3 \\
-0.0363636363636367 0.199999999999999 \\
-0.0363636363636367 0.0999999999999996 \\
-0.0363636363636367 0 \\
-0.0363636363636367 -0.0999999999999996 \\
0 -0.136363636363636 \\
0.0636363636363632 -0.199999999999999 \\
0 -0.263636363636363 \\
-0.0363636363636367 -0.3 \\
-0.0363636363636367 -0.399999999999999 \\
-0.0363636363636367 -0.5 \\
-0.0363636363636367 -0.6 \\
0 -0.636363636363636 \\
0.0636363636363632 -0.699999999999999 \\
0 -0.763636363636363 \\
-0.0363636363636367 -0.8 \\
-0.0363636363636367 -0.899999999999999 \\
-0.0363636363636367 -1 \\
-0.0363636363636367 -1.1 \\
0 -1.13636363636364 \\
0.0636363636363632 -1.2 \\
0 -1.26363636363636 \\
-0.0363636363636367 -1.3 \\
-0.0363636363636367 -1.4 \\
-0.0363636363636367 -1.5 \\
-0.0363636363636367 -1.6 \\
0 -1.63636363636364 \\
0.0636363636363632 -1.7 \\
0 -1.76363636363636 \\
-0.0363636363636367 -1.8 \\
-0.0363636363636367 -1.9 \\
-0.0363636363636367 -2 \\
0 -2.03636363636364 \\
0.0636363636363632 -2.1 \\
0 -2.16363636363636 \\
-0.0363636363636367 -2.2 \\
-0.0363636363636367 -2.3 \\
-0.0363636363636367 -2.4 \\
-0.0363636363636367 -2.5 \\
-0.0363636363636367 -2.6 \\
-0.0363636363636367 -2.7 \\
-0.0363636363636367 -2.8 \\
-0.0363636363636367 -2.9 \\
-0.0363636363636367 -3 \\
0 -3.03636363636364 \\
0.0636363636363632 -3.1 \\
0 -3.16363636363636 \\
-0.0363636363636367 -3.2 \\
-0.0363636363636367 -3.3 \\
-0.0363636363636367 -3.4 \\
-0.0363636363636367 -3.5 \\
-0.0363636363636367 -3.6 \\
-0.0363636363636367 -3.7 \\
-0.0363636363636367 -3.8 \\
0 -3.83636363636364 \\
0.0636363636363632 -3.9 \\
0 -3.96363636363636 \\
-0.0363636363636367 -4 \\
-0.0363636363636367 -4.1 \\
0 -4.13636363636364 \\
0.0636363636363632 -4.2 \\
0 -4.26363636363636 \\
-0.0363636363636367 -4.3 \\
0 -4.33636363636364 \\
0.0636363636363632 -4.4 \\
0.0636363636363632 -4.5 \\
0.0636363636363632 -4.6 \\
0.0636363636363632 -4.7 \\
0 -4.76363636363636 \\
-0.0363636363636367 -4.8 \\
0 -4.83636363636364 \\
0.0636363636363632 -4.9 \\
0 -4.96363636363636 \\
-0.0363636363636367 -5 \\
-0.0363636363636367 -5.1 \\
-0.0363636363636367 -5.2 \\
-0.0363636363636367 -5.3 \\
-0.0363636363636367 -5.4 \\
0 -5.43636363636364 \\
0.0636363636363632 -5.5 \\
0.0636363636363632 -5.6 \\
0 -5.66363636363636 \\
-0.0363636363636367 -5.7 \\
-0.0363636363636367 -5.8 \\
-0.0363636363636367 -5.9 \\
0 -5.93636363636364 \\
0.0636363636363632 -6 \\
NaN NaN \\
};


\addplot[dashed,line width=2.0pt,draw=green,forget plot]
table[row sep=crcr]{
x y\\
0.0545454545454546 6 \\
0 5.94545454545455 \\
-0.0454545454545454 5.9 \\
-0.0454545454545454 5.8 \\
-0.0454545454545454 5.7 \\
0 5.65454545454545 \\
0.0545454545454546 5.6 \\
0.0545454545454546 5.5 \\
0 5.44545454545455 \\
-0.0454545454545454 5.4 \\
-0.0454545454545454 5.3 \\
0 5.25454545454545 \\
0.0545454545454546 5.2 \\
0 5.14545454545455 \\
-0.0454545454545454 5.1 \\
-0.0454545454545454 5 \\
0 4.95454545454546 \\
0.0545454545454546 4.9 \\
0 4.84545454545455 \\
-0.0454545454545454 4.8 \\
-0.0454545454545454 4.7 \\
-0.0454545454545454 4.6 \\
0 4.55454545454545 \\
0.0545454545454546 4.5 \\
0.0545454545454546 4.4 \\
0.0545454545454546 4.3 \\
0.0545454545454546 4.2 \\
0.0545454545454546 4.1 \\
0 4.04545454545455 \\
-0.0454545454545454 4 \\
-0.0454545454545454 3.9 \\
-0.0454545454545454 3.8 \\
-0.0454545454545454 3.7 \\
0 3.65454545454545 \\
0.0545454545454546 3.6 \\
0 3.54545454545455 \\
-0.0454545454545454 3.5 \\
-0.0454545454545454 3.4 \\
0 3.35454545454545 \\
0.0545454545454546 3.3 \\
0.0545454545454546 3.2 \\
0.0545454545454546 3.1 \\
0 3.04545454545455 \\
-0.0454545454545454 3 \\
0 2.95454545454545 \\
0.0545454545454546 2.9 \\
0 2.84545454545455 \\
-0.0454545454545454 2.8 \\
-0.0454545454545454 2.7 \\
-0.0454545454545454 2.6 \\
-0.0454545454545454 2.5 \\
-0.0454545454545454 2.4 \\
-0.0454545454545454 2.3 \\
0 2.25454545454545 \\
0.0545454545454546 2.2 \\
0.0545454545454546 2.1 \\
0 2.04545454545455 \\
-0.0454545454545454 2 \\
-0.0454545454545454 1.9 \\
-0.0454545454545454 1.8 \\
0 1.75454545454545 \\
0.0545454545454546 1.7 \\
0.0545454545454546 1.6 \\
0 1.54545454545455 \\
-0.0454545454545454 1.5 \\
-0.0454545454545454 1.4 \\
0 1.35454545454545 \\
0.0545454545454546 1.3 \\
0 1.24545454545454 \\
-0.0454545454545454 1.2 \\
-0.0454545454545454 1.1 \\
-0.0454545454545454 1 \\
-0.0454545454545454 0.899999999999999 \\
0 0.854545454545454 \\
0.0545454545454546 0.8 \\
0 0.745454545454545 \\
-0.0454545454545454 0.699999999999999 \\
-0.0454545454545454 0.6 \\
-0.0454545454545454 0.5 \\
-0.0454545454545454 0.399999999999999 \\
-0.0454545454545454 0.3 \\
-0.0454545454545454 0.199999999999999 \\
-0.0454545454545454 0.0999999999999996 \\
-0.0454545454545454 0 \\
-0.0454545454545454 -0.0999999999999996 \\
0 -0.145454545454545 \\
0.0545454545454546 -0.199999999999999 \\
0 -0.254545454545454 \\
-0.0454545454545454 -0.3 \\
-0.0454545454545454 -0.399999999999999 \\
-0.0454545454545454 -0.5 \\
-0.0454545454545454 -0.6 \\
0 -0.645454545454545 \\
0.0545454545454546 -0.699999999999999 \\
0 -0.754545454545454 \\
-0.0454545454545454 -0.8 \\
-0.0454545454545454 -0.899999999999999 \\
-0.0454545454545454 -1 \\
-0.0454545454545454 -1.1 \\
0 -1.14545454545454 \\
0.0545454545454546 -1.2 \\
0 -1.25454545454545 \\
-0.0454545454545454 -1.3 \\
-0.0454545454545454 -1.4 \\
-0.0454545454545454 -1.5 \\
-0.0454545454545454 -1.6 \\
0 -1.64545454545455 \\
0.0545454545454546 -1.7 \\
0 -1.75454545454545 \\
-0.0454545454545454 -1.8 \\
-0.0454545454545454 -1.9 \\
-0.0454545454545454 -2 \\
0 -2.04545454545455 \\
0.0545454545454546 -2.1 \\
0 -2.15454545454545 \\
-0.0454545454545454 -2.2 \\
-0.0454545454545454 -2.3 \\
-0.0454545454545454 -2.4 \\
-0.0454545454545454 -2.5 \\
-0.0454545454545454 -2.6 \\
-0.0454545454545454 -2.7 \\
-0.0454545454545454 -2.8 \\
-0.0454545454545454 -2.9 \\
-0.0454545454545454 -3 \\
0 -3.04545454545455 \\
0.0545454545454546 -3.1 \\
0 -3.15454545454545 \\
-0.0454545454545454 -3.2 \\
-0.0454545454545454 -3.3 \\
-0.0454545454545454 -3.4 \\
-0.0454545454545454 -3.5 \\
-0.0454545454545454 -3.6 \\
-0.0454545454545454 -3.7 \\
-0.0454545454545454 -3.8 \\
0 -3.84545454545455 \\
0.0545454545454546 -3.9 \\
0 -3.95454545454545 \\
-0.0454545454545454 -4 \\
-0.0454545454545454 -4.1 \\
0 -4.14545454545455 \\
0.0545454545454546 -4.2 \\
0 -4.25454545454545 \\
-0.0454545454545454 -4.3 \\
0 -4.34545454545455 \\
0.0545454545454546 -4.4 \\
0.0545454545454546 -4.5 \\
0.0545454545454546 -4.6 \\
0.0545454545454546 -4.7 \\
0 -4.75454545454545 \\
-0.0454545454545454 -4.8 \\
0 -4.84545454545455 \\
0.0545454545454546 -4.9 \\
0 -4.95454545454546 \\
-0.0454545454545454 -5 \\
-0.0454545454545454 -5.1 \\
-0.0454545454545454 -5.2 \\
-0.0454545454545454 -5.3 \\
-0.0454545454545454 -5.4 \\
0 -5.44545454545455 \\
0.0545454545454546 -5.5 \\
0.0545454545454546 -5.6 \\
0 -5.65454545454545 \\
-0.0454545454545454 -5.7 \\
-0.0454545454545454 -5.8 \\
-0.0454545454545454 -5.9 \\
0 -5.94545454545455 \\
0.0545454545454546 -6 \\
NaN NaN \\
};


\addplot[dashed,line width=2.0pt,draw=green,forget plot]
table[row sep=crcr]{
x y\\
0.0454545454545453 6 \\
0 5.95454545454546 \\
-0.0545454545454547 5.9 \\
-0.0545454545454547 5.8 \\
-0.0545454545454547 5.7 \\
0 5.64545454545455 \\
0.0454545454545453 5.6 \\
0.0454545454545453 5.5 \\
0 5.45454545454546 \\
-0.0545454545454547 5.4 \\
-0.0545454545454547 5.3 \\
0 5.24545454545455 \\
0.0454545454545453 5.2 \\
0 5.15454545454545 \\
-0.0545454545454547 5.1 \\
-0.0545454545454547 5 \\
0 4.94545454545455 \\
0.0454545454545453 4.9 \\
0 4.85454545454545 \\
-0.0545454545454547 4.8 \\
-0.0545454545454547 4.7 \\
-0.0545454545454547 4.6 \\
0 4.54545454545455 \\
0.0454545454545453 4.5 \\
0.0454545454545453 4.4 \\
0.0454545454545453 4.3 \\
0.0454545454545453 4.2 \\
0.0454545454545453 4.1 \\
0 4.05454545454545 \\
-0.0545454545454547 4 \\
-0.0545454545454547 3.9 \\
-0.0545454545454547 3.8 \\
-0.0545454545454547 3.7 \\
0 3.64545454545455 \\
0.0454545454545453 3.6 \\
0 3.55454545454545 \\
-0.0545454545454547 3.5 \\
-0.0545454545454547 3.4 \\
0 3.34545454545455 \\
0.0454545454545453 3.3 \\
0.0454545454545453 3.2 \\
0.0454545454545453 3.1 \\
0 3.05454545454545 \\
-0.0545454545454547 3 \\
0 2.94545454545455 \\
0.0454545454545453 2.9 \\
0 2.85454545454545 \\
-0.0545454545454547 2.8 \\
-0.0545454545454547 2.7 \\
-0.0545454545454547 2.6 \\
-0.0545454545454547 2.5 \\
-0.0545454545454547 2.4 \\
-0.0545454545454547 2.3 \\
0 2.24545454545455 \\
0.0454545454545453 2.2 \\
0.0454545454545453 2.1 \\
0 2.05454545454545 \\
-0.0545454545454547 2 \\
-0.0545454545454547 1.9 \\
-0.0545454545454547 1.8 \\
0 1.74545454545455 \\
0.0454545454545453 1.7 \\
0.0454545454545453 1.6 \\
0 1.55454545454545 \\
-0.0545454545454547 1.5 \\
-0.0545454545454547 1.4 \\
0 1.34545454545455 \\
0.0454545454545453 1.3 \\
0 1.25454545454545 \\
-0.0545454545454547 1.2 \\
-0.0545454545454547 1.1 \\
-0.0545454545454547 1 \\
-0.0545454545454547 0.899999999999999 \\
0 0.845454545454545 \\
0.0454545454545453 0.8 \\
0 0.754545454545454 \\
-0.0545454545454547 0.699999999999999 \\
-0.0545454545454547 0.6 \\
-0.0545454545454547 0.5 \\
-0.0545454545454547 0.399999999999999 \\
-0.0545454545454547 0.3 \\
-0.0545454545454547 0.199999999999999 \\
-0.0545454545454547 0.0999999999999996 \\
-0.0545454545454547 0 \\
-0.0545454545454547 -0.0999999999999996 \\
0 -0.154545454545454 \\
0.0454545454545453 -0.199999999999999 \\
0 -0.245454545454545 \\
-0.0545454545454547 -0.3 \\
-0.0545454545454547 -0.399999999999999 \\
-0.0545454545454547 -0.5 \\
-0.0545454545454547 -0.6 \\
0 -0.654545454545454 \\
0.0454545454545453 -0.699999999999999 \\
0 -0.745454545454545 \\
-0.0545454545454547 -0.8 \\
-0.0545454545454547 -0.899999999999999 \\
-0.0545454545454547 -1 \\
-0.0545454545454547 -1.1 \\
0 -1.15454545454545 \\
0.0454545454545453 -1.2 \\
0 -1.24545454545454 \\
-0.0545454545454547 -1.3 \\
-0.0545454545454547 -1.4 \\
-0.0545454545454547 -1.5 \\
-0.0545454545454547 -1.6 \\
0 -1.65454545454545 \\
0.0454545454545453 -1.7 \\
0 -1.74545454545455 \\
-0.0545454545454547 -1.8 \\
-0.0545454545454547 -1.9 \\
-0.0545454545454547 -2 \\
0 -2.05454545454545 \\
0.0454545454545453 -2.1 \\
0 -2.14545454545455 \\
-0.0545454545454547 -2.2 \\
-0.0545454545454547 -2.3 \\
-0.0545454545454547 -2.4 \\
-0.0545454545454547 -2.5 \\
-0.0545454545454547 -2.6 \\
-0.0545454545454547 -2.7 \\
-0.0545454545454547 -2.8 \\
-0.0545454545454547 -2.9 \\
-0.0545454545454547 -3 \\
0 -3.05454545454545 \\
0.0454545454545453 -3.1 \\
0 -3.14545454545455 \\
-0.0545454545454547 -3.2 \\
-0.0545454545454547 -3.3 \\
-0.0545454545454547 -3.4 \\
-0.0545454545454547 -3.5 \\
-0.0545454545454547 -3.6 \\
-0.0545454545454547 -3.7 \\
-0.0545454545454547 -3.8 \\
0 -3.85454545454545 \\
0.0454545454545453 -3.9 \\
0 -3.94545454545455 \\
-0.0545454545454547 -4 \\
-0.0545454545454547 -4.1 \\
0 -4.15454545454545 \\
0.0454545454545453 -4.2 \\
0 -4.24545454545455 \\
-0.0545454545454547 -4.3 \\
0 -4.35454545454545 \\
0.0454545454545453 -4.4 \\
0.0454545454545453 -4.5 \\
0.0454545454545453 -4.6 \\
0.0454545454545453 -4.7 \\
0 -4.74545454545455 \\
-0.0545454545454547 -4.8 \\
0 -4.85454545454545 \\
0.0454545454545453 -4.9 \\
0 -4.94545454545455 \\
-0.0545454545454547 -5 \\
-0.0545454545454547 -5.1 \\
-0.0545454545454547 -5.2 \\
-0.0545454545454547 -5.3 \\
-0.0545454545454547 -5.4 \\
0 -5.45454545454546 \\
0.0454545454545453 -5.5 \\
0.0454545454545453 -5.6 \\
0 -5.64545454545454 \\
-0.0545454545454547 -5.7 \\
-0.0545454545454547 -5.8 \\
-0.0545454545454547 -5.9 \\
0 -5.95454545454546 \\
0.0454545454545453 -6 \\
NaN NaN \\
};


\addplot[dashed,line width=2.0pt,draw=green,forget plot]
table[row sep=crcr]{
x y\\
0.0363636363636366 6 \\
0 5.96363636363636 \\
-0.0636363636363634 5.9 \\
-0.0636363636363634 5.8 \\
-0.0636363636363634 5.7 \\
0 5.63636363636364 \\
0.0363636363636366 5.6 \\
0.0363636363636366 5.5 \\
0 5.46363636363636 \\
-0.0636363636363634 5.4 \\
-0.0636363636363634 5.3 \\
0 5.23636363636364 \\
0.0363636363636366 5.2 \\
0 5.16363636363636 \\
-0.0636363636363634 5.1 \\
-0.0636363636363634 5 \\
0 4.93636363636364 \\
0.0363636363636366 4.9 \\
0 4.86363636363636 \\
-0.0636363636363634 4.8 \\
-0.0636363636363634 4.7 \\
-0.0636363636363634 4.6 \\
0 4.53636363636364 \\
0.0363636363636366 4.5 \\
0.0363636363636366 4.4 \\
0.0363636363636366 4.3 \\
0.0363636363636366 4.2 \\
0.0363636363636366 4.1 \\
0 4.06363636363636 \\
-0.0636363636363634 4 \\
-0.0636363636363634 3.9 \\
-0.0636363636363634 3.8 \\
-0.0636363636363634 3.7 \\
0 3.63636363636364 \\
0.0363636363636366 3.6 \\
0 3.56363636363636 \\
-0.0636363636363634 3.5 \\
-0.0636363636363634 3.4 \\
0 3.33636363636364 \\
0.0363636363636366 3.3 \\
0.0363636363636366 3.2 \\
0.0363636363636366 3.1 \\
0 3.06363636363636 \\
-0.0636363636363634 3 \\
0 2.93636363636364 \\
0.0363636363636366 2.9 \\
0 2.86363636363636 \\
-0.0636363636363634 2.8 \\
-0.0636363636363634 2.7 \\
-0.0636363636363634 2.6 \\
-0.0636363636363634 2.5 \\
-0.0636363636363634 2.4 \\
-0.0636363636363634 2.3 \\
0 2.23636363636364 \\
0.0363636363636366 2.2 \\
0.0363636363636366 2.1 \\
0 2.06363636363636 \\
-0.0636363636363634 2 \\
-0.0636363636363634 1.9 \\
-0.0636363636363634 1.8 \\
0 1.73636363636364 \\
0.0363636363636366 1.7 \\
0.0363636363636366 1.6 \\
0 1.56363636363636 \\
-0.0636363636363634 1.5 \\
-0.0636363636363634 1.4 \\
0 1.33636363636364 \\
0.0363636363636366 1.3 \\
0 1.26363636363636 \\
-0.0636363636363634 1.2 \\
-0.0636363636363634 1.1 \\
-0.0636363636363634 1 \\
-0.0636363636363634 0.899999999999999 \\
0 0.836363636363636 \\
0.0363636363636366 0.8 \\
0 0.763636363636364 \\
-0.0636363636363634 0.699999999999999 \\
-0.0636363636363634 0.6 \\
-0.0636363636363634 0.5 \\
-0.0636363636363634 0.399999999999999 \\
-0.0636363636363634 0.3 \\
-0.0636363636363634 0.199999999999999 \\
-0.0636363636363634 0.0999999999999996 \\
-0.0636363636363634 0 \\
-0.0636363636363634 -0.0999999999999996 \\
0 -0.163636363636363 \\
0.0363636363636366 -0.199999999999999 \\
0 -0.236363636363636 \\
-0.0636363636363634 -0.3 \\
-0.0636363636363634 -0.399999999999999 \\
-0.0636363636363634 -0.5 \\
-0.0636363636363634 -0.6 \\
0 -0.663636363636363 \\
0.0363636363636366 -0.699999999999999 \\
0 -0.736363636363636 \\
-0.0636363636363634 -0.8 \\
-0.0636363636363634 -0.899999999999999 \\
-0.0636363636363634 -1 \\
-0.0636363636363634 -1.1 \\
0 -1.16363636363636 \\
0.0363636363636366 -1.2 \\
0 -1.23636363636364 \\
-0.0636363636363634 -1.3 \\
-0.0636363636363634 -1.4 \\
-0.0636363636363634 -1.5 \\
-0.0636363636363634 -1.6 \\
0 -1.66363636363636 \\
0.0363636363636366 -1.7 \\
0 -1.73636363636364 \\
-0.0636363636363634 -1.8 \\
-0.0636363636363634 -1.9 \\
-0.0636363636363634 -2 \\
0 -2.06363636363636 \\
0.0363636363636366 -2.1 \\
0 -2.13636363636364 \\
-0.0636363636363634 -2.2 \\
-0.0636363636363634 -2.3 \\
-0.0636363636363634 -2.4 \\
-0.0636363636363634 -2.5 \\
-0.0636363636363634 -2.6 \\
-0.0636363636363634 -2.7 \\
-0.0636363636363634 -2.8 \\
-0.0636363636363634 -2.9 \\
-0.0636363636363634 -3 \\
0 -3.06363636363636 \\
0.0363636363636366 -3.1 \\
0 -3.13636363636364 \\
-0.0636363636363634 -3.2 \\
-0.0636363636363634 -3.3 \\
-0.0636363636363634 -3.4 \\
-0.0636363636363634 -3.5 \\
-0.0636363636363634 -3.6 \\
-0.0636363636363634 -3.7 \\
-0.0636363636363634 -3.8 \\
0 -3.86363636363636 \\
0.0363636363636366 -3.9 \\
0 -3.93636363636364 \\
-0.0636363636363634 -4 \\
-0.0636363636363634 -4.1 \\
0 -4.16363636363636 \\
0.0363636363636366 -4.2 \\
0 -4.23636363636364 \\
-0.0636363636363634 -4.3 \\
0 -4.36363636363636 \\
0.0363636363636366 -4.4 \\
0.0363636363636366 -4.5 \\
0.0363636363636366 -4.6 \\
0.0363636363636366 -4.7 \\
0 -4.73636363636364 \\
-0.0636363636363634 -4.8 \\
0 -4.86363636363636 \\
0.0363636363636366 -4.9 \\
0 -4.93636363636364 \\
-0.0636363636363634 -5 \\
-0.0636363636363634 -5.1 \\
-0.0636363636363634 -5.2 \\
-0.0636363636363634 -5.3 \\
-0.0636363636363634 -5.4 \\
0 -5.46363636363636 \\
0.0363636363636366 -5.5 \\
0.0363636363636366 -5.6 \\
0 -5.63636363636364 \\
-0.0636363636363634 -5.7 \\
-0.0636363636363634 -5.8 \\
-0.0636363636363634 -5.9 \\
0 -5.96363636363636 \\
0.0363636363636366 -6 \\
NaN NaN \\
};


\addplot[dashed,line width=2.0pt,draw=green,forget plot]
table[row sep=crcr]{
x y\\
0.0272727272727273 6 \\
0 5.97272727272727 \\
-0.0727272727272727 5.9 \\
-0.0727272727272727 5.8 \\
-0.0727272727272727 5.7 \\
0 5.62727272727273 \\
0.0272727272727273 5.6 \\
0.0272727272727273 5.5 \\
0 5.47272727272727 \\
-0.0727272727272727 5.4 \\
-0.0727272727272727 5.3 \\
0 5.22727272727273 \\
0.0272727272727273 5.2 \\
0 5.17272727272727 \\
-0.0727272727272727 5.1 \\
-0.0727272727272727 5 \\
0 4.92727272727273 \\
0.0272727272727273 4.9 \\
0 4.87272727272727 \\
-0.0727272727272727 4.8 \\
-0.0727272727272727 4.7 \\
-0.0727272727272727 4.6 \\
0 4.52727272727273 \\
0.0272727272727273 4.5 \\
0.0272727272727273 4.4 \\
0.0272727272727273 4.3 \\
0.0272727272727273 4.2 \\
0.0272727272727273 4.1 \\
0 4.07272727272727 \\
-0.0727272727272727 4 \\
-0.0727272727272727 3.9 \\
-0.0727272727272727 3.8 \\
-0.0727272727272727 3.7 \\
0 3.62727272727273 \\
0.0272727272727273 3.6 \\
0 3.57272727272727 \\
-0.0727272727272727 3.5 \\
-0.0727272727272727 3.4 \\
0 3.32727272727273 \\
0.0272727272727273 3.3 \\
0.0272727272727273 3.2 \\
0.0272727272727273 3.1 \\
0 3.07272727272727 \\
-0.0727272727272727 3 \\
0 2.92727272727273 \\
0.0272727272727273 2.9 \\
0 2.87272727272727 \\
-0.0727272727272727 2.8 \\
-0.0727272727272727 2.7 \\
-0.0727272727272727 2.6 \\
-0.0727272727272727 2.5 \\
-0.0727272727272727 2.4 \\
-0.0727272727272727 2.3 \\
0 2.22727272727273 \\
0.0272727272727273 2.2 \\
0.0272727272727273 2.1 \\
0 2.07272727272727 \\
-0.0727272727272727 2 \\
-0.0727272727272727 1.9 \\
-0.0727272727272727 1.8 \\
0 1.72727272727273 \\
0.0272727272727273 1.7 \\
0.0272727272727273 1.6 \\
0 1.57272727272727 \\
-0.0727272727272727 1.5 \\
-0.0727272727272727 1.4 \\
0 1.32727272727273 \\
0.0272727272727273 1.3 \\
0 1.27272727272727 \\
-0.0727272727272727 1.2 \\
-0.0727272727272727 1.1 \\
-0.0727272727272727 1 \\
-0.0727272727272727 0.899999999999999 \\
0 0.827272727272726 \\
0.0272727272727273 0.8 \\
0 0.772727272727273 \\
-0.0727272727272727 0.699999999999999 \\
-0.0727272727272727 0.6 \\
-0.0727272727272727 0.5 \\
-0.0727272727272727 0.399999999999999 \\
-0.0727272727272727 0.3 \\
-0.0727272727272727 0.199999999999999 \\
-0.0727272727272727 0.0999999999999996 \\
-0.0727272727272727 0 \\
-0.0727272727272727 -0.0999999999999996 \\
0 -0.172727272727272 \\
0.0272727272727273 -0.199999999999999 \\
0 -0.227272727272727 \\
-0.0727272727272727 -0.3 \\
-0.0727272727272727 -0.399999999999999 \\
-0.0727272727272727 -0.5 \\
-0.0727272727272727 -0.6 \\
0 -0.672727272727272 \\
0.0272727272727273 -0.699999999999999 \\
0 -0.727272727272727 \\
-0.0727272727272727 -0.8 \\
-0.0727272727272727 -0.899999999999999 \\
-0.0727272727272727 -1 \\
-0.0727272727272727 -1.1 \\
0 -1.17272727272727 \\
0.0272727272727273 -1.2 \\
0 -1.22727272727273 \\
-0.0727272727272727 -1.3 \\
-0.0727272727272727 -1.4 \\
-0.0727272727272727 -1.5 \\
-0.0727272727272727 -1.6 \\
0 -1.67272727272727 \\
0.0272727272727273 -1.7 \\
0 -1.72727272727273 \\
-0.0727272727272727 -1.8 \\
-0.0727272727272727 -1.9 \\
-0.0727272727272727 -2 \\
0 -2.07272727272727 \\
0.0272727272727273 -2.1 \\
0 -2.12727272727273 \\
-0.0727272727272727 -2.2 \\
-0.0727272727272727 -2.3 \\
-0.0727272727272727 -2.4 \\
-0.0727272727272727 -2.5 \\
-0.0727272727272727 -2.6 \\
-0.0727272727272727 -2.7 \\
-0.0727272727272727 -2.8 \\
-0.0727272727272727 -2.9 \\
-0.0727272727272727 -3 \\
0 -3.07272727272727 \\
0.0272727272727273 -3.1 \\
0 -3.12727272727273 \\
-0.0727272727272727 -3.2 \\
-0.0727272727272727 -3.3 \\
-0.0727272727272727 -3.4 \\
-0.0727272727272727 -3.5 \\
-0.0727272727272727 -3.6 \\
-0.0727272727272727 -3.7 \\
-0.0727272727272727 -3.8 \\
0 -3.87272727272727 \\
0.0272727272727273 -3.9 \\
0 -3.92727272727273 \\
-0.0727272727272727 -4 \\
-0.0727272727272727 -4.1 \\
0 -4.17272727272727 \\
0.0272727272727273 -4.2 \\
0 -4.22727272727273 \\
-0.0727272727272727 -4.3 \\
0 -4.37272727272727 \\
0.0272727272727273 -4.4 \\
0.0272727272727273 -4.5 \\
0.0272727272727273 -4.6 \\
0.0272727272727273 -4.7 \\
0 -4.72727272727273 \\
-0.0727272727272727 -4.8 \\
0 -4.87272727272727 \\
0.0272727272727273 -4.9 \\
0 -4.92727272727273 \\
-0.0727272727272727 -5 \\
-0.0727272727272727 -5.1 \\
-0.0727272727272727 -5.2 \\
-0.0727272727272727 -5.3 \\
-0.0727272727272727 -5.4 \\
0 -5.47272727272727 \\
0.0272727272727273 -5.5 \\
0.0272727272727273 -5.6 \\
0 -5.62727272727273 \\
-0.0727272727272727 -5.7 \\
-0.0727272727272727 -5.8 \\
-0.0727272727272727 -5.9 \\
0 -5.97272727272727 \\
0.0272727272727273 -6 \\
NaN NaN \\
};


\addplot[dashed,line width=2.0pt,draw=green,forget plot]
table[row sep=crcr]{
x y\\
0.018181818181818 6 \\
0 5.98181818181818 \\
-0.0818181818181821 5.9 \\
-0.0818181818181821 5.8 \\
-0.0818181818181821 5.7 \\
0 5.61818181818182 \\
0.018181818181818 5.6 \\
0.018181818181818 5.5 \\
0 5.48181818181818 \\
-0.0818181818181821 5.4 \\
-0.0818181818181821 5.3 \\
0 5.21818181818182 \\
0.018181818181818 5.2 \\
0 5.18181818181818 \\
-0.0818181818181821 5.1 \\
-0.0818181818181821 5 \\
0 4.91818181818182 \\
0.018181818181818 4.9 \\
0 4.88181818181818 \\
-0.0818181818181821 4.8 \\
-0.0818181818181821 4.7 \\
-0.0818181818181821 4.6 \\
0 4.51818181818182 \\
0.018181818181818 4.5 \\
0.018181818181818 4.4 \\
0.018181818181818 4.3 \\
0.018181818181818 4.2 \\
0.018181818181818 4.1 \\
0 4.08181818181818 \\
-0.0818181818181821 4 \\
-0.0818181818181821 3.9 \\
-0.0818181818181821 3.8 \\
-0.0818181818181821 3.7 \\
0 3.61818181818182 \\
0.018181818181818 3.6 \\
0 3.58181818181818 \\
-0.0818181818181821 3.5 \\
-0.0818181818181821 3.4 \\
0 3.31818181818182 \\
0.018181818181818 3.3 \\
0.018181818181818 3.2 \\
0.018181818181818 3.1 \\
0 3.08181818181818 \\
-0.0818181818181821 3 \\
0 2.91818181818182 \\
0.018181818181818 2.9 \\
0 2.88181818181818 \\
-0.0818181818181821 2.8 \\
-0.0818181818181821 2.7 \\
-0.0818181818181821 2.6 \\
-0.0818181818181821 2.5 \\
-0.0818181818181821 2.4 \\
-0.0818181818181821 2.3 \\
0 2.21818181818182 \\
0.018181818181818 2.2 \\
0.018181818181818 2.1 \\
0 2.08181818181818 \\
-0.0818181818181821 2 \\
-0.0818181818181821 1.9 \\
-0.0818181818181821 1.8 \\
0 1.71818181818182 \\
0.018181818181818 1.7 \\
0.018181818181818 1.6 \\
0 1.58181818181818 \\
-0.0818181818181821 1.5 \\
-0.0818181818181821 1.4 \\
0 1.31818181818182 \\
0.018181818181818 1.3 \\
0 1.28181818181818 \\
-0.0818181818181821 1.2 \\
-0.0818181818181821 1.1 \\
-0.0818181818181821 1 \\
-0.0818181818181821 0.899999999999999 \\
0 0.818181818181818 \\
0.018181818181818 0.8 \\
0 0.781818181818181 \\
-0.0818181818181821 0.699999999999999 \\
-0.0818181818181821 0.6 \\
-0.0818181818181821 0.5 \\
-0.0818181818181821 0.399999999999999 \\
-0.0818181818181821 0.3 \\
-0.0818181818181821 0.199999999999999 \\
-0.0818181818181821 0.0999999999999996 \\
-0.0818181818181821 0 \\
-0.0818181818181821 -0.0999999999999996 \\
0 -0.181818181818181 \\
0.018181818181818 -0.199999999999999 \\
0 -0.218181818181817 \\
-0.0818181818181821 -0.3 \\
-0.0818181818181821 -0.399999999999999 \\
-0.0818181818181821 -0.5 \\
-0.0818181818181821 -0.6 \\
0 -0.681818181818181 \\
0.018181818181818 -0.699999999999999 \\
0 -0.718181818181817 \\
-0.0818181818181821 -0.8 \\
-0.0818181818181821 -0.899999999999999 \\
-0.0818181818181821 -1 \\
-0.0818181818181821 -1.1 \\
0 -1.18181818181818 \\
0.018181818181818 -1.2 \\
0 -1.21818181818182 \\
-0.0818181818181821 -1.3 \\
-0.0818181818181821 -1.4 \\
-0.0818181818181821 -1.5 \\
-0.0818181818181821 -1.6 \\
0 -1.68181818181818 \\
0.018181818181818 -1.7 \\
0 -1.71818181818182 \\
-0.0818181818181821 -1.8 \\
-0.0818181818181821 -1.9 \\
-0.0818181818181821 -2 \\
0 -2.08181818181818 \\
0.018181818181818 -2.1 \\
0 -2.11818181818182 \\
-0.0818181818181821 -2.2 \\
-0.0818181818181821 -2.3 \\
-0.0818181818181821 -2.4 \\
-0.0818181818181821 -2.5 \\
-0.0818181818181821 -2.6 \\
-0.0818181818181821 -2.7 \\
-0.0818181818181821 -2.8 \\
-0.0818181818181821 -2.9 \\
-0.0818181818181821 -3 \\
0 -3.08181818181818 \\
0.018181818181818 -3.1 \\
0 -3.11818181818182 \\
-0.0818181818181821 -3.2 \\
-0.0818181818181821 -3.3 \\
-0.0818181818181821 -3.4 \\
-0.0818181818181821 -3.5 \\
-0.0818181818181821 -3.6 \\
-0.0818181818181821 -3.7 \\
-0.0818181818181821 -3.8 \\
0 -3.88181818181818 \\
0.018181818181818 -3.9 \\
0 -3.91818181818182 \\
-0.0818181818181821 -4 \\
-0.0818181818181821 -4.1 \\
0 -4.18181818181818 \\
0.018181818181818 -4.2 \\
0 -4.21818181818182 \\
-0.0818181818181821 -4.3 \\
0 -4.38181818181818 \\
0.018181818181818 -4.4 \\
0.018181818181818 -4.5 \\
0.018181818181818 -4.6 \\
0.018181818181818 -4.7 \\
0 -4.71818181818182 \\
-0.0818181818181821 -4.8 \\
0 -4.88181818181818 \\
0.018181818181818 -4.9 \\
0 -4.91818181818182 \\
-0.0818181818181821 -5 \\
-0.0818181818181821 -5.1 \\
-0.0818181818181821 -5.2 \\
-0.0818181818181821 -5.3 \\
-0.0818181818181821 -5.4 \\
0 -5.48181818181818 \\
0.018181818181818 -5.5 \\
0.018181818181818 -5.6 \\
0 -5.61818181818182 \\
-0.0818181818181821 -5.7 \\
-0.0818181818181821 -5.8 \\
-0.0818181818181821 -5.9 \\
0 -5.98181818181818 \\
0.018181818181818 -6 \\
NaN NaN \\
};


\addplot[dashed,line width=2.0pt,draw=green,forget plot]
table[row sep=crcr]{
x y\\
0.00909090909090934 6 \\
0 5.99090909090909 \\
-0.0909090909090907 5.9 \\
-0.0909090909090907 5.8 \\
-0.0909090909090907 5.7 \\
0 5.60909090909091 \\
0.00909090909090934 5.6 \\
0.00909090909090934 5.5 \\
0 5.49090909090909 \\
-0.0909090909090907 5.4 \\
-0.0909090909090907 5.3 \\
0 5.20909090909091 \\
0.00909090909090934 5.2 \\
0 5.19090909090909 \\
-0.0909090909090907 5.1 \\
-0.0909090909090907 5 \\
0 4.90909090909091 \\
0.00909090909090934 4.9 \\
0 4.89090909090909 \\
-0.0909090909090907 4.8 \\
-0.0909090909090907 4.7 \\
-0.0909090909090907 4.6 \\
0 4.50909090909091 \\
0.00909090909090934 4.5 \\
0.00909090909090934 4.4 \\
0.00909090909090934 4.3 \\
0.00909090909090934 4.2 \\
0.00909090909090934 4.1 \\
0 4.09090909090909 \\
-0.0909090909090907 4 \\
-0.0909090909090907 3.9 \\
-0.0909090909090907 3.8 \\
-0.0909090909090907 3.7 \\
0 3.60909090909091 \\
0.00909090909090934 3.6 \\
0 3.59090909090909 \\
-0.0909090909090907 3.5 \\
-0.0909090909090907 3.4 \\
0 3.30909090909091 \\
0.00909090909090934 3.3 \\
0.00909090909090934 3.2 \\
0.00909090909090934 3.1 \\
0 3.09090909090909 \\
-0.0909090909090907 3 \\
0 2.90909090909091 \\
0.00909090909090934 2.9 \\
0 2.89090909090909 \\
-0.0909090909090907 2.8 \\
-0.0909090909090907 2.7 \\
-0.0909090909090907 2.6 \\
-0.0909090909090907 2.5 \\
-0.0909090909090907 2.4 \\
-0.0909090909090907 2.3 \\
0 2.20909090909091 \\
0.00909090909090934 2.2 \\
0.00909090909090934 2.1 \\
0 2.09090909090909 \\
-0.0909090909090907 2 \\
-0.0909090909090907 1.9 \\
-0.0909090909090907 1.8 \\
0 1.70909090909091 \\
0.00909090909090934 1.7 \\
0.00909090909090934 1.6 \\
0 1.59090909090909 \\
-0.0909090909090907 1.5 \\
-0.0909090909090907 1.4 \\
0 1.30909090909091 \\
0.00909090909090934 1.3 \\
0 1.29090909090909 \\
-0.0909090909090907 1.2 \\
-0.0909090909090907 1.1 \\
-0.0909090909090907 1 \\
-0.0909090909090907 0.899999999999999 \\
0 0.809090909090909 \\
0.00909090909090934 0.8 \\
0 0.79090909090909 \\
-0.0909090909090907 0.699999999999999 \\
-0.0909090909090907 0.6 \\
-0.0909090909090907 0.5 \\
-0.0909090909090907 0.399999999999999 \\
-0.0909090909090907 0.3 \\
-0.0909090909090907 0.199999999999999 \\
-0.0909090909090907 0.0999999999999996 \\
-0.0909090909090907 0 \\
-0.0909090909090907 -0.0999999999999996 \\
0 -0.19090909090909 \\
0.00909090909090934 -0.199999999999999 \\
0 -0.209090909090909 \\
-0.0909090909090907 -0.3 \\
-0.0909090909090907 -0.399999999999999 \\
-0.0909090909090907 -0.5 \\
-0.0909090909090907 -0.6 \\
0 -0.69090909090909 \\
0.00909090909090934 -0.699999999999999 \\
0 -0.709090909090909 \\
-0.0909090909090907 -0.8 \\
-0.0909090909090907 -0.899999999999999 \\
-0.0909090909090907 -1 \\
-0.0909090909090907 -1.1 \\
0 -1.19090909090909 \\
0.00909090909090934 -1.2 \\
0 -1.20909090909091 \\
-0.0909090909090907 -1.3 \\
-0.0909090909090907 -1.4 \\
-0.0909090909090907 -1.5 \\
-0.0909090909090907 -1.6 \\
0 -1.69090909090909 \\
0.00909090909090934 -1.7 \\
0 -1.70909090909091 \\
-0.0909090909090907 -1.8 \\
-0.0909090909090907 -1.9 \\
-0.0909090909090907 -2 \\
0 -2.09090909090909 \\
0.00909090909090934 -2.1 \\
0 -2.10909090909091 \\
-0.0909090909090907 -2.2 \\
-0.0909090909090907 -2.3 \\
-0.0909090909090907 -2.4 \\
-0.0909090909090907 -2.5 \\
-0.0909090909090907 -2.6 \\
-0.0909090909090907 -2.7 \\
-0.0909090909090907 -2.8 \\
-0.0909090909090907 -2.9 \\
-0.0909090909090907 -3 \\
0 -3.09090909090909 \\
0.00909090909090934 -3.1 \\
0 -3.10909090909091 \\
-0.0909090909090907 -3.2 \\
-0.0909090909090907 -3.3 \\
-0.0909090909090907 -3.4 \\
-0.0909090909090907 -3.5 \\
-0.0909090909090907 -3.6 \\
-0.0909090909090907 -3.7 \\
-0.0909090909090907 -3.8 \\
0 -3.89090909090909 \\
0.00909090909090934 -3.9 \\
0 -3.90909090909091 \\
-0.0909090909090907 -4 \\
-0.0909090909090907 -4.1 \\
0 -4.19090909090909 \\
0.00909090909090934 -4.2 \\
0 -4.20909090909091 \\
-0.0909090909090907 -4.3 \\
0 -4.39090909090909 \\
0.00909090909090934 -4.4 \\
0.00909090909090934 -4.5 \\
0.00909090909090934 -4.6 \\
0.00909090909090934 -4.7 \\
0 -4.70909090909091 \\
-0.0909090909090907 -4.8 \\
0 -4.89090909090909 \\
0.00909090909090934 -4.9 \\
0 -4.90909090909091 \\
-0.0909090909090907 -5 \\
-0.0909090909090907 -5.1 \\
-0.0909090909090907 -5.2 \\
-0.0909090909090907 -5.3 \\
-0.0909090909090907 -5.4 \\
0 -5.49090909090909 \\
0.00909090909090934 -5.5 \\
0.00909090909090934 -5.6 \\
0 -5.60909090909091 \\
-0.0909090909090907 -5.7 \\
-0.0909090909090907 -5.8 \\
-0.0909090909090907 -5.9 \\
0 -5.99090909090909 \\
0.00909090909090934 -6 \\
NaN NaN \\
};


\addplot[dashed,line width=2.0pt,draw=green,forget plot]
table[row sep=crcr]{
x y\\
0 6 \\
0 6 \\
-0.1 5.9 \\
-0.1 5.8 \\
-0.1 5.7 \\
0 5.6 \\
0 5.6 \\
0 5.5 \\
0 5.5 \\
-0.1 5.4 \\
-0.1 5.3 \\
0 5.2 \\
0 5.2 \\
0 5.2 \\
-0.1 5.1 \\
-0.1 5 \\
0 4.9 \\
0 4.9 \\
0 4.9 \\
-0.1 4.8 \\
-0.1 4.7 \\
-0.1 4.6 \\
0 4.5 \\
0 4.5 \\
0 4.4 \\
0 4.3 \\
0 4.2 \\
0 4.1 \\
0 4.1 \\
-0.1 4 \\
-0.1 3.9 \\
-0.1 3.8 \\
-0.1 3.7 \\
0 3.6 \\
0 3.6 \\
0 3.6 \\
-0.1 3.5 \\
-0.1 3.4 \\
0 3.3 \\
0 3.3 \\
0 3.2 \\
0 3.1 \\
0 3.1 \\
-0.1 3 \\
0 2.9 \\
0 2.9 \\
0 2.9 \\
-0.1 2.8 \\
-0.1 2.7 \\
-0.1 2.6 \\
-0.1 2.5 \\
-0.1 2.4 \\
-0.1 2.3 \\
0 2.2 \\
0 2.2 \\
0 2.1 \\
0 2.1 \\
-0.1 2 \\
-0.1 1.9 \\
-0.1 1.8 \\
0 1.7 \\
0 1.7 \\
0 1.6 \\
0 1.6 \\
-0.1 1.5 \\
-0.1 1.4 \\
0 1.3 \\
0 1.3 \\
0 1.3 \\
-0.1 1.2 \\
-0.1 1.1 \\
-0.1 1 \\
-0.1 0.899999999999999 \\
0 0.8 \\
0 0.8 \\
0 0.8 \\
-0.1 0.699999999999999 \\
-0.1 0.6 \\
-0.1 0.5 \\
-0.1 0.399999999999999 \\
-0.1 0.3 \\
-0.1 0.199999999999999 \\
-0.1 0.0999999999999996 \\
-0.1 0 \\
-0.1 -0.0999999999999996 \\
0 -0.199999999999999 \\
0 -0.199999999999999 \\
0 -0.199999999999999 \\
-0.1 -0.3 \\
-0.1 -0.399999999999999 \\
-0.1 -0.5 \\
-0.1 -0.6 \\
0 -0.699999999999999 \\
0 -0.699999999999999 \\
0 -0.699999999999999 \\
-0.1 -0.8 \\
-0.1 -0.899999999999999 \\
-0.1 -1 \\
-0.1 -1.1 \\
0 -1.2 \\
0 -1.2 \\
0 -1.2 \\
-0.1 -1.3 \\
-0.1 -1.4 \\
-0.1 -1.5 \\
-0.1 -1.6 \\
0 -1.7 \\
0 -1.7 \\
0 -1.7 \\
-0.1 -1.8 \\
-0.1 -1.9 \\
-0.1 -2 \\
0 -2.1 \\
0 -2.1 \\
0 -2.1 \\
-0.1 -2.2 \\
-0.1 -2.3 \\
-0.1 -2.4 \\
-0.1 -2.5 \\
-0.1 -2.6 \\
-0.1 -2.7 \\
-0.1 -2.8 \\
-0.1 -2.9 \\
-0.1 -3 \\
0 -3.1 \\
0 -3.1 \\
0 -3.1 \\
-0.1 -3.2 \\
-0.1 -3.3 \\
-0.1 -3.4 \\
-0.1 -3.5 \\
-0.1 -3.6 \\
-0.1 -3.7 \\
-0.1 -3.8 \\
0 -3.9 \\
0 -3.9 \\
0 -3.9 \\
-0.1 -4 \\
-0.1 -4.1 \\
0 -4.2 \\
0 -4.2 \\
0 -4.2 \\
-0.1 -4.3 \\
0 -4.4 \\
0 -4.4 \\
0 -4.5 \\
0 -4.6 \\
0 -4.7 \\
0 -4.7 \\
-0.1 -4.8 \\
0 -4.9 \\
0 -4.9 \\
0 -4.9 \\
-0.1 -5 \\
-0.1 -5.1 \\
-0.1 -5.2 \\
-0.1 -5.3 \\
-0.1 -5.4 \\
0 -5.5 \\
0 -5.5 \\
0 -5.6 \\
0 -5.6 \\
-0.1 -5.7 \\
-0.1 -5.8 \\
-0.1 -5.9 \\
0 -6 \\
0 -6 \\
NaN NaN \\
};


\addplot[dashed,line width=2.0pt,draw=blue,forget plot]
table[row sep=crcr]{
x y\\
-1.9 1.09999999 \\
-1.89999999 1.1 \\
-1.89999999 1.2 \\
-1.8 1.29999999 \\
-1.79999999 1.3 \\
-1.7 1.39999999 \\
-1.69999999 1.4 \\
-1.6 1.49999999 \\
-1.59999999 1.5 \\
-1.5 1.59999999 \\
-1.49999999 1.6 \\
-1.4 1.69999999 \\
-1.3 1.69999999 \\
-1.29999999 1.7 \\
-1.2 1.79999999 \\
-1.1 1.79999999 \\
-1 1.79999999 \\
-0.9 1.79999999 \\
-0.8 1.79999999 \\
-0.70000001 1.7 \\
-0.7 1.69999999 \\
-0.6 1.69999999 \\
-0.50000001 1.6 \\
-0.5 1.59999999 \\
-0.40000001 1.5 \\
-0.4 1.49999999 \\
-0.30000001 1.4 \\
-0.3 1.39999999 \\
-0.20000001 1.3 \\
-0.2 1.29999999 \\
-0.10000001 1.2 \\
-0.10000001 1.1 \\
-0.1 1.09999999 \\
-1.0000000116861e-008 1 \\
-1.0000000116861e-008 0.899999999999999 \\
-1.0000000116861e-008 0.8 \\
-1.0000000116861e-008 0.699999999999999 \\
-1.0000000116861e-008 0.6 \\
-1.0000000116861e-008 0.5 \\
-1.0000000116861e-008 0.399999999999999 \\
-1.0000000116861e-008 0.3 \\
-1.0000000116861e-008 0.199999999999999 \\
-1.0000000116861e-008 0.0999999999999996 \\
-1.0000000116861e-008 0 \\
-1.0000000116861e-008 -0.0999999999999996 \\
-1.0000000116861e-008 -0.199999999999999 \\
-1.0000000116861e-008 -0.3 \\
-1.0000000116861e-008 -0.399999999999999 \\
-1.0000000116861e-008 -0.5 \\
-1.0000000116861e-008 -0.6 \\
-1.0000000116861e-008 -0.699999999999999 \\
-1.0000000116861e-008 -0.8 \\
-1.0000000116861e-008 -0.899999999999999 \\
-1.0000000116861e-008 -1 \\
-0.1 -1.09999999 \\
-0.10000001 -1.1 \\
-0.10000001 -1.2 \\
-0.2 -1.29999999 \\
-0.20000001 -1.3 \\
-0.3 -1.39999999 \\
-0.30000001 -1.4 \\
-0.4 -1.49999999 \\
-0.40000001 -1.5 \\
-0.5 -1.59999999 \\
-0.50000001 -1.6 \\
-0.6 -1.69999999 \\
-0.7 -1.69999999 \\
-0.70000001 -1.7 \\
-0.8 -1.79999999 \\
-0.9 -1.79999999 \\
-1 -1.79999999 \\
-1.1 -1.79999999 \\
-1.2 -1.79999999 \\
-1.29999999 -1.7 \\
-1.3 -1.69999999 \\
-1.4 -1.69999999 \\
-1.49999999 -1.6 \\
-1.5 -1.59999999 \\
-1.59999999 -1.5 \\
-1.6 -1.49999999 \\
-1.69999999 -1.4 \\
-1.7 -1.39999999 \\
-1.79999999 -1.3 \\
-1.8 -1.29999999 \\
-1.89999999 -1.2 \\
-1.89999999 -1.1 \\
-1.9 -1.09999999 \\
-1.99999999 -1 \\
-1.99999999 -0.899999999999999 \\
-1.99999999 -0.8 \\
-1.99999999 -0.699999999999999 \\
-1.99999999 -0.6 \\
-1.99999999 -0.5 \\
-1.99999999 -0.399999999999999 \\
-1.99999999 -0.3 \\
-1.99999999 -0.199999999999999 \\
-1.99999999 -0.0999999999999996 \\
-1.99999999 0 \\
-1.99999999 0.0999999999999996 \\
-1.99999999 0.199999999999999 \\
-1.99999999 0.3 \\
-1.99999999 0.399999999999999 \\
-1.99999999 0.5 \\
-1.99999999 0.6 \\
-1.99999999 0.699999999999999 \\
-1.99999999 0.8 \\
-1.99999999 0.899999999999999 \\
-1.99999999 1 \\
-1.9 1.09999999 \\
NaN NaN \\
};


\addplot[dashed,line width=2.0pt,draw=blue,forget plot]
table[row sep=crcr]{
x y\\
-1.9 1.09090909090909 \\
-1.89090909090909 1.1 \\
-1.89090909090909 1.2 \\
-1.8 1.29090909090909 \\
-1.79090909090909 1.3 \\
-1.7 1.39090909090909 \\
-1.69090909090909 1.4 \\
-1.6 1.49090909090909 \\
-1.59090909090909 1.5 \\
-1.5 1.59090909090909 \\
-1.49090909090909 1.6 \\
-1.4 1.69090909090909 \\
-1.3 1.69090909090909 \\
-1.29090909090909 1.7 \\
-1.2 1.79090909090909 \\
-1.1 1.79090909090909 \\
-1 1.79090909090909 \\
-0.9 1.79090909090909 \\
-0.8 1.79090909090909 \\
-0.70909090909091 1.7 \\
-0.7 1.69090909090909 \\
-0.6 1.69090909090909 \\
-0.509090909090909 1.6 \\
-0.5 1.59090909090909 \\
-0.40909090909091 1.5 \\
-0.4 1.49090909090909 \\
-0.30909090909091 1.4 \\
-0.3 1.39090909090909 \\
-0.20909090909091 1.3 \\
-0.2 1.29090909090909 \\
-0.109090909090909 1.2 \\
-0.109090909090909 1.1 \\
-0.1 1.09090909090909 \\
-0.00909090909090936 1 \\
-0.00909090909090936 0.899999999999999 \\
-0.00909090909090936 0.8 \\
-0.00909090909090936 0.699999999999999 \\
-0.00909090909090936 0.6 \\
-0.00909090909090936 0.5 \\
-0.00909090909090936 0.399999999999999 \\
-0.00909090909090936 0.3 \\
-0.00909090909090936 0.199999999999999 \\
-0.00909090909090936 0.0999999999999996 \\
-0.00909090909090936 0 \\
-0.00909090909090936 -0.0999999999999996 \\
-0.00909090909090936 -0.199999999999999 \\
-0.00909090909090936 -0.3 \\
-0.00909090909090936 -0.399999999999999 \\
-0.00909090909090936 -0.5 \\
-0.00909090909090936 -0.6 \\
-0.00909090909090936 -0.699999999999999 \\
-0.00909090909090936 -0.8 \\
-0.00909090909090936 -0.899999999999999 \\
-0.00909090909090936 -1 \\
-0.1 -1.09090909090909 \\
-0.109090909090909 -1.1 \\
-0.109090909090909 -1.2 \\
-0.2 -1.29090909090909 \\
-0.20909090909091 -1.3 \\
-0.3 -1.39090909090909 \\
-0.30909090909091 -1.4 \\
-0.4 -1.49090909090909 \\
-0.40909090909091 -1.5 \\
-0.5 -1.59090909090909 \\
-0.509090909090909 -1.6 \\
-0.6 -1.69090909090909 \\
-0.7 -1.69090909090909 \\
-0.70909090909091 -1.7 \\
-0.8 -1.79090909090909 \\
-0.9 -1.79090909090909 \\
-1 -1.79090909090909 \\
-1.1 -1.79090909090909 \\
-1.2 -1.79090909090909 \\
-1.29090909090909 -1.7 \\
-1.3 -1.69090909090909 \\
-1.4 -1.69090909090909 \\
-1.49090909090909 -1.6 \\
-1.5 -1.59090909090909 \\
-1.59090909090909 -1.5 \\
-1.6 -1.49090909090909 \\
-1.69090909090909 -1.4 \\
-1.7 -1.39090909090909 \\
-1.79090909090909 -1.3 \\
-1.8 -1.29090909090909 \\
-1.89090909090909 -1.2 \\
-1.89090909090909 -1.1 \\
-1.9 -1.09090909090909 \\
-1.99090909090909 -1 \\
-1.99090909090909 -0.899999999999999 \\
-1.99090909090909 -0.8 \\
-1.99090909090909 -0.699999999999999 \\
-1.99090909090909 -0.6 \\
-1.99090909090909 -0.5 \\
-1.99090909090909 -0.399999999999999 \\
-1.99090909090909 -0.3 \\
-1.99090909090909 -0.199999999999999 \\
-1.99090909090909 -0.0999999999999996 \\
-1.99090909090909 0 \\
-1.99090909090909 0.0999999999999996 \\
-1.99090909090909 0.199999999999999 \\
-1.99090909090909 0.3 \\
-1.99090909090909 0.399999999999999 \\
-1.99090909090909 0.5 \\
-1.99090909090909 0.6 \\
-1.99090909090909 0.699999999999999 \\
-1.99090909090909 0.8 \\
-1.99090909090909 0.899999999999999 \\
-1.99090909090909 1 \\
-1.9 1.09090909090909 \\
NaN NaN \\
};


\addplot[dashed,line width=2.0pt,draw=blue,forget plot]
table[row sep=crcr]{
x y\\
-1.9 1.08181818181818 \\
-1.88181818181818 1.1 \\
-1.88181818181818 1.2 \\
-1.8 1.28181818181818 \\
-1.78181818181818 1.3 \\
-1.7 1.38181818181818 \\
-1.68181818181818 1.4 \\
-1.6 1.48181818181818 \\
-1.58181818181818 1.5 \\
-1.5 1.58181818181818 \\
-1.48181818181818 1.6 \\
-1.4 1.68181818181818 \\
-1.3 1.68181818181818 \\
-1.28181818181818 1.7 \\
-1.2 1.78181818181818 \\
-1.1 1.78181818181818 \\
-1 1.78181818181818 \\
-0.9 1.78181818181818 \\
-0.8 1.78181818181818 \\
-0.718181818181818 1.7 \\
-0.7 1.68181818181818 \\
-0.6 1.68181818181818 \\
-0.518181818181818 1.6 \\
-0.5 1.58181818181818 \\
-0.418181818181818 1.5 \\
-0.4 1.48181818181818 \\
-0.318181818181818 1.4 \\
-0.3 1.38181818181818 \\
-0.218181818181818 1.3 \\
-0.2 1.28181818181818 \\
-0.118181818181818 1.2 \\
-0.118181818181818 1.1 \\
-0.1 1.08181818181818 \\
-0.018181818181818 1 \\
-0.018181818181818 0.899999999999999 \\
-0.018181818181818 0.8 \\
-0.018181818181818 0.699999999999999 \\
-0.018181818181818 0.6 \\
-0.018181818181818 0.5 \\
-0.018181818181818 0.399999999999999 \\
-0.018181818181818 0.3 \\
-0.018181818181818 0.199999999999999 \\
-0.018181818181818 0.0999999999999996 \\
-0.018181818181818 0 \\
-0.018181818181818 -0.0999999999999996 \\
-0.018181818181818 -0.199999999999999 \\
-0.018181818181818 -0.3 \\
-0.018181818181818 -0.399999999999999 \\
-0.018181818181818 -0.5 \\
-0.018181818181818 -0.6 \\
-0.018181818181818 -0.699999999999999 \\
-0.018181818181818 -0.8 \\
-0.018181818181818 -0.899999999999999 \\
-0.018181818181818 -1 \\
-0.1 -1.08181818181818 \\
-0.118181818181818 -1.1 \\
-0.118181818181818 -1.2 \\
-0.2 -1.28181818181818 \\
-0.218181818181818 -1.3 \\
-0.3 -1.38181818181818 \\
-0.318181818181818 -1.4 \\
-0.4 -1.48181818181818 \\
-0.418181818181818 -1.5 \\
-0.5 -1.58181818181818 \\
-0.518181818181818 -1.6 \\
-0.6 -1.68181818181818 \\
-0.7 -1.68181818181818 \\
-0.718181818181818 -1.7 \\
-0.8 -1.78181818181818 \\
-0.9 -1.78181818181818 \\
-1 -1.78181818181818 \\
-1.1 -1.78181818181818 \\
-1.2 -1.78181818181818 \\
-1.28181818181818 -1.7 \\
-1.3 -1.68181818181818 \\
-1.4 -1.68181818181818 \\
-1.48181818181818 -1.6 \\
-1.5 -1.58181818181818 \\
-1.58181818181818 -1.5 \\
-1.6 -1.48181818181818 \\
-1.68181818181818 -1.4 \\
-1.7 -1.38181818181818 \\
-1.78181818181818 -1.3 \\
-1.8 -1.28181818181818 \\
-1.88181818181818 -1.2 \\
-1.88181818181818 -1.1 \\
-1.9 -1.08181818181818 \\
-1.98181818181818 -1 \\
-1.98181818181818 -0.899999999999999 \\
-1.98181818181818 -0.8 \\
-1.98181818181818 -0.699999999999999 \\
-1.98181818181818 -0.6 \\
-1.98181818181818 -0.5 \\
-1.98181818181818 -0.399999999999999 \\
-1.98181818181818 -0.3 \\
-1.98181818181818 -0.199999999999999 \\
-1.98181818181818 -0.0999999999999996 \\
-1.98181818181818 0 \\
-1.98181818181818 0.0999999999999996 \\
-1.98181818181818 0.199999999999999 \\
-1.98181818181818 0.3 \\
-1.98181818181818 0.399999999999999 \\
-1.98181818181818 0.5 \\
-1.98181818181818 0.6 \\
-1.98181818181818 0.699999999999999 \\
-1.98181818181818 0.8 \\
-1.98181818181818 0.899999999999999 \\
-1.98181818181818 1 \\
-1.9 1.08181818181818 \\
NaN NaN \\
};


\addplot[dashed,line width=2.0pt,draw=blue,forget plot]
table[row sep=crcr]{
x y\\
-1.9 1.07272727272727 \\
-1.87272727272727 1.1 \\
-1.87272727272727 1.2 \\
-1.8 1.27272727272727 \\
-1.77272727272727 1.3 \\
-1.7 1.37272727272727 \\
-1.67272727272727 1.4 \\
-1.6 1.47272727272727 \\
-1.57272727272727 1.5 \\
-1.5 1.57272727272727 \\
-1.47272727272727 1.6 \\
-1.4 1.67272727272727 \\
-1.3 1.67272727272727 \\
-1.27272727272727 1.7 \\
-1.2 1.77272727272727 \\
-1.1 1.77272727272727 \\
-1 1.77272727272727 \\
-0.9 1.77272727272727 \\
-0.8 1.77272727272727 \\
-0.727272727272728 1.7 \\
-0.7 1.67272727272727 \\
-0.6 1.67272727272727 \\
-0.527272727272727 1.6 \\
-0.5 1.57272727272727 \\
-0.427272727272728 1.5 \\
-0.4 1.47272727272727 \\
-0.327272727272728 1.4 \\
-0.3 1.37272727272727 \\
-0.227272727272728 1.3 \\
-0.2 1.27272727272727 \\
-0.127272727272727 1.2 \\
-0.127272727272727 1.1 \\
-0.1 1.07272727272727 \\
-0.0272727272727274 1 \\
-0.0272727272727274 0.899999999999999 \\
-0.0272727272727274 0.8 \\
-0.0272727272727274 0.699999999999999 \\
-0.0272727272727274 0.6 \\
-0.0272727272727274 0.5 \\
-0.0272727272727274 0.399999999999999 \\
-0.0272727272727274 0.3 \\
-0.0272727272727274 0.199999999999999 \\
-0.0272727272727274 0.0999999999999996 \\
-0.0272727272727274 0 \\
-0.0272727272727274 -0.0999999999999996 \\
-0.0272727272727274 -0.199999999999999 \\
-0.0272727272727274 -0.3 \\
-0.0272727272727274 -0.399999999999999 \\
-0.0272727272727274 -0.5 \\
-0.0272727272727274 -0.6 \\
-0.0272727272727274 -0.699999999999999 \\
-0.0272727272727274 -0.8 \\
-0.0272727272727274 -0.899999999999999 \\
-0.0272727272727274 -1 \\
-0.1 -1.07272727272727 \\
-0.127272727272727 -1.1 \\
-0.127272727272727 -1.2 \\
-0.2 -1.27272727272727 \\
-0.227272727272728 -1.3 \\
-0.3 -1.37272727272727 \\
-0.327272727272728 -1.4 \\
-0.4 -1.47272727272727 \\
-0.427272727272728 -1.5 \\
-0.5 -1.57272727272727 \\
-0.527272727272727 -1.6 \\
-0.6 -1.67272727272727 \\
-0.7 -1.67272727272727 \\
-0.727272727272728 -1.7 \\
-0.8 -1.77272727272727 \\
-0.9 -1.77272727272727 \\
-1 -1.77272727272727 \\
-1.1 -1.77272727272727 \\
-1.2 -1.77272727272727 \\
-1.27272727272727 -1.7 \\
-1.3 -1.67272727272727 \\
-1.4 -1.67272727272727 \\
-1.47272727272727 -1.6 \\
-1.5 -1.57272727272727 \\
-1.57272727272727 -1.5 \\
-1.6 -1.47272727272727 \\
-1.67272727272727 -1.4 \\
-1.7 -1.37272727272727 \\
-1.77272727272727 -1.3 \\
-1.8 -1.27272727272727 \\
-1.87272727272727 -1.2 \\
-1.87272727272727 -1.1 \\
-1.9 -1.07272727272727 \\
-1.97272727272727 -1 \\
-1.97272727272727 -0.899999999999999 \\
-1.97272727272727 -0.8 \\
-1.97272727272727 -0.699999999999999 \\
-1.97272727272727 -0.6 \\
-1.97272727272727 -0.5 \\
-1.97272727272727 -0.399999999999999 \\
-1.97272727272727 -0.3 \\
-1.97272727272727 -0.199999999999999 \\
-1.97272727272727 -0.0999999999999996 \\
-1.97272727272727 0 \\
-1.97272727272727 0.0999999999999996 \\
-1.97272727272727 0.199999999999999 \\
-1.97272727272727 0.3 \\
-1.97272727272727 0.399999999999999 \\
-1.97272727272727 0.5 \\
-1.97272727272727 0.6 \\
-1.97272727272727 0.699999999999999 \\
-1.97272727272727 0.8 \\
-1.97272727272727 0.899999999999999 \\
-1.97272727272727 1 \\
-1.9 1.07272727272727 \\
NaN NaN \\
};


\addplot[dashed,line width=2.0pt,draw=blue,forget plot]
table[row sep=crcr]{
x y\\
-1.9 1.06363636363636 \\
-1.86363636363636 1.1 \\
-1.86363636363636 1.2 \\
-1.8 1.26363636363636 \\
-1.76363636363636 1.3 \\
-1.7 1.36363636363636 \\
-1.66363636363636 1.4 \\
-1.6 1.46363636363636 \\
-1.56363636363636 1.5 \\
-1.5 1.56363636363636 \\
-1.46363636363636 1.6 \\
-1.4 1.66363636363636 \\
-1.3 1.66363636363636 \\
-1.26363636363636 1.7 \\
-1.2 1.76363636363636 \\
-1.1 1.76363636363636 \\
-1 1.76363636363636 \\
-0.9 1.76363636363636 \\
-0.8 1.76363636363636 \\
-0.736363636363637 1.7 \\
-0.7 1.66363636363636 \\
-0.6 1.66363636363636 \\
-0.536363636363637 1.6 \\
-0.5 1.56363636363636 \\
-0.436363636363637 1.5 \\
-0.4 1.46363636363636 \\
-0.336363636363637 1.4 \\
-0.3 1.36363636363636 \\
-0.236363636363637 1.3 \\
-0.2 1.26363636363636 \\
-0.136363636363637 1.2 \\
-0.136363636363637 1.1 \\
-0.1 1.06363636363636 \\
-0.0363636363636367 1 \\
-0.0363636363636367 0.899999999999999 \\
-0.0363636363636367 0.8 \\
-0.0363636363636367 0.699999999999999 \\
-0.0363636363636367 0.6 \\
-0.0363636363636367 0.5 \\
-0.0363636363636367 0.399999999999999 \\
-0.0363636363636367 0.3 \\
-0.0363636363636367 0.199999999999999 \\
-0.0363636363636367 0.0999999999999996 \\
-0.0363636363636367 0 \\
-0.0363636363636367 -0.0999999999999996 \\
-0.0363636363636367 -0.199999999999999 \\
-0.0363636363636367 -0.3 \\
-0.0363636363636367 -0.399999999999999 \\
-0.0363636363636367 -0.5 \\
-0.0363636363636367 -0.6 \\
-0.0363636363636367 -0.699999999999999 \\
-0.0363636363636367 -0.8 \\
-0.0363636363636367 -0.899999999999999 \\
-0.0363636363636367 -1 \\
-0.1 -1.06363636363636 \\
-0.136363636363637 -1.1 \\
-0.136363636363637 -1.2 \\
-0.2 -1.26363636363636 \\
-0.236363636363637 -1.3 \\
-0.3 -1.36363636363636 \\
-0.336363636363637 -1.4 \\
-0.4 -1.46363636363636 \\
-0.436363636363637 -1.5 \\
-0.5 -1.56363636363636 \\
-0.536363636363637 -1.6 \\
-0.6 -1.66363636363636 \\
-0.7 -1.66363636363636 \\
-0.736363636363637 -1.7 \\
-0.8 -1.76363636363636 \\
-0.9 -1.76363636363636 \\
-1 -1.76363636363636 \\
-1.1 -1.76363636363636 \\
-1.2 -1.76363636363636 \\
-1.26363636363636 -1.7 \\
-1.3 -1.66363636363636 \\
-1.4 -1.66363636363636 \\
-1.46363636363636 -1.6 \\
-1.5 -1.56363636363636 \\
-1.56363636363636 -1.5 \\
-1.6 -1.46363636363636 \\
-1.66363636363636 -1.4 \\
-1.7 -1.36363636363636 \\
-1.76363636363636 -1.3 \\
-1.8 -1.26363636363636 \\
-1.86363636363636 -1.2 \\
-1.86363636363636 -1.1 \\
-1.9 -1.06363636363636 \\
-1.96363636363636 -1 \\
-1.96363636363636 -0.899999999999999 \\
-1.96363636363636 -0.8 \\
-1.96363636363636 -0.699999999999999 \\
-1.96363636363636 -0.6 \\
-1.96363636363636 -0.5 \\
-1.96363636363636 -0.399999999999999 \\
-1.96363636363636 -0.3 \\
-1.96363636363636 -0.199999999999999 \\
-1.96363636363636 -0.0999999999999996 \\
-1.96363636363636 0 \\
-1.96363636363636 0.0999999999999996 \\
-1.96363636363636 0.199999999999999 \\
-1.96363636363636 0.3 \\
-1.96363636363636 0.399999999999999 \\
-1.96363636363636 0.5 \\
-1.96363636363636 0.6 \\
-1.96363636363636 0.699999999999999 \\
-1.96363636363636 0.8 \\
-1.96363636363636 0.899999999999999 \\
-1.96363636363636 1 \\
-1.9 1.06363636363636 \\
NaN NaN \\
};


\addplot[dashed,line width=2.0pt,draw=blue,forget plot]
table[row sep=crcr]{
x y\\
-1.9 1.05454545454545 \\
-1.85454545454545 1.1 \\
-1.85454545454545 1.2 \\
-1.8 1.25454545454545 \\
-1.75454545454545 1.3 \\
-1.7 1.35454545454545 \\
-1.65454545454545 1.4 \\
-1.6 1.45454545454545 \\
-1.55454545454545 1.5 \\
-1.5 1.55454545454545 \\
-1.45454545454545 1.6 \\
-1.4 1.65454545454545 \\
-1.3 1.65454545454545 \\
-1.25454545454545 1.7 \\
-1.2 1.75454545454545 \\
-1.1 1.75454545454545 \\
-1 1.75454545454545 \\
-0.9 1.75454545454545 \\
-0.8 1.75454545454545 \\
-0.745454545454546 1.7 \\
-0.7 1.65454545454545 \\
-0.6 1.65454545454545 \\
-0.545454545454545 1.6 \\
-0.5 1.55454545454545 \\
-0.445454545454546 1.5 \\
-0.4 1.45454545454545 \\
-0.345454545454546 1.4 \\
-0.3 1.35454545454545 \\
-0.245454545454546 1.3 \\
-0.2 1.25454545454545 \\
-0.145454545454545 1.2 \\
-0.145454545454545 1.1 \\
-0.1 1.05454545454545 \\
-0.0454545454545454 1 \\
-0.0454545454545454 0.899999999999999 \\
-0.0454545454545454 0.8 \\
-0.0454545454545454 0.699999999999999 \\
-0.0454545454545454 0.6 \\
-0.0454545454545454 0.5 \\
-0.0454545454545454 0.399999999999999 \\
-0.0454545454545454 0.3 \\
-0.0454545454545454 0.199999999999999 \\
-0.0454545454545454 0.0999999999999996 \\
-0.0454545454545454 0 \\
-0.0454545454545454 -0.0999999999999996 \\
-0.0454545454545454 -0.199999999999999 \\
-0.0454545454545454 -0.3 \\
-0.0454545454545454 -0.399999999999999 \\
-0.0454545454545454 -0.5 \\
-0.0454545454545454 -0.6 \\
-0.0454545454545454 -0.699999999999999 \\
-0.0454545454545454 -0.8 \\
-0.0454545454545454 -0.899999999999999 \\
-0.0454545454545454 -1 \\
-0.1 -1.05454545454545 \\
-0.145454545454545 -1.1 \\
-0.145454545454545 -1.2 \\
-0.2 -1.25454545454545 \\
-0.245454545454546 -1.3 \\
-0.3 -1.35454545454545 \\
-0.345454545454546 -1.4 \\
-0.4 -1.45454545454545 \\
-0.445454545454546 -1.5 \\
-0.5 -1.55454545454545 \\
-0.545454545454545 -1.6 \\
-0.6 -1.65454545454545 \\
-0.7 -1.65454545454545 \\
-0.745454545454546 -1.7 \\
-0.8 -1.75454545454545 \\
-0.9 -1.75454545454545 \\
-1 -1.75454545454545 \\
-1.1 -1.75454545454545 \\
-1.2 -1.75454545454545 \\
-1.25454545454545 -1.7 \\
-1.3 -1.65454545454545 \\
-1.4 -1.65454545454545 \\
-1.45454545454545 -1.6 \\
-1.5 -1.55454545454545 \\
-1.55454545454545 -1.5 \\
-1.6 -1.45454545454545 \\
-1.65454545454545 -1.4 \\
-1.7 -1.35454545454545 \\
-1.75454545454545 -1.3 \\
-1.8 -1.25454545454545 \\
-1.85454545454545 -1.2 \\
-1.85454545454545 -1.1 \\
-1.9 -1.05454545454545 \\
-1.95454545454545 -1 \\
-1.95454545454545 -0.899999999999999 \\
-1.95454545454545 -0.8 \\
-1.95454545454545 -0.699999999999999 \\
-1.95454545454545 -0.6 \\
-1.95454545454545 -0.5 \\
-1.95454545454545 -0.399999999999999 \\
-1.95454545454545 -0.3 \\
-1.95454545454545 -0.199999999999999 \\
-1.95454545454545 -0.0999999999999996 \\
-1.95454545454545 0 \\
-1.95454545454545 0.0999999999999996 \\
-1.95454545454545 0.199999999999999 \\
-1.95454545454545 0.3 \\
-1.95454545454545 0.399999999999999 \\
-1.95454545454545 0.5 \\
-1.95454545454545 0.6 \\
-1.95454545454545 0.699999999999999 \\
-1.95454545454545 0.8 \\
-1.95454545454545 0.899999999999999 \\
-1.95454545454545 1 \\
-1.9 1.05454545454545 \\
NaN NaN \\
};


\addplot[dashed,line width=2.0pt,draw=blue,forget plot]
table[row sep=crcr]{
x y\\
-1.9 1.04545454545455 \\
-1.84545454545455 1.1 \\
-1.84545454545455 1.2 \\
-1.8 1.24545454545454 \\
-1.74545454545455 1.3 \\
-1.7 1.34545454545455 \\
-1.64545454545455 1.4 \\
-1.6 1.44545454545454 \\
-1.54545454545455 1.5 \\
-1.5 1.54545454545455 \\
-1.44545454545455 1.6 \\
-1.4 1.64545454545455 \\
-1.3 1.64545454545455 \\
-1.24545454545455 1.7 \\
-1.2 1.74545454545455 \\
-1.1 1.74545454545455 \\
-1 1.74545454545455 \\
-0.9 1.74545454545455 \\
-0.8 1.74545454545455 \\
-0.754545454545455 1.7 \\
-0.7 1.64545454545455 \\
-0.6 1.64545454545455 \\
-0.554545454545455 1.6 \\
-0.5 1.54545454545455 \\
-0.454545454545455 1.5 \\
-0.4 1.44545454545454 \\
-0.354545454545455 1.4 \\
-0.3 1.34545454545455 \\
-0.254545454545455 1.3 \\
-0.2 1.24545454545454 \\
-0.154545454545455 1.2 \\
-0.154545454545455 1.1 \\
-0.1 1.04545454545455 \\
-0.0545454545454547 1 \\
-0.0545454545454547 0.899999999999999 \\
-0.0545454545454547 0.8 \\
-0.0545454545454547 0.699999999999999 \\
-0.0545454545454547 0.6 \\
-0.0545454545454547 0.5 \\
-0.0545454545454547 0.399999999999999 \\
-0.0545454545454547 0.3 \\
-0.0545454545454547 0.199999999999999 \\
-0.0545454545454547 0.0999999999999996 \\
-0.0545454545454547 0 \\
-0.0545454545454547 -0.0999999999999996 \\
-0.0545454545454547 -0.199999999999999 \\
-0.0545454545454547 -0.3 \\
-0.0545454545454547 -0.399999999999999 \\
-0.0545454545454547 -0.5 \\
-0.0545454545454547 -0.6 \\
-0.0545454545454547 -0.699999999999999 \\
-0.0545454545454547 -0.8 \\
-0.0545454545454547 -0.899999999999999 \\
-0.0545454545454547 -1 \\
-0.1 -1.04545454545455 \\
-0.154545454545455 -1.1 \\
-0.154545454545455 -1.2 \\
-0.2 -1.24545454545454 \\
-0.254545454545455 -1.3 \\
-0.3 -1.34545454545455 \\
-0.354545454545455 -1.4 \\
-0.4 -1.44545454545454 \\
-0.454545454545455 -1.5 \\
-0.5 -1.54545454545455 \\
-0.554545454545455 -1.6 \\
-0.6 -1.64545454545455 \\
-0.7 -1.64545454545455 \\
-0.754545454545455 -1.7 \\
-0.8 -1.74545454545455 \\
-0.9 -1.74545454545455 \\
-1 -1.74545454545455 \\
-1.1 -1.74545454545455 \\
-1.2 -1.74545454545455 \\
-1.24545454545455 -1.7 \\
-1.3 -1.64545454545455 \\
-1.4 -1.64545454545455 \\
-1.44545454545455 -1.6 \\
-1.5 -1.54545454545455 \\
-1.54545454545455 -1.5 \\
-1.6 -1.44545454545454 \\
-1.64545454545455 -1.4 \\
-1.7 -1.34545454545455 \\
-1.74545454545455 -1.3 \\
-1.8 -1.24545454545454 \\
-1.84545454545455 -1.2 \\
-1.84545454545455 -1.1 \\
-1.9 -1.04545454545455 \\
-1.94545454545455 -1 \\
-1.94545454545455 -0.899999999999999 \\
-1.94545454545455 -0.8 \\
-1.94545454545455 -0.699999999999999 \\
-1.94545454545455 -0.6 \\
-1.94545454545455 -0.5 \\
-1.94545454545455 -0.399999999999999 \\
-1.94545454545455 -0.3 \\
-1.94545454545455 -0.199999999999999 \\
-1.94545454545455 -0.0999999999999996 \\
-1.94545454545455 0 \\
-1.94545454545455 0.0999999999999996 \\
-1.94545454545455 0.199999999999999 \\
-1.94545454545455 0.3 \\
-1.94545454545455 0.399999999999999 \\
-1.94545454545455 0.5 \\
-1.94545454545455 0.6 \\
-1.94545454545455 0.699999999999999 \\
-1.94545454545455 0.8 \\
-1.94545454545455 0.899999999999999 \\
-1.94545454545455 1 \\
-1.9 1.04545454545455 \\
NaN NaN \\
};


\addplot[dashed,line width=2.0pt,draw=blue,forget plot]
table[row sep=crcr]{
x y\\
-1.9 1.03636363636364 \\
-1.83636363636364 1.1 \\
-1.83636363636364 1.2 \\
-1.8 1.23636363636364 \\
-1.73636363636364 1.3 \\
-1.7 1.33636363636364 \\
-1.63636363636364 1.4 \\
-1.6 1.43636363636364 \\
-1.53636363636364 1.5 \\
-1.5 1.53636363636364 \\
-1.43636363636364 1.6 \\
-1.4 1.63636363636364 \\
-1.3 1.63636363636364 \\
-1.23636363636364 1.7 \\
-1.2 1.73636363636364 \\
-1.1 1.73636363636364 \\
-1 1.73636363636364 \\
-0.9 1.73636363636364 \\
-0.8 1.73636363636364 \\
-0.763636363636364 1.7 \\
-0.7 1.63636363636364 \\
-0.6 1.63636363636364 \\
-0.563636363636363 1.6 \\
-0.5 1.53636363636364 \\
-0.463636363636363 1.5 \\
-0.4 1.43636363636364 \\
-0.363636363636364 1.4 \\
-0.3 1.33636363636364 \\
-0.263636363636364 1.3 \\
-0.2 1.23636363636364 \\
-0.163636363636363 1.2 \\
-0.163636363636363 1.1 \\
-0.1 1.03636363636364 \\
-0.0636363636363634 1 \\
-0.0636363636363634 0.899999999999999 \\
-0.0636363636363634 0.8 \\
-0.0636363636363634 0.699999999999999 \\
-0.0636363636363634 0.6 \\
-0.0636363636363634 0.5 \\
-0.0636363636363634 0.399999999999999 \\
-0.0636363636363634 0.3 \\
-0.0636363636363634 0.199999999999999 \\
-0.0636363636363634 0.0999999999999996 \\
-0.0636363636363634 0 \\
-0.0636363636363634 -0.0999999999999996 \\
-0.0636363636363634 -0.199999999999999 \\
-0.0636363636363634 -0.3 \\
-0.0636363636363634 -0.399999999999999 \\
-0.0636363636363634 -0.5 \\
-0.0636363636363634 -0.6 \\
-0.0636363636363634 -0.699999999999999 \\
-0.0636363636363634 -0.8 \\
-0.0636363636363634 -0.899999999999999 \\
-0.0636363636363634 -1 \\
-0.1 -1.03636363636364 \\
-0.163636363636363 -1.1 \\
-0.163636363636363 -1.2 \\
-0.2 -1.23636363636364 \\
-0.263636363636364 -1.3 \\
-0.3 -1.33636363636364 \\
-0.363636363636364 -1.4 \\
-0.4 -1.43636363636364 \\
-0.463636363636363 -1.5 \\
-0.5 -1.53636363636364 \\
-0.563636363636363 -1.6 \\
-0.6 -1.63636363636364 \\
-0.7 -1.63636363636364 \\
-0.763636363636364 -1.7 \\
-0.8 -1.73636363636364 \\
-0.9 -1.73636363636364 \\
-1 -1.73636363636364 \\
-1.1 -1.73636363636364 \\
-1.2 -1.73636363636364 \\
-1.23636363636364 -1.7 \\
-1.3 -1.63636363636364 \\
-1.4 -1.63636363636364 \\
-1.43636363636364 -1.6 \\
-1.5 -1.53636363636364 \\
-1.53636363636364 -1.5 \\
-1.6 -1.43636363636364 \\
-1.63636363636364 -1.4 \\
-1.7 -1.33636363636364 \\
-1.73636363636364 -1.3 \\
-1.8 -1.23636363636364 \\
-1.83636363636364 -1.2 \\
-1.83636363636364 -1.1 \\
-1.9 -1.03636363636364 \\
-1.93636363636364 -1 \\
-1.93636363636364 -0.899999999999999 \\
-1.93636363636364 -0.8 \\
-1.93636363636364 -0.699999999999999 \\
-1.93636363636364 -0.6 \\
-1.93636363636364 -0.5 \\
-1.93636363636364 -0.399999999999999 \\
-1.93636363636364 -0.3 \\
-1.93636363636364 -0.199999999999999 \\
-1.93636363636364 -0.0999999999999996 \\
-1.93636363636364 0 \\
-1.93636363636364 0.0999999999999996 \\
-1.93636363636364 0.199999999999999 \\
-1.93636363636364 0.3 \\
-1.93636363636364 0.399999999999999 \\
-1.93636363636364 0.5 \\
-1.93636363636364 0.6 \\
-1.93636363636364 0.699999999999999 \\
-1.93636363636364 0.8 \\
-1.93636363636364 0.899999999999999 \\
-1.93636363636364 1 \\
-1.9 1.03636363636364 \\
NaN NaN \\
};


\addplot[dashed,line width=2.0pt,draw=blue,forget plot]
table[row sep=crcr]{
x y\\
-1.9 1.02727272727273 \\
-1.82727272727273 1.1 \\
-1.82727272727273 1.2 \\
-1.8 1.22727272727273 \\
-1.72727272727273 1.3 \\
-1.7 1.32727272727273 \\
-1.62727272727273 1.4 \\
-1.6 1.42727272727273 \\
-1.52727272727273 1.5 \\
-1.5 1.52727272727273 \\
-1.42727272727273 1.6 \\
-1.4 1.62727272727273 \\
-1.3 1.62727272727273 \\
-1.22727272727273 1.7 \\
-1.2 1.72727272727273 \\
-1.1 1.72727272727273 \\
-1 1.72727272727273 \\
-0.9 1.72727272727273 \\
-0.8 1.72727272727273 \\
-0.772727272727273 1.7 \\
-0.7 1.62727272727273 \\
-0.6 1.62727272727273 \\
-0.572727272727273 1.6 \\
-0.5 1.52727272727273 \\
-0.472727272727273 1.5 \\
-0.4 1.42727272727273 \\
-0.372727272727273 1.4 \\
-0.3 1.32727272727273 \\
-0.272727272727273 1.3 \\
-0.2 1.22727272727273 \\
-0.172727272727273 1.2 \\
-0.172727272727273 1.1 \\
-0.1 1.02727272727273 \\
-0.0727272727272727 1 \\
-0.0727272727272727 0.899999999999999 \\
-0.0727272727272727 0.8 \\
-0.0727272727272727 0.699999999999999 \\
-0.0727272727272727 0.6 \\
-0.0727272727272727 0.5 \\
-0.0727272727272727 0.399999999999999 \\
-0.0727272727272727 0.3 \\
-0.0727272727272727 0.199999999999999 \\
-0.0727272727272727 0.0999999999999996 \\
-0.0727272727272727 0 \\
-0.0727272727272727 -0.0999999999999996 \\
-0.0727272727272727 -0.199999999999999 \\
-0.0727272727272727 -0.3 \\
-0.0727272727272727 -0.399999999999999 \\
-0.0727272727272727 -0.5 \\
-0.0727272727272727 -0.6 \\
-0.0727272727272727 -0.699999999999999 \\
-0.0727272727272727 -0.8 \\
-0.0727272727272727 -0.899999999999999 \\
-0.0727272727272727 -1 \\
-0.1 -1.02727272727273 \\
-0.172727272727273 -1.1 \\
-0.172727272727273 -1.2 \\
-0.2 -1.22727272727273 \\
-0.272727272727273 -1.3 \\
-0.3 -1.32727272727273 \\
-0.372727272727273 -1.4 \\
-0.4 -1.42727272727273 \\
-0.472727272727273 -1.5 \\
-0.5 -1.52727272727273 \\
-0.572727272727273 -1.6 \\
-0.6 -1.62727272727273 \\
-0.7 -1.62727272727273 \\
-0.772727272727273 -1.7 \\
-0.8 -1.72727272727273 \\
-0.9 -1.72727272727273 \\
-1 -1.72727272727273 \\
-1.1 -1.72727272727273 \\
-1.2 -1.72727272727273 \\
-1.22727272727273 -1.7 \\
-1.3 -1.62727272727273 \\
-1.4 -1.62727272727273 \\
-1.42727272727273 -1.6 \\
-1.5 -1.52727272727273 \\
-1.52727272727273 -1.5 \\
-1.6 -1.42727272727273 \\
-1.62727272727273 -1.4 \\
-1.7 -1.32727272727273 \\
-1.72727272727273 -1.3 \\
-1.8 -1.22727272727273 \\
-1.82727272727273 -1.2 \\
-1.82727272727273 -1.1 \\
-1.9 -1.02727272727273 \\
-1.92727272727273 -1 \\
-1.92727272727273 -0.899999999999999 \\
-1.92727272727273 -0.8 \\
-1.92727272727273 -0.699999999999999 \\
-1.92727272727273 -0.6 \\
-1.92727272727273 -0.5 \\
-1.92727272727273 -0.399999999999999 \\
-1.92727272727273 -0.3 \\
-1.92727272727273 -0.199999999999999 \\
-1.92727272727273 -0.0999999999999996 \\
-1.92727272727273 0 \\
-1.92727272727273 0.0999999999999996 \\
-1.92727272727273 0.199999999999999 \\
-1.92727272727273 0.3 \\
-1.92727272727273 0.399999999999999 \\
-1.92727272727273 0.5 \\
-1.92727272727273 0.6 \\
-1.92727272727273 0.699999999999999 \\
-1.92727272727273 0.8 \\
-1.92727272727273 0.899999999999999 \\
-1.92727272727273 1 \\
-1.9 1.02727272727273 \\
NaN NaN \\
};


\addplot[dashed,line width=2.0pt,draw=blue,forget plot]
table[row sep=crcr]{
x y\\
-1.9 1.01818181818182 \\
-1.81818181818182 1.1 \\
-1.81818181818182 1.2 \\
-1.8 1.21818181818182 \\
-1.71818181818182 1.3 \\
-1.7 1.31818181818182 \\
-1.61818181818182 1.4 \\
-1.6 1.41818181818182 \\
-1.51818181818182 1.5 \\
-1.5 1.51818181818182 \\
-1.41818181818182 1.6 \\
-1.4 1.61818181818182 \\
-1.3 1.61818181818182 \\
-1.21818181818182 1.7 \\
-1.2 1.71818181818182 \\
-1.1 1.71818181818182 \\
-1 1.71818181818182 \\
-0.9 1.71818181818182 \\
-0.8 1.71818181818182 \\
-0.781818181818182 1.7 \\
-0.7 1.61818181818182 \\
-0.6 1.61818181818182 \\
-0.581818181818182 1.6 \\
-0.5 1.51818181818182 \\
-0.481818181818182 1.5 \\
-0.4 1.41818181818182 \\
-0.381818181818182 1.4 \\
-0.3 1.31818181818182 \\
-0.281818181818182 1.3 \\
-0.2 1.21818181818182 \\
-0.181818181818182 1.2 \\
-0.181818181818182 1.1 \\
-0.1 1.01818181818182 \\
-0.0818181818181821 1 \\
-0.0818181818181821 0.899999999999999 \\
-0.0818181818181821 0.8 \\
-0.0818181818181821 0.699999999999999 \\
-0.0818181818181821 0.6 \\
-0.0818181818181821 0.5 \\
-0.0818181818181821 0.399999999999999 \\
-0.0818181818181821 0.3 \\
-0.0818181818181821 0.199999999999999 \\
-0.0818181818181821 0.0999999999999996 \\
-0.0818181818181821 0 \\
-0.0818181818181821 -0.0999999999999996 \\
-0.0818181818181821 -0.199999999999999 \\
-0.0818181818181821 -0.3 \\
-0.0818181818181821 -0.399999999999999 \\
-0.0818181818181821 -0.5 \\
-0.0818181818181821 -0.6 \\
-0.0818181818181821 -0.699999999999999 \\
-0.0818181818181821 -0.8 \\
-0.0818181818181821 -0.899999999999999 \\
-0.0818181818181821 -1 \\
-0.1 -1.01818181818182 \\
-0.181818181818182 -1.1 \\
-0.181818181818182 -1.2 \\
-0.2 -1.21818181818182 \\
-0.281818181818182 -1.3 \\
-0.3 -1.31818181818182 \\
-0.381818181818182 -1.4 \\
-0.4 -1.41818181818182 \\
-0.481818181818182 -1.5 \\
-0.5 -1.51818181818182 \\
-0.581818181818182 -1.6 \\
-0.6 -1.61818181818182 \\
-0.7 -1.61818181818182 \\
-0.781818181818182 -1.7 \\
-0.8 -1.71818181818182 \\
-0.9 -1.71818181818182 \\
-1 -1.71818181818182 \\
-1.1 -1.71818181818182 \\
-1.2 -1.71818181818182 \\
-1.21818181818182 -1.7 \\
-1.3 -1.61818181818182 \\
-1.4 -1.61818181818182 \\
-1.41818181818182 -1.6 \\
-1.5 -1.51818181818182 \\
-1.51818181818182 -1.5 \\
-1.6 -1.41818181818182 \\
-1.61818181818182 -1.4 \\
-1.7 -1.31818181818182 \\
-1.71818181818182 -1.3 \\
-1.8 -1.21818181818182 \\
-1.81818181818182 -1.2 \\
-1.81818181818182 -1.1 \\
-1.9 -1.01818181818182 \\
-1.91818181818182 -1 \\
-1.91818181818182 -0.899999999999999 \\
-1.91818181818182 -0.8 \\
-1.91818181818182 -0.699999999999999 \\
-1.91818181818182 -0.6 \\
-1.91818181818182 -0.5 \\
-1.91818181818182 -0.399999999999999 \\
-1.91818181818182 -0.3 \\
-1.91818181818182 -0.199999999999999 \\
-1.91818181818182 -0.0999999999999996 \\
-1.91818181818182 0 \\
-1.91818181818182 0.0999999999999996 \\
-1.91818181818182 0.199999999999999 \\
-1.91818181818182 0.3 \\
-1.91818181818182 0.399999999999999 \\
-1.91818181818182 0.5 \\
-1.91818181818182 0.6 \\
-1.91818181818182 0.699999999999999 \\
-1.91818181818182 0.8 \\
-1.91818181818182 0.899999999999999 \\
-1.91818181818182 1 \\
-1.9 1.01818181818182 \\
NaN NaN \\
};


\addplot[dashed,line width=2.0pt,draw=blue,forget plot]
table[row sep=crcr]{
x y\\
-1.9 1.00909090909091 \\
-1.80909090909091 1.1 \\
-1.80909090909091 1.2 \\
-1.8 1.20909090909091 \\
-1.70909090909091 1.3 \\
-1.7 1.30909090909091 \\
-1.60909090909091 1.4 \\
-1.6 1.40909090909091 \\
-1.50909090909091 1.5 \\
-1.5 1.50909090909091 \\
-1.40909090909091 1.6 \\
-1.4 1.60909090909091 \\
-1.3 1.60909090909091 \\
-1.20909090909091 1.7 \\
-1.2 1.70909090909091 \\
-1.1 1.70909090909091 \\
-1 1.70909090909091 \\
-0.9 1.70909090909091 \\
-0.8 1.70909090909091 \\
-0.790909090909091 1.7 \\
-0.7 1.60909090909091 \\
-0.6 1.60909090909091 \\
-0.590909090909091 1.6 \\
-0.5 1.50909090909091 \\
-0.490909090909091 1.5 \\
-0.4 1.40909090909091 \\
-0.390909090909091 1.4 \\
-0.3 1.30909090909091 \\
-0.290909090909091 1.3 \\
-0.2 1.20909090909091 \\
-0.190909090909091 1.2 \\
-0.190909090909091 1.1 \\
-0.1 1.00909090909091 \\
-0.0909090909090907 1 \\
-0.0909090909090907 0.899999999999999 \\
-0.0909090909090907 0.8 \\
-0.0909090909090907 0.699999999999999 \\
-0.0909090909090907 0.6 \\
-0.0909090909090907 0.5 \\
-0.0909090909090907 0.399999999999999 \\
-0.0909090909090907 0.3 \\
-0.0909090909090907 0.199999999999999 \\
-0.0909090909090907 0.0999999999999996 \\
-0.0909090909090907 0 \\
-0.0909090909090907 -0.0999999999999996 \\
-0.0909090909090907 -0.199999999999999 \\
-0.0909090909090907 -0.3 \\
-0.0909090909090907 -0.399999999999999 \\
-0.0909090909090907 -0.5 \\
-0.0909090909090907 -0.6 \\
-0.0909090909090907 -0.699999999999999 \\
-0.0909090909090907 -0.8 \\
-0.0909090909090907 -0.899999999999999 \\
-0.0909090909090907 -1 \\
-0.1 -1.00909090909091 \\
-0.190909090909091 -1.1 \\
-0.190909090909091 -1.2 \\
-0.2 -1.20909090909091 \\
-0.290909090909091 -1.3 \\
-0.3 -1.30909090909091 \\
-0.390909090909091 -1.4 \\
-0.4 -1.40909090909091 \\
-0.490909090909091 -1.5 \\
-0.5 -1.50909090909091 \\
-0.590909090909091 -1.6 \\
-0.6 -1.60909090909091 \\
-0.7 -1.60909090909091 \\
-0.790909090909091 -1.7 \\
-0.8 -1.70909090909091 \\
-0.9 -1.70909090909091 \\
-1 -1.70909090909091 \\
-1.1 -1.70909090909091 \\
-1.2 -1.70909090909091 \\
-1.20909090909091 -1.7 \\
-1.3 -1.60909090909091 \\
-1.4 -1.60909090909091 \\
-1.40909090909091 -1.6 \\
-1.5 -1.50909090909091 \\
-1.50909090909091 -1.5 \\
-1.6 -1.40909090909091 \\
-1.60909090909091 -1.4 \\
-1.7 -1.30909090909091 \\
-1.70909090909091 -1.3 \\
-1.8 -1.20909090909091 \\
-1.80909090909091 -1.2 \\
-1.80909090909091 -1.1 \\
-1.9 -1.00909090909091 \\
-1.90909090909091 -1 \\
-1.90909090909091 -0.899999999999999 \\
-1.90909090909091 -0.8 \\
-1.90909090909091 -0.699999999999999 \\
-1.90909090909091 -0.6 \\
-1.90909090909091 -0.5 \\
-1.90909090909091 -0.399999999999999 \\
-1.90909090909091 -0.3 \\
-1.90909090909091 -0.199999999999999 \\
-1.90909090909091 -0.0999999999999996 \\
-1.90909090909091 0 \\
-1.90909090909091 0.0999999999999996 \\
-1.90909090909091 0.199999999999999 \\
-1.90909090909091 0.3 \\
-1.90909090909091 0.399999999999999 \\
-1.90909090909091 0.5 \\
-1.90909090909091 0.6 \\
-1.90909090909091 0.699999999999999 \\
-1.90909090909091 0.8 \\
-1.90909090909091 0.899999999999999 \\
-1.90909090909091 1 \\
-1.9 1.00909090909091 \\
NaN NaN \\
};


\addplot[dashed,line width=2.0pt,draw=blue,forget plot]
table[row sep=crcr]{
x y\\
-1.9 1 \\
-1.8 1.1 \\
-1.8 1.2 \\
-1.8 1.2 \\
-1.7 1.3 \\
-1.7 1.3 \\
-1.6 1.4 \\
-1.6 1.4 \\
-1.5 1.5 \\
-1.5 1.5 \\
-1.4 1.6 \\
-1.4 1.6 \\
-1.3 1.6 \\
-1.2 1.7 \\
-1.2 1.7 \\
-1.1 1.7 \\
-1 1.7 \\
-0.9 1.7 \\
-0.8 1.7 \\
-0.8 1.7 \\
-0.7 1.6 \\
-0.6 1.6 \\
-0.6 1.6 \\
-0.5 1.5 \\
-0.5 1.5 \\
-0.4 1.4 \\
-0.4 1.4 \\
-0.3 1.3 \\
-0.3 1.3 \\
-0.2 1.2 \\
-0.2 1.2 \\
-0.2 1.1 \\
-0.1 1 \\
-0.1 1 \\
-0.1 0.899999999999999 \\
-0.1 0.8 \\
-0.1 0.699999999999999 \\
-0.1 0.6 \\
-0.1 0.5 \\
-0.1 0.399999999999999 \\
-0.1 0.3 \\
-0.1 0.199999999999999 \\
-0.1 0.0999999999999996 \\
-0.1 0 \\
-0.1 -0.0999999999999996 \\
-0.1 -0.199999999999999 \\
-0.1 -0.3 \\
-0.1 -0.399999999999999 \\
-0.1 -0.5 \\
-0.1 -0.6 \\
-0.1 -0.699999999999999 \\
-0.1 -0.8 \\
-0.1 -0.899999999999999 \\
-0.1 -1 \\
-0.1 -1 \\
-0.2 -1.1 \\
-0.2 -1.2 \\
-0.2 -1.2 \\
-0.3 -1.3 \\
-0.3 -1.3 \\
-0.4 -1.4 \\
-0.4 -1.4 \\
-0.5 -1.5 \\
-0.5 -1.5 \\
-0.6 -1.6 \\
-0.6 -1.6 \\
-0.7 -1.6 \\
-0.8 -1.7 \\
-0.8 -1.7 \\
-0.9 -1.7 \\
-1 -1.7 \\
-1.1 -1.7 \\
-1.2 -1.7 \\
-1.2 -1.7 \\
-1.3 -1.6 \\
-1.4 -1.6 \\
-1.4 -1.6 \\
-1.5 -1.5 \\
-1.5 -1.5 \\
-1.6 -1.4 \\
-1.6 -1.4 \\
-1.7 -1.3 \\
-1.7 -1.3 \\
-1.8 -1.2 \\
-1.8 -1.2 \\
-1.8 -1.1 \\
-1.9 -1 \\
-1.9 -1 \\
-1.9 -0.899999999999999 \\
-1.9 -0.8 \\
-1.9 -0.699999999999999 \\
-1.9 -0.6 \\
-1.9 -0.5 \\
-1.9 -0.399999999999999 \\
-1.9 -0.3 \\
-1.9 -0.199999999999999 \\
-1.9 -0.0999999999999996 \\
-1.9 0 \\
-1.9 0.0999999999999996 \\
-1.9 0.199999999999999 \\
-1.9 0.3 \\
-1.9 0.399999999999999 \\
-1.9 0.5 \\
-1.9 0.6 \\
-1.9 0.699999999999999 \\
-1.9 0.8 \\
-1.9 0.899999999999999 \\
-1.9 1 \\
-1.9 1 \\
NaN NaN \\
};


\addplot[dashed,line width=2.0pt,draw=mycolor1,forget plot]
table[row sep=crcr]{
x y\\
-2.7 0.89999999 \\
-2.69999999 0.899999999999999 \\
-2.69999999 1 \\
-2.69999999 1.1 \\
-2.6 1.19999999 \\
-2.59999999 1.2 \\
-2.59999999 1.3 \\
-2.5 1.39999999 \\
-2.49999999 1.4 \\
-2.4 1.49999999 \\
-2.39999999 1.5 \\
-2.39999999 1.6 \\
-2.3 1.69999999 \\
-2.29999999 1.7 \\
-2.2 1.79999999 \\
-2.1 1.79999999 \\
-2.09999999 1.8 \\
-2 1.89999999 \\
-1.99999999 1.9 \\
-1.9 1.99999999 \\
-1.8 1.99999999 \\
-1.79999999 2 \\
-1.7 2.09999999 \\
-1.6 2.09999999 \\
-1.59999999 2.1 \\
-1.5 2.19999999 \\
-1.49999999 2.2 \\
-1.4 2.29999999 \\
-1.3 2.29999999 \\
-1.29999999 2.3 \\
-1.2 2.39999999 \\
-1.19999999 2.4 \\
-1.1 2.49999999 \\
-1.09999999 2.5 \\
-1 2.59999999 \\
-0.99999999 2.6 \\
-0.9 2.69999999 \\
-0.89999999 2.7 \\
-0.8 2.79999999 \\
-0.79999999 2.8 \\
-0.7 2.89999999 \\
-0.6 2.89999999 \\
-0.59999999 2.9 \\
-0.5 2.99999999 \\
-0.4 2.99999999 \\
-0.3 2.99999999 \\
-0.2 2.99999999 \\
-0.10000001 2.9 \\
-0.1 2.89999999 \\
0 2.89999999 \\
0.0999999899999997 2.8 \\
0.0999999999999999 2.79999999 \\
0.19999999 2.7 \\
0.19999999 2.6 \\
0.2 2.59999999 \\
0.29999999 2.5 \\
0.29999999 2.4 \\
0.29999999 2.3 \\
0.29999999 2.2 \\
0.29999999 2.1 \\
0.2 2.00000001 \\
0.19999999 2 \\
0.19999999 1.9 \\
0.19999999 1.8 \\
0.19999999 1.7 \\
0.0999999999999999 1.60000001 \\
0.0999999899999997 1.6 \\
0.0999999899999997 1.5 \\
0.0999999899999997 1.4 \\
0.0999999899999997 1.3 \\
0.0999999899999997 1.2 \\
0.0999999899999997 1.1 \\
0.0999999899999997 1 \\
0.0999999899999997 0.899999999999999 \\
0.0999999899999997 0.8 \\
0.0999999899999997 0.699999999999999 \\
0.0999999899999997 0.6 \\
0.0999999899999997 0.5 \\
0.0999999899999997 0.399999999999999 \\
0.0999999899999997 0.3 \\
0.0999999899999997 0.199999999999999 \\
0.0999999899999997 0.0999999999999996 \\
0 1.00000001168609e-008 \\
-1.0000000116861e-008 0 \\
0 -1.00000001168609e-008 \\
0.0999999899999997 -0.0999999999999996 \\
0.0999999899999997 -0.199999999999999 \\
0.0999999899999997 -0.3 \\
0.0999999899999997 -0.399999999999999 \\
0.0999999899999997 -0.5 \\
0.0999999899999997 -0.6 \\
0.0999999899999997 -0.699999999999999 \\
0.0999999899999997 -0.8 \\
0.0999999899999997 -0.899999999999999 \\
0.0999999899999997 -1 \\
0.0999999899999997 -1.1 \\
0.0999999899999997 -1.2 \\
0.0999999899999997 -1.3 \\
0.0999999899999997 -1.4 \\
0.0999999899999997 -1.5 \\
0.0999999899999997 -1.6 \\
0.0999999999999999 -1.60000001 \\
0.19999999 -1.7 \\
0.19999999 -1.8 \\
0.19999999 -1.9 \\
0.19999999 -2 \\
0.2 -2.00000001 \\
0.29999999 -2.1 \\
0.29999999 -2.2 \\
0.29999999 -2.3 \\
0.29999999 -2.4 \\
0.29999999 -2.5 \\
0.2 -2.59999999 \\
0.19999999 -2.6 \\
0.19999999 -2.7 \\
0.0999999999999999 -2.79999999 \\
0.0999999899999997 -2.8 \\
0 -2.89999999 \\
-0.1 -2.89999999 \\
-0.10000001 -2.9 \\
-0.2 -2.99999999 \\
-0.3 -2.99999999 \\
-0.4 -2.99999999 \\
-0.5 -2.99999999 \\
-0.59999999 -2.9 \\
-0.6 -2.89999999 \\
-0.7 -2.89999999 \\
-0.79999999 -2.8 \\
-0.8 -2.79999999 \\
-0.89999999 -2.7 \\
-0.9 -2.69999999 \\
-0.99999999 -2.6 \\
-1 -2.59999999 \\
-1.09999999 -2.5 \\
-1.1 -2.49999999 \\
-1.19999999 -2.4 \\
-1.2 -2.39999999 \\
-1.29999999 -2.3 \\
-1.3 -2.29999999 \\
-1.4 -2.29999999 \\
-1.49999999 -2.2 \\
-1.5 -2.19999999 \\
-1.59999999 -2.1 \\
-1.6 -2.09999999 \\
-1.7 -2.09999999 \\
-1.79999999 -2 \\
-1.8 -1.99999999 \\
-1.9 -1.99999999 \\
-1.99999999 -1.9 \\
-2 -1.89999999 \\
-2.09999999 -1.8 \\
-2.1 -1.79999999 \\
-2.2 -1.79999999 \\
-2.29999999 -1.7 \\
-2.3 -1.69999999 \\
-2.39999999 -1.6 \\
-2.39999999 -1.5 \\
-2.4 -1.49999999 \\
-2.49999999 -1.4 \\
-2.5 -1.39999999 \\
-2.59999999 -1.3 \\
-2.59999999 -1.2 \\
-2.6 -1.19999999 \\
-2.69999999 -1.1 \\
-2.69999999 -1 \\
-2.69999999 -0.899999999999999 \\
-2.7 -0.899999989999999 \\
-2.79999999 -0.8 \\
-2.79999999 -0.699999999999999 \\
-2.79999999 -0.6 \\
-2.79999999 -0.5 \\
-2.79999999 -0.399999999999999 \\
-2.79999999 -0.3 \\
-2.79999999 -0.199999999999999 \\
-2.79999999 -0.0999999999999996 \\
-2.79999999 0 \\
-2.79999999 0.0999999999999996 \\
-2.79999999 0.199999999999999 \\
-2.79999999 0.3 \\
-2.79999999 0.399999999999999 \\
-2.79999999 0.5 \\
-2.79999999 0.6 \\
-2.79999999 0.699999999999999 \\
-2.79999999 0.8 \\
-2.7 0.89999999 \\
NaN NaN \\
};


\addplot[dashed,line width=2.0pt,draw=mycolor1,forget plot]
table[row sep=crcr]{
x y\\
-2.7 0.89090909090909 \\
-2.69090909090909 0.899999999999999 \\
-2.69090909090909 1 \\
-2.69090909090909 1.1 \\
-2.6 1.19090909090909 \\
-2.59090909090909 1.2 \\
-2.59090909090909 1.3 \\
-2.5 1.39090909090909 \\
-2.49090909090909 1.4 \\
-2.4 1.49090909090909 \\
-2.39090909090909 1.5 \\
-2.39090909090909 1.6 \\
-2.3 1.69090909090909 \\
-2.29090909090909 1.7 \\
-2.2 1.79090909090909 \\
-2.1 1.79090909090909 \\
-2.09090909090909 1.8 \\
-2 1.89090909090909 \\
-1.99090909090909 1.9 \\
-1.9 1.99090909090909 \\
-1.8 1.99090909090909 \\
-1.79090909090909 2 \\
-1.7 2.09090909090909 \\
-1.6 2.09090909090909 \\
-1.59090909090909 2.1 \\
-1.5 2.19090909090909 \\
-1.49090909090909 2.2 \\
-1.4 2.29090909090909 \\
-1.3 2.29090909090909 \\
-1.29090909090909 2.3 \\
-1.2 2.39090909090909 \\
-1.19090909090909 2.4 \\
-1.1 2.49090909090909 \\
-1.09090909090909 2.5 \\
-1 2.59090909090909 \\
-0.990909090909091 2.6 \\
-0.9 2.69090909090909 \\
-0.890909090909091 2.7 \\
-0.8 2.79090909090909 \\
-0.790909090909091 2.8 \\
-0.7 2.89090909090909 \\
-0.6 2.89090909090909 \\
-0.590909090909091 2.9 \\
-0.5 2.99090909090909 \\
-0.4 2.99090909090909 \\
-0.3 2.99090909090909 \\
-0.2 2.99090909090909 \\
-0.109090909090909 2.9 \\
-0.1 2.89090909090909 \\
0 2.89090909090909 \\
0.0909090909090905 2.8 \\
0.0999999999999999 2.79090909090909 \\
0.190909090909091 2.7 \\
0.190909090909091 2.6 \\
0.2 2.59090909090909 \\
0.29090909090909 2.5 \\
0.29090909090909 2.4 \\
0.29090909090909 2.3 \\
0.29090909090909 2.2 \\
0.29090909090909 2.1 \\
0.2 2.00909090909091 \\
0.190909090909091 2 \\
0.190909090909091 1.9 \\
0.190909090909091 1.8 \\
0.190909090909091 1.7 \\
0.0999999999999999 1.60909090909091 \\
0.0909090909090905 1.6 \\
0.0909090909090905 1.5 \\
0.0909090909090905 1.4 \\
0.0909090909090905 1.3 \\
0.0909090909090905 1.2 \\
0.0909090909090905 1.1 \\
0.0909090909090905 1 \\
0.0909090909090905 0.899999999999999 \\
0.0909090909090905 0.8 \\
0.0909090909090905 0.699999999999999 \\
0.0909090909090905 0.6 \\
0.0909090909090905 0.5 \\
0.0909090909090905 0.399999999999999 \\
0.0909090909090905 0.3 \\
0.0909090909090905 0.199999999999999 \\
0.0909090909090905 0.0999999999999996 \\
0 0.00909090909090932 \\
-0.00909090909090936 0 \\
0 -0.00909090909090932 \\
0.0909090909090905 -0.0999999999999996 \\
0.0909090909090905 -0.199999999999999 \\
0.0909090909090905 -0.3 \\
0.0909090909090905 -0.399999999999999 \\
0.0909090909090905 -0.5 \\
0.0909090909090905 -0.6 \\
0.0909090909090905 -0.699999999999999 \\
0.0909090909090905 -0.8 \\
0.0909090909090905 -0.899999999999999 \\
0.0909090909090905 -1 \\
0.0909090909090905 -1.1 \\
0.0909090909090905 -1.2 \\
0.0909090909090905 -1.3 \\
0.0909090909090905 -1.4 \\
0.0909090909090905 -1.5 \\
0.0909090909090905 -1.6 \\
0.0999999999999999 -1.60909090909091 \\
0.190909090909091 -1.7 \\
0.190909090909091 -1.8 \\
0.190909090909091 -1.9 \\
0.190909090909091 -2 \\
0.2 -2.00909090909091 \\
0.29090909090909 -2.1 \\
0.29090909090909 -2.2 \\
0.29090909090909 -2.3 \\
0.29090909090909 -2.4 \\
0.29090909090909 -2.5 \\
0.2 -2.59090909090909 \\
0.190909090909091 -2.6 \\
0.190909090909091 -2.7 \\
0.0999999999999999 -2.79090909090909 \\
0.0909090909090905 -2.8 \\
0 -2.89090909090909 \\
-0.1 -2.89090909090909 \\
-0.109090909090909 -2.9 \\
-0.2 -2.99090909090909 \\
-0.3 -2.99090909090909 \\
-0.4 -2.99090909090909 \\
-0.5 -2.99090909090909 \\
-0.590909090909091 -2.9 \\
-0.6 -2.89090909090909 \\
-0.7 -2.89090909090909 \\
-0.790909090909091 -2.8 \\
-0.8 -2.79090909090909 \\
-0.890909090909091 -2.7 \\
-0.9 -2.69090909090909 \\
-0.990909090909091 -2.6 \\
-1 -2.59090909090909 \\
-1.09090909090909 -2.5 \\
-1.1 -2.49090909090909 \\
-1.19090909090909 -2.4 \\
-1.2 -2.39090909090909 \\
-1.29090909090909 -2.3 \\
-1.3 -2.29090909090909 \\
-1.4 -2.29090909090909 \\
-1.49090909090909 -2.2 \\
-1.5 -2.19090909090909 \\
-1.59090909090909 -2.1 \\
-1.6 -2.09090909090909 \\
-1.7 -2.09090909090909 \\
-1.79090909090909 -2 \\
-1.8 -1.99090909090909 \\
-1.9 -1.99090909090909 \\
-1.99090909090909 -1.9 \\
-2 -1.89090909090909 \\
-2.09090909090909 -1.8 \\
-2.1 -1.79090909090909 \\
-2.2 -1.79090909090909 \\
-2.29090909090909 -1.7 \\
-2.3 -1.69090909090909 \\
-2.39090909090909 -1.6 \\
-2.39090909090909 -1.5 \\
-2.4 -1.49090909090909 \\
-2.49090909090909 -1.4 \\
-2.5 -1.39090909090909 \\
-2.59090909090909 -1.3 \\
-2.59090909090909 -1.2 \\
-2.6 -1.19090909090909 \\
-2.69090909090909 -1.1 \\
-2.69090909090909 -1 \\
-2.69090909090909 -0.899999999999999 \\
-2.7 -0.89090909090909 \\
-2.79090909090909 -0.8 \\
-2.79090909090909 -0.699999999999999 \\
-2.79090909090909 -0.6 \\
-2.79090909090909 -0.5 \\
-2.79090909090909 -0.399999999999999 \\
-2.79090909090909 -0.3 \\
-2.79090909090909 -0.199999999999999 \\
-2.79090909090909 -0.0999999999999996 \\
-2.79090909090909 0 \\
-2.79090909090909 0.0999999999999996 \\
-2.79090909090909 0.199999999999999 \\
-2.79090909090909 0.3 \\
-2.79090909090909 0.399999999999999 \\
-2.79090909090909 0.5 \\
-2.79090909090909 0.6 \\
-2.79090909090909 0.699999999999999 \\
-2.79090909090909 0.8 \\
-2.7 0.89090909090909 \\
NaN NaN \\
};


\addplot[dashed,line width=2.0pt,draw=mycolor1,forget plot]
table[row sep=crcr]{
x y\\
-2.7 0.881818181818181 \\
-2.68181818181818 0.899999999999999 \\
-2.68181818181818 1 \\
-2.68181818181818 1.1 \\
-2.6 1.18181818181818 \\
-2.58181818181818 1.2 \\
-2.58181818181818 1.3 \\
-2.5 1.38181818181818 \\
-2.48181818181818 1.4 \\
-2.4 1.48181818181818 \\
-2.38181818181818 1.5 \\
-2.38181818181818 1.6 \\
-2.3 1.68181818181818 \\
-2.28181818181818 1.7 \\
-2.2 1.78181818181818 \\
-2.1 1.78181818181818 \\
-2.08181818181818 1.8 \\
-2 1.88181818181818 \\
-1.98181818181818 1.9 \\
-1.9 1.98181818181818 \\
-1.8 1.98181818181818 \\
-1.78181818181818 2 \\
-1.7 2.08181818181818 \\
-1.6 2.08181818181818 \\
-1.58181818181818 2.1 \\
-1.5 2.18181818181818 \\
-1.48181818181818 2.2 \\
-1.4 2.28181818181818 \\
-1.3 2.28181818181818 \\
-1.28181818181818 2.3 \\
-1.2 2.38181818181818 \\
-1.18181818181818 2.4 \\
-1.1 2.48181818181818 \\
-1.08181818181818 2.5 \\
-1 2.58181818181818 \\
-0.981818181818182 2.6 \\
-0.9 2.68181818181818 \\
-0.881818181818182 2.7 \\
-0.8 2.78181818181818 \\
-0.781818181818182 2.8 \\
-0.7 2.88181818181818 \\
-0.6 2.88181818181818 \\
-0.581818181818182 2.9 \\
-0.5 2.98181818181818 \\
-0.4 2.98181818181818 \\
-0.3 2.98181818181818 \\
-0.2 2.98181818181818 \\
-0.118181818181818 2.9 \\
-0.1 2.88181818181818 \\
0 2.88181818181818 \\
0.0818181818181819 2.8 \\
0.0999999999999999 2.78181818181818 \\
0.181818181818182 2.7 \\
0.181818181818182 2.6 \\
0.2 2.58181818181818 \\
0.281818181818182 2.5 \\
0.281818181818182 2.4 \\
0.281818181818182 2.3 \\
0.281818181818182 2.2 \\
0.281818181818182 2.1 \\
0.2 2.01818181818182 \\
0.181818181818182 2 \\
0.181818181818182 1.9 \\
0.181818181818182 1.8 \\
0.181818181818182 1.7 \\
0.0999999999999999 1.61818181818182 \\
0.0818181818181819 1.6 \\
0.0818181818181819 1.5 \\
0.0818181818181819 1.4 \\
0.0818181818181819 1.3 \\
0.0818181818181819 1.2 \\
0.0818181818181819 1.1 \\
0.0818181818181819 1 \\
0.0818181818181819 0.899999999999999 \\
0.0818181818181819 0.8 \\
0.0818181818181819 0.699999999999999 \\
0.0818181818181819 0.6 \\
0.0818181818181819 0.5 \\
0.0818181818181819 0.399999999999999 \\
0.0818181818181819 0.3 \\
0.0818181818181819 0.199999999999999 \\
0.0818181818181819 0.0999999999999996 \\
0 0.0181818181818179 \\
-0.018181818181818 0 \\
0 -0.0181818181818179 \\
0.0818181818181819 -0.0999999999999996 \\
0.0818181818181819 -0.199999999999999 \\
0.0818181818181819 -0.3 \\
0.0818181818181819 -0.399999999999999 \\
0.0818181818181819 -0.5 \\
0.0818181818181819 -0.6 \\
0.0818181818181819 -0.699999999999999 \\
0.0818181818181819 -0.8 \\
0.0818181818181819 -0.899999999999999 \\
0.0818181818181819 -1 \\
0.0818181818181819 -1.1 \\
0.0818181818181819 -1.2 \\
0.0818181818181819 -1.3 \\
0.0818181818181819 -1.4 \\
0.0818181818181819 -1.5 \\
0.0818181818181819 -1.6 \\
0.0999999999999999 -1.61818181818182 \\
0.181818181818182 -1.7 \\
0.181818181818182 -1.8 \\
0.181818181818182 -1.9 \\
0.181818181818182 -2 \\
0.2 -2.01818181818182 \\
0.281818181818182 -2.1 \\
0.281818181818182 -2.2 \\
0.281818181818182 -2.3 \\
0.281818181818182 -2.4 \\
0.281818181818182 -2.5 \\
0.2 -2.58181818181818 \\
0.181818181818182 -2.6 \\
0.181818181818182 -2.7 \\
0.0999999999999999 -2.78181818181818 \\
0.0818181818181819 -2.8 \\
0 -2.88181818181818 \\
-0.1 -2.88181818181818 \\
-0.118181818181818 -2.9 \\
-0.2 -2.98181818181818 \\
-0.3 -2.98181818181818 \\
-0.4 -2.98181818181818 \\
-0.5 -2.98181818181818 \\
-0.581818181818182 -2.9 \\
-0.6 -2.88181818181818 \\
-0.7 -2.88181818181818 \\
-0.781818181818182 -2.8 \\
-0.8 -2.78181818181818 \\
-0.881818181818182 -2.7 \\
-0.9 -2.68181818181818 \\
-0.981818181818182 -2.6 \\
-1 -2.58181818181818 \\
-1.08181818181818 -2.5 \\
-1.1 -2.48181818181818 \\
-1.18181818181818 -2.4 \\
-1.2 -2.38181818181818 \\
-1.28181818181818 -2.3 \\
-1.3 -2.28181818181818 \\
-1.4 -2.28181818181818 \\
-1.48181818181818 -2.2 \\
-1.5 -2.18181818181818 \\
-1.58181818181818 -2.1 \\
-1.6 -2.08181818181818 \\
-1.7 -2.08181818181818 \\
-1.78181818181818 -2 \\
-1.8 -1.98181818181818 \\
-1.9 -1.98181818181818 \\
-1.98181818181818 -1.9 \\
-2 -1.88181818181818 \\
-2.08181818181818 -1.8 \\
-2.1 -1.78181818181818 \\
-2.2 -1.78181818181818 \\
-2.28181818181818 -1.7 \\
-2.3 -1.68181818181818 \\
-2.38181818181818 -1.6 \\
-2.38181818181818 -1.5 \\
-2.4 -1.48181818181818 \\
-2.48181818181818 -1.4 \\
-2.5 -1.38181818181818 \\
-2.58181818181818 -1.3 \\
-2.58181818181818 -1.2 \\
-2.6 -1.18181818181818 \\
-2.68181818181818 -1.1 \\
-2.68181818181818 -1 \\
-2.68181818181818 -0.899999999999999 \\
-2.7 -0.881818181818182 \\
-2.78181818181818 -0.8 \\
-2.78181818181818 -0.699999999999999 \\
-2.78181818181818 -0.6 \\
-2.78181818181818 -0.5 \\
-2.78181818181818 -0.399999999999999 \\
-2.78181818181818 -0.3 \\
-2.78181818181818 -0.199999999999999 \\
-2.78181818181818 -0.0999999999999996 \\
-2.78181818181818 0 \\
-2.78181818181818 0.0999999999999996 \\
-2.78181818181818 0.199999999999999 \\
-2.78181818181818 0.3 \\
-2.78181818181818 0.399999999999999 \\
-2.78181818181818 0.5 \\
-2.78181818181818 0.6 \\
-2.78181818181818 0.699999999999999 \\
-2.78181818181818 0.8 \\
-2.7 0.881818181818181 \\
NaN NaN \\
};


\addplot[dashed,line width=2.0pt,draw=mycolor1,forget plot]
table[row sep=crcr]{
x y\\
-2.7 0.872727272727273 \\
-2.67272727272727 0.899999999999999 \\
-2.67272727272727 1 \\
-2.67272727272727 1.1 \\
-2.6 1.17272727272727 \\
-2.57272727272727 1.2 \\
-2.57272727272727 1.3 \\
-2.5 1.37272727272727 \\
-2.47272727272727 1.4 \\
-2.4 1.47272727272727 \\
-2.37272727272727 1.5 \\
-2.37272727272727 1.6 \\
-2.3 1.67272727272727 \\
-2.27272727272727 1.7 \\
-2.2 1.77272727272727 \\
-2.1 1.77272727272727 \\
-2.07272727272727 1.8 \\
-2 1.87272727272727 \\
-1.97272727272727 1.9 \\
-1.9 1.97272727272727 \\
-1.8 1.97272727272727 \\
-1.77272727272727 2 \\
-1.7 2.07272727272727 \\
-1.6 2.07272727272727 \\
-1.57272727272727 2.1 \\
-1.5 2.17272727272727 \\
-1.47272727272727 2.2 \\
-1.4 2.27272727272727 \\
-1.3 2.27272727272727 \\
-1.27272727272727 2.3 \\
-1.2 2.37272727272727 \\
-1.17272727272727 2.4 \\
-1.1 2.47272727272727 \\
-1.07272727272727 2.5 \\
-1 2.57272727272727 \\
-0.972727272727273 2.6 \\
-0.9 2.67272727272727 \\
-0.872727272727273 2.7 \\
-0.8 2.77272727272727 \\
-0.772727272727273 2.8 \\
-0.7 2.87272727272727 \\
-0.6 2.87272727272727 \\
-0.572727272727273 2.9 \\
-0.5 2.97272727272727 \\
-0.4 2.97272727272727 \\
-0.3 2.97272727272727 \\
-0.2 2.97272727272727 \\
-0.127272727272727 2.9 \\
-0.1 2.87272727272727 \\
0 2.87272727272727 \\
0.0727272727272726 2.8 \\
0.0999999999999999 2.77272727272727 \\
0.172727272727273 2.7 \\
0.172727272727273 2.6 \\
0.2 2.57272727272727 \\
0.272727272727273 2.5 \\
0.272727272727273 2.4 \\
0.272727272727273 2.3 \\
0.272727272727273 2.2 \\
0.272727272727273 2.1 \\
0.2 2.02727272727273 \\
0.172727272727273 2 \\
0.172727272727273 1.9 \\
0.172727272727273 1.8 \\
0.172727272727273 1.7 \\
0.0999999999999999 1.62727272727273 \\
0.0727272727272726 1.6 \\
0.0727272727272726 1.5 \\
0.0727272727272726 1.4 \\
0.0727272727272726 1.3 \\
0.0727272727272726 1.2 \\
0.0727272727272726 1.1 \\
0.0727272727272726 1 \\
0.0727272727272726 0.899999999999999 \\
0.0727272727272726 0.8 \\
0.0727272727272726 0.699999999999999 \\
0.0727272727272726 0.6 \\
0.0727272727272726 0.5 \\
0.0727272727272726 0.399999999999999 \\
0.0727272727272726 0.3 \\
0.0727272727272726 0.199999999999999 \\
0.0727272727272726 0.0999999999999996 \\
0 0.0272727272727272 \\
-0.0272727272727274 0 \\
0 -0.0272727272727272 \\
0.0727272727272726 -0.0999999999999996 \\
0.0727272727272726 -0.199999999999999 \\
0.0727272727272726 -0.3 \\
0.0727272727272726 -0.399999999999999 \\
0.0727272727272726 -0.5 \\
0.0727272727272726 -0.6 \\
0.0727272727272726 -0.699999999999999 \\
0.0727272727272726 -0.8 \\
0.0727272727272726 -0.899999999999999 \\
0.0727272727272726 -1 \\
0.0727272727272726 -1.1 \\
0.0727272727272726 -1.2 \\
0.0727272727272726 -1.3 \\
0.0727272727272726 -1.4 \\
0.0727272727272726 -1.5 \\
0.0727272727272726 -1.6 \\
0.0999999999999999 -1.62727272727273 \\
0.172727272727273 -1.7 \\
0.172727272727273 -1.8 \\
0.172727272727273 -1.9 \\
0.172727272727273 -2 \\
0.2 -2.02727272727273 \\
0.272727272727273 -2.1 \\
0.272727272727273 -2.2 \\
0.272727272727273 -2.3 \\
0.272727272727273 -2.4 \\
0.272727272727273 -2.5 \\
0.2 -2.57272727272727 \\
0.172727272727273 -2.6 \\
0.172727272727273 -2.7 \\
0.0999999999999999 -2.77272727272727 \\
0.0727272727272726 -2.8 \\
0 -2.87272727272727 \\
-0.1 -2.87272727272727 \\
-0.127272727272727 -2.9 \\
-0.2 -2.97272727272727 \\
-0.3 -2.97272727272727 \\
-0.4 -2.97272727272727 \\
-0.5 -2.97272727272727 \\
-0.572727272727273 -2.9 \\
-0.6 -2.87272727272727 \\
-0.7 -2.87272727272727 \\
-0.772727272727273 -2.8 \\
-0.8 -2.77272727272727 \\
-0.872727272727273 -2.7 \\
-0.9 -2.67272727272727 \\
-0.972727272727273 -2.6 \\
-1 -2.57272727272727 \\
-1.07272727272727 -2.5 \\
-1.1 -2.47272727272727 \\
-1.17272727272727 -2.4 \\
-1.2 -2.37272727272727 \\
-1.27272727272727 -2.3 \\
-1.3 -2.27272727272727 \\
-1.4 -2.27272727272727 \\
-1.47272727272727 -2.2 \\
-1.5 -2.17272727272727 \\
-1.57272727272727 -2.1 \\
-1.6 -2.07272727272727 \\
-1.7 -2.07272727272727 \\
-1.77272727272727 -2 \\
-1.8 -1.97272727272727 \\
-1.9 -1.97272727272727 \\
-1.97272727272727 -1.9 \\
-2 -1.87272727272727 \\
-2.07272727272727 -1.8 \\
-2.1 -1.77272727272727 \\
-2.2 -1.77272727272727 \\
-2.27272727272727 -1.7 \\
-2.3 -1.67272727272727 \\
-2.37272727272727 -1.6 \\
-2.37272727272727 -1.5 \\
-2.4 -1.47272727272727 \\
-2.47272727272727 -1.4 \\
-2.5 -1.37272727272727 \\
-2.57272727272727 -1.3 \\
-2.57272727272727 -1.2 \\
-2.6 -1.17272727272727 \\
-2.67272727272727 -1.1 \\
-2.67272727272727 -1 \\
-2.67272727272727 -0.899999999999999 \\
-2.7 -0.872727272727272 \\
-2.77272727272727 -0.8 \\
-2.77272727272727 -0.699999999999999 \\
-2.77272727272727 -0.6 \\
-2.77272727272727 -0.5 \\
-2.77272727272727 -0.399999999999999 \\
-2.77272727272727 -0.3 \\
-2.77272727272727 -0.199999999999999 \\
-2.77272727272727 -0.0999999999999996 \\
-2.77272727272727 0 \\
-2.77272727272727 0.0999999999999996 \\
-2.77272727272727 0.199999999999999 \\
-2.77272727272727 0.3 \\
-2.77272727272727 0.399999999999999 \\
-2.77272727272727 0.5 \\
-2.77272727272727 0.6 \\
-2.77272727272727 0.699999999999999 \\
-2.77272727272727 0.8 \\
-2.7 0.872727272727273 \\
NaN NaN \\
};


\addplot[dashed,line width=2.0pt,draw=mycolor1,forget plot]
table[row sep=crcr]{
x y\\
-2.7 0.863636363636364 \\
-2.66363636363636 0.899999999999999 \\
-2.66363636363636 1 \\
-2.66363636363636 1.1 \\
-2.6 1.16363636363636 \\
-2.56363636363636 1.2 \\
-2.56363636363636 1.3 \\
-2.5 1.36363636363636 \\
-2.46363636363636 1.4 \\
-2.4 1.46363636363636 \\
-2.36363636363636 1.5 \\
-2.36363636363636 1.6 \\
-2.3 1.66363636363636 \\
-2.26363636363636 1.7 \\
-2.2 1.76363636363636 \\
-2.1 1.76363636363636 \\
-2.06363636363636 1.8 \\
-2 1.86363636363636 \\
-1.96363636363636 1.9 \\
-1.9 1.96363636363636 \\
-1.8 1.96363636363636 \\
-1.76363636363636 2 \\
-1.7 2.06363636363636 \\
-1.6 2.06363636363636 \\
-1.56363636363636 2.1 \\
-1.5 2.16363636363636 \\
-1.46363636363636 2.2 \\
-1.4 2.26363636363636 \\
-1.3 2.26363636363636 \\
-1.26363636363636 2.3 \\
-1.2 2.36363636363636 \\
-1.16363636363636 2.4 \\
-1.1 2.46363636363636 \\
-1.06363636363636 2.5 \\
-1 2.56363636363636 \\
-0.963636363636364 2.6 \\
-0.9 2.66363636363636 \\
-0.863636363636364 2.7 \\
-0.8 2.76363636363636 \\
-0.763636363636364 2.8 \\
-0.7 2.86363636363636 \\
-0.6 2.86363636363636 \\
-0.563636363636363 2.9 \\
-0.5 2.96363636363636 \\
-0.4 2.96363636363636 \\
-0.3 2.96363636363636 \\
-0.2 2.96363636363636 \\
-0.136363636363637 2.9 \\
-0.1 2.86363636363636 \\
0 2.86363636363636 \\
0.0636363636363632 2.8 \\
0.0999999999999999 2.76363636363636 \\
0.163636363636363 2.7 \\
0.163636363636363 2.6 \\
0.2 2.56363636363636 \\
0.263636363636363 2.5 \\
0.263636363636363 2.4 \\
0.263636363636363 2.3 \\
0.263636363636363 2.2 \\
0.263636363636363 2.1 \\
0.2 2.03636363636364 \\
0.163636363636363 2 \\
0.163636363636363 1.9 \\
0.163636363636363 1.8 \\
0.163636363636363 1.7 \\
0.0999999999999999 1.63636363636364 \\
0.0636363636363632 1.6 \\
0.0636363636363632 1.5 \\
0.0636363636363632 1.4 \\
0.0636363636363632 1.3 \\
0.0636363636363632 1.2 \\
0.0636363636363632 1.1 \\
0.0636363636363632 1 \\
0.0636363636363632 0.899999999999999 \\
0.0636363636363632 0.8 \\
0.0636363636363632 0.699999999999999 \\
0.0636363636363632 0.6 \\
0.0636363636363632 0.5 \\
0.0636363636363632 0.399999999999999 \\
0.0636363636363632 0.3 \\
0.0636363636363632 0.199999999999999 \\
0.0636363636363632 0.0999999999999996 \\
0 0.0363636363636366 \\
-0.0363636363636367 0 \\
0 -0.0363636363636366 \\
0.0636363636363632 -0.0999999999999996 \\
0.0636363636363632 -0.199999999999999 \\
0.0636363636363632 -0.3 \\
0.0636363636363632 -0.399999999999999 \\
0.0636363636363632 -0.5 \\
0.0636363636363632 -0.6 \\
0.0636363636363632 -0.699999999999999 \\
0.0636363636363632 -0.8 \\
0.0636363636363632 -0.899999999999999 \\
0.0636363636363632 -1 \\
0.0636363636363632 -1.1 \\
0.0636363636363632 -1.2 \\
0.0636363636363632 -1.3 \\
0.0636363636363632 -1.4 \\
0.0636363636363632 -1.5 \\
0.0636363636363632 -1.6 \\
0.0999999999999999 -1.63636363636364 \\
0.163636363636363 -1.7 \\
0.163636363636363 -1.8 \\
0.163636363636363 -1.9 \\
0.163636363636363 -2 \\
0.2 -2.03636363636364 \\
0.263636363636363 -2.1 \\
0.263636363636363 -2.2 \\
0.263636363636363 -2.3 \\
0.263636363636363 -2.4 \\
0.263636363636363 -2.5 \\
0.2 -2.56363636363636 \\
0.163636363636363 -2.6 \\
0.163636363636363 -2.7 \\
0.0999999999999999 -2.76363636363636 \\
0.0636363636363632 -2.8 \\
0 -2.86363636363636 \\
-0.1 -2.86363636363636 \\
-0.136363636363637 -2.9 \\
-0.2 -2.96363636363636 \\
-0.3 -2.96363636363636 \\
-0.4 -2.96363636363636 \\
-0.5 -2.96363636363636 \\
-0.563636363636363 -2.9 \\
-0.6 -2.86363636363636 \\
-0.7 -2.86363636363636 \\
-0.763636363636364 -2.8 \\
-0.8 -2.76363636363636 \\
-0.863636363636364 -2.7 \\
-0.9 -2.66363636363636 \\
-0.963636363636364 -2.6 \\
-1 -2.56363636363636 \\
-1.06363636363636 -2.5 \\
-1.1 -2.46363636363636 \\
-1.16363636363636 -2.4 \\
-1.2 -2.36363636363636 \\
-1.26363636363636 -2.3 \\
-1.3 -2.26363636363636 \\
-1.4 -2.26363636363636 \\
-1.46363636363636 -2.2 \\
-1.5 -2.16363636363636 \\
-1.56363636363636 -2.1 \\
-1.6 -2.06363636363636 \\
-1.7 -2.06363636363636 \\
-1.76363636363636 -2 \\
-1.8 -1.96363636363636 \\
-1.9 -1.96363636363636 \\
-1.96363636363636 -1.9 \\
-2 -1.86363636363636 \\
-2.06363636363636 -1.8 \\
-2.1 -1.76363636363636 \\
-2.2 -1.76363636363636 \\
-2.26363636363636 -1.7 \\
-2.3 -1.66363636363636 \\
-2.36363636363636 -1.6 \\
-2.36363636363636 -1.5 \\
-2.4 -1.46363636363636 \\
-2.46363636363636 -1.4 \\
-2.5 -1.36363636363636 \\
-2.56363636363636 -1.3 \\
-2.56363636363636 -1.2 \\
-2.6 -1.16363636363636 \\
-2.66363636363636 -1.1 \\
-2.66363636363636 -1 \\
-2.66363636363636 -0.899999999999999 \\
-2.7 -0.863636363636363 \\
-2.76363636363636 -0.8 \\
-2.76363636363636 -0.699999999999999 \\
-2.76363636363636 -0.6 \\
-2.76363636363636 -0.5 \\
-2.76363636363636 -0.399999999999999 \\
-2.76363636363636 -0.3 \\
-2.76363636363636 -0.199999999999999 \\
-2.76363636363636 -0.0999999999999996 \\
-2.76363636363636 0 \\
-2.76363636363636 0.0999999999999996 \\
-2.76363636363636 0.199999999999999 \\
-2.76363636363636 0.3 \\
-2.76363636363636 0.399999999999999 \\
-2.76363636363636 0.5 \\
-2.76363636363636 0.6 \\
-2.76363636363636 0.699999999999999 \\
-2.76363636363636 0.8 \\
-2.7 0.863636363636364 \\
NaN NaN \\
};


\addplot[dashed,line width=2.0pt,draw=mycolor1,forget plot]
table[row sep=crcr]{
x y\\
-2.7 0.854545454545454 \\
-2.65454545454545 0.899999999999999 \\
-2.65454545454545 1 \\
-2.65454545454545 1.1 \\
-2.6 1.15454545454545 \\
-2.55454545454545 1.2 \\
-2.55454545454545 1.3 \\
-2.5 1.35454545454545 \\
-2.45454545454545 1.4 \\
-2.4 1.45454545454545 \\
-2.35454545454545 1.5 \\
-2.35454545454545 1.6 \\
-2.3 1.65454545454545 \\
-2.25454545454545 1.7 \\
-2.2 1.75454545454545 \\
-2.1 1.75454545454545 \\
-2.05454545454545 1.8 \\
-2 1.85454545454545 \\
-1.95454545454545 1.9 \\
-1.9 1.95454545454545 \\
-1.8 1.95454545454545 \\
-1.75454545454545 2 \\
-1.7 2.05454545454545 \\
-1.6 2.05454545454545 \\
-1.55454545454545 2.1 \\
-1.5 2.15454545454545 \\
-1.45454545454545 2.2 \\
-1.4 2.25454545454545 \\
-1.3 2.25454545454545 \\
-1.25454545454545 2.3 \\
-1.2 2.35454545454545 \\
-1.15454545454545 2.4 \\
-1.1 2.45454545454545 \\
-1.05454545454545 2.5 \\
-1 2.55454545454545 \\
-0.954545454545455 2.6 \\
-0.9 2.65454545454545 \\
-0.854545454545455 2.7 \\
-0.8 2.75454545454545 \\
-0.754545454545455 2.8 \\
-0.7 2.85454545454545 \\
-0.6 2.85454545454545 \\
-0.554545454545455 2.9 \\
-0.5 2.95454545454545 \\
-0.4 2.95454545454545 \\
-0.3 2.95454545454545 \\
-0.2 2.95454545454545 \\
-0.145454545454545 2.9 \\
-0.1 2.85454545454545 \\
0 2.85454545454545 \\
0.0545454545454546 2.8 \\
0.0999999999999999 2.75454545454545 \\
0.154545454545455 2.7 \\
0.154545454545455 2.6 \\
0.2 2.55454545454545 \\
0.254545454545455 2.5 \\
0.254545454545455 2.4 \\
0.254545454545455 2.3 \\
0.254545454545455 2.2 \\
0.254545454545455 2.1 \\
0.2 2.04545454545455 \\
0.154545454545455 2 \\
0.154545454545455 1.9 \\
0.154545454545455 1.8 \\
0.154545454545455 1.7 \\
0.0999999999999999 1.64545454545455 \\
0.0545454545454546 1.6 \\
0.0545454545454546 1.5 \\
0.0545454545454546 1.4 \\
0.0545454545454546 1.3 \\
0.0545454545454546 1.2 \\
0.0545454545454546 1.1 \\
0.0545454545454546 1 \\
0.0545454545454546 0.899999999999999 \\
0.0545454545454546 0.8 \\
0.0545454545454546 0.699999999999999 \\
0.0545454545454546 0.6 \\
0.0545454545454546 0.5 \\
0.0545454545454546 0.399999999999999 \\
0.0545454545454546 0.3 \\
0.0545454545454546 0.199999999999999 \\
0.0545454545454546 0.0999999999999996 \\
0 0.0454545454545452 \\
-0.0454545454545454 0 \\
0 -0.0454545454545452 \\
0.0545454545454546 -0.0999999999999996 \\
0.0545454545454546 -0.199999999999999 \\
0.0545454545454546 -0.3 \\
0.0545454545454546 -0.399999999999999 \\
0.0545454545454546 -0.5 \\
0.0545454545454546 -0.6 \\
0.0545454545454546 -0.699999999999999 \\
0.0545454545454546 -0.8 \\
0.0545454545454546 -0.899999999999999 \\
0.0545454545454546 -1 \\
0.0545454545454546 -1.1 \\
0.0545454545454546 -1.2 \\
0.0545454545454546 -1.3 \\
0.0545454545454546 -1.4 \\
0.0545454545454546 -1.5 \\
0.0545454545454546 -1.6 \\
0.0999999999999999 -1.64545454545455 \\
0.154545454545455 -1.7 \\
0.154545454545455 -1.8 \\
0.154545454545455 -1.9 \\
0.154545454545455 -2 \\
0.2 -2.04545454545455 \\
0.254545454545455 -2.1 \\
0.254545454545455 -2.2 \\
0.254545454545455 -2.3 \\
0.254545454545455 -2.4 \\
0.254545454545455 -2.5 \\
0.2 -2.55454545454545 \\
0.154545454545455 -2.6 \\
0.154545454545455 -2.7 \\
0.0999999999999999 -2.75454545454545 \\
0.0545454545454546 -2.8 \\
0 -2.85454545454545 \\
-0.1 -2.85454545454545 \\
-0.145454545454545 -2.9 \\
-0.2 -2.95454545454545 \\
-0.3 -2.95454545454545 \\
-0.4 -2.95454545454545 \\
-0.5 -2.95454545454545 \\
-0.554545454545455 -2.9 \\
-0.6 -2.85454545454545 \\
-0.7 -2.85454545454545 \\
-0.754545454545455 -2.8 \\
-0.8 -2.75454545454545 \\
-0.854545454545455 -2.7 \\
-0.9 -2.65454545454545 \\
-0.954545454545455 -2.6 \\
-1 -2.55454545454545 \\
-1.05454545454545 -2.5 \\
-1.1 -2.45454545454545 \\
-1.15454545454545 -2.4 \\
-1.2 -2.35454545454545 \\
-1.25454545454545 -2.3 \\
-1.3 -2.25454545454545 \\
-1.4 -2.25454545454545 \\
-1.45454545454545 -2.2 \\
-1.5 -2.15454545454545 \\
-1.55454545454545 -2.1 \\
-1.6 -2.05454545454545 \\
-1.7 -2.05454545454545 \\
-1.75454545454545 -2 \\
-1.8 -1.95454545454545 \\
-1.9 -1.95454545454545 \\
-1.95454545454545 -1.9 \\
-2 -1.85454545454545 \\
-2.05454545454545 -1.8 \\
-2.1 -1.75454545454545 \\
-2.2 -1.75454545454545 \\
-2.25454545454545 -1.7 \\
-2.3 -1.65454545454545 \\
-2.35454545454545 -1.6 \\
-2.35454545454545 -1.5 \\
-2.4 -1.45454545454545 \\
-2.45454545454545 -1.4 \\
-2.5 -1.35454545454545 \\
-2.55454545454545 -1.3 \\
-2.55454545454545 -1.2 \\
-2.6 -1.15454545454545 \\
-2.65454545454545 -1.1 \\
-2.65454545454545 -1 \\
-2.65454545454545 -0.899999999999999 \\
-2.7 -0.854545454545454 \\
-2.75454545454545 -0.8 \\
-2.75454545454545 -0.699999999999999 \\
-2.75454545454545 -0.6 \\
-2.75454545454545 -0.5 \\
-2.75454545454545 -0.399999999999999 \\
-2.75454545454545 -0.3 \\
-2.75454545454545 -0.199999999999999 \\
-2.75454545454545 -0.0999999999999996 \\
-2.75454545454545 0 \\
-2.75454545454545 0.0999999999999996 \\
-2.75454545454545 0.199999999999999 \\
-2.75454545454545 0.3 \\
-2.75454545454545 0.399999999999999 \\
-2.75454545454545 0.5 \\
-2.75454545454545 0.6 \\
-2.75454545454545 0.699999999999999 \\
-2.75454545454545 0.8 \\
-2.7 0.854545454545454 \\
NaN NaN \\
};


\addplot[dashed,line width=2.0pt,draw=mycolor1,forget plot]
table[row sep=crcr]{
x y\\
-2.7 0.845454545454545 \\
-2.64545454545455 0.899999999999999 \\
-2.64545454545455 1 \\
-2.64545454545455 1.1 \\
-2.6 1.14545454545454 \\
-2.54545454545455 1.2 \\
-2.54545454545455 1.3 \\
-2.5 1.34545454545455 \\
-2.44545454545455 1.4 \\
-2.4 1.44545454545454 \\
-2.34545454545455 1.5 \\
-2.34545454545455 1.6 \\
-2.3 1.64545454545455 \\
-2.24545454545455 1.7 \\
-2.2 1.74545454545455 \\
-2.1 1.74545454545455 \\
-2.04545454545455 1.8 \\
-2 1.84545454545455 \\
-1.94545454545455 1.9 \\
-1.9 1.94545454545455 \\
-1.8 1.94545454545455 \\
-1.74545454545455 2 \\
-1.7 2.04545454545455 \\
-1.6 2.04545454545455 \\
-1.54545454545455 2.1 \\
-1.5 2.14545454545455 \\
-1.44545454545455 2.2 \\
-1.4 2.24545454545455 \\
-1.3 2.24545454545455 \\
-1.24545454545455 2.3 \\
-1.2 2.34545454545455 \\
-1.14545454545455 2.4 \\
-1.1 2.44545454545455 \\
-1.04545454545455 2.5 \\
-1 2.54545454545455 \\
-0.945454545454546 2.6 \\
-0.9 2.64545454545455 \\
-0.845454545454546 2.7 \\
-0.8 2.74545454545455 \\
-0.745454545454546 2.8 \\
-0.7 2.84545454545455 \\
-0.6 2.84545454545455 \\
-0.545454545454545 2.9 \\
-0.5 2.94545454545455 \\
-0.4 2.94545454545455 \\
-0.3 2.94545454545455 \\
-0.2 2.94545454545455 \\
-0.154545454545455 2.9 \\
-0.1 2.84545454545455 \\
0 2.84545454545455 \\
0.0454545454545453 2.8 \\
0.0999999999999999 2.74545454545455 \\
0.145454545454545 2.7 \\
0.145454545454545 2.6 \\
0.2 2.54545454545455 \\
0.245454545454545 2.5 \\
0.245454545454545 2.4 \\
0.245454545454545 2.3 \\
0.245454545454545 2.2 \\
0.245454545454545 2.1 \\
0.2 2.05454545454545 \\
0.145454545454545 2 \\
0.145454545454545 1.9 \\
0.145454545454545 1.8 \\
0.145454545454545 1.7 \\
0.0999999999999999 1.65454545454545 \\
0.0454545454545453 1.6 \\
0.0454545454545453 1.5 \\
0.0454545454545453 1.4 \\
0.0454545454545453 1.3 \\
0.0454545454545453 1.2 \\
0.0454545454545453 1.1 \\
0.0454545454545453 1 \\
0.0454545454545453 0.899999999999999 \\
0.0454545454545453 0.8 \\
0.0454545454545453 0.699999999999999 \\
0.0454545454545453 0.6 \\
0.0454545454545453 0.5 \\
0.0454545454545453 0.399999999999999 \\
0.0454545454545453 0.3 \\
0.0454545454545453 0.199999999999999 \\
0.0454545454545453 0.0999999999999996 \\
0 0.0545454545454545 \\
-0.0545454545454547 0 \\
0 -0.0545454545454545 \\
0.0454545454545453 -0.0999999999999996 \\
0.0454545454545453 -0.199999999999999 \\
0.0454545454545453 -0.3 \\
0.0454545454545453 -0.399999999999999 \\
0.0454545454545453 -0.5 \\
0.0454545454545453 -0.6 \\
0.0454545454545453 -0.699999999999999 \\
0.0454545454545453 -0.8 \\
0.0454545454545453 -0.899999999999999 \\
0.0454545454545453 -1 \\
0.0454545454545453 -1.1 \\
0.0454545454545453 -1.2 \\
0.0454545454545453 -1.3 \\
0.0454545454545453 -1.4 \\
0.0454545454545453 -1.5 \\
0.0454545454545453 -1.6 \\
0.0999999999999999 -1.65454545454545 \\
0.145454545454545 -1.7 \\
0.145454545454545 -1.8 \\
0.145454545454545 -1.9 \\
0.145454545454545 -2 \\
0.2 -2.05454545454545 \\
0.245454545454545 -2.1 \\
0.245454545454545 -2.2 \\
0.245454545454545 -2.3 \\
0.245454545454545 -2.4 \\
0.245454545454545 -2.5 \\
0.2 -2.54545454545455 \\
0.145454545454545 -2.6 \\
0.145454545454545 -2.7 \\
0.0999999999999999 -2.74545454545455 \\
0.0454545454545453 -2.8 \\
0 -2.84545454545455 \\
-0.1 -2.84545454545455 \\
-0.154545454545455 -2.9 \\
-0.2 -2.94545454545455 \\
-0.3 -2.94545454545455 \\
-0.4 -2.94545454545455 \\
-0.5 -2.94545454545455 \\
-0.545454545454545 -2.9 \\
-0.6 -2.84545454545455 \\
-0.7 -2.84545454545455 \\
-0.745454545454546 -2.8 \\
-0.8 -2.74545454545455 \\
-0.845454545454546 -2.7 \\
-0.9 -2.64545454545455 \\
-0.945454545454546 -2.6 \\
-1 -2.54545454545455 \\
-1.04545454545455 -2.5 \\
-1.1 -2.44545454545455 \\
-1.14545454545455 -2.4 \\
-1.2 -2.34545454545455 \\
-1.24545454545455 -2.3 \\
-1.3 -2.24545454545455 \\
-1.4 -2.24545454545455 \\
-1.44545454545455 -2.2 \\
-1.5 -2.14545454545455 \\
-1.54545454545455 -2.1 \\
-1.6 -2.04545454545455 \\
-1.7 -2.04545454545455 \\
-1.74545454545455 -2 \\
-1.8 -1.94545454545455 \\
-1.9 -1.94545454545455 \\
-1.94545454545455 -1.9 \\
-2 -1.84545454545455 \\
-2.04545454545455 -1.8 \\
-2.1 -1.74545454545455 \\
-2.2 -1.74545454545455 \\
-2.24545454545455 -1.7 \\
-2.3 -1.64545454545455 \\
-2.34545454545455 -1.6 \\
-2.34545454545455 -1.5 \\
-2.4 -1.44545454545454 \\
-2.44545454545455 -1.4 \\
-2.5 -1.34545454545455 \\
-2.54545454545455 -1.3 \\
-2.54545454545455 -1.2 \\
-2.6 -1.14545454545454 \\
-2.64545454545455 -1.1 \\
-2.64545454545455 -1 \\
-2.64545454545455 -0.899999999999999 \\
-2.7 -0.845454545454545 \\
-2.74545454545455 -0.8 \\
-2.74545454545455 -0.699999999999999 \\
-2.74545454545455 -0.6 \\
-2.74545454545455 -0.5 \\
-2.74545454545455 -0.399999999999999 \\
-2.74545454545455 -0.3 \\
-2.74545454545455 -0.199999999999999 \\
-2.74545454545455 -0.0999999999999996 \\
-2.74545454545455 0 \\
-2.74545454545455 0.0999999999999996 \\
-2.74545454545455 0.199999999999999 \\
-2.74545454545455 0.3 \\
-2.74545454545455 0.399999999999999 \\
-2.74545454545455 0.5 \\
-2.74545454545455 0.6 \\
-2.74545454545455 0.699999999999999 \\
-2.74545454545455 0.8 \\
-2.7 0.845454545454545 \\
NaN NaN \\
};


\addplot[dashed,line width=2.0pt,draw=mycolor1,forget plot]
table[row sep=crcr]{
x y\\
-2.7 0.836363636363636 \\
-2.63636363636364 0.899999999999999 \\
-2.63636363636364 1 \\
-2.63636363636364 1.1 \\
-2.6 1.13636363636364 \\
-2.53636363636364 1.2 \\
-2.53636363636364 1.3 \\
-2.5 1.33636363636364 \\
-2.43636363636364 1.4 \\
-2.4 1.43636363636364 \\
-2.33636363636364 1.5 \\
-2.33636363636364 1.6 \\
-2.3 1.63636363636364 \\
-2.23636363636364 1.7 \\
-2.2 1.73636363636364 \\
-2.1 1.73636363636364 \\
-2.03636363636364 1.8 \\
-2 1.83636363636364 \\
-1.93636363636364 1.9 \\
-1.9 1.93636363636364 \\
-1.8 1.93636363636364 \\
-1.73636363636364 2 \\
-1.7 2.03636363636364 \\
-1.6 2.03636363636364 \\
-1.53636363636364 2.1 \\
-1.5 2.13636363636364 \\
-1.43636363636364 2.2 \\
-1.4 2.23636363636364 \\
-1.3 2.23636363636364 \\
-1.23636363636364 2.3 \\
-1.2 2.33636363636364 \\
-1.13636363636364 2.4 \\
-1.1 2.43636363636364 \\
-1.03636363636364 2.5 \\
-1 2.53636363636364 \\
-0.936363636363637 2.6 \\
-0.9 2.63636363636364 \\
-0.836363636363637 2.7 \\
-0.8 2.73636363636364 \\
-0.736363636363637 2.8 \\
-0.7 2.83636363636364 \\
-0.6 2.83636363636364 \\
-0.536363636363637 2.9 \\
-0.5 2.93636363636364 \\
-0.4 2.93636363636364 \\
-0.3 2.93636363636364 \\
-0.2 2.93636363636364 \\
-0.163636363636363 2.9 \\
-0.1 2.83636363636364 \\
0 2.83636363636364 \\
0.0363636363636366 2.8 \\
0.0999999999999999 2.73636363636364 \\
0.136363636363637 2.7 \\
0.136363636363637 2.6 \\
0.2 2.53636363636364 \\
0.236363636363637 2.5 \\
0.236363636363637 2.4 \\
0.236363636363637 2.3 \\
0.236363636363637 2.2 \\
0.236363636363637 2.1 \\
0.2 2.06363636363636 \\
0.136363636363637 2 \\
0.136363636363637 1.9 \\
0.136363636363637 1.8 \\
0.136363636363637 1.7 \\
0.0999999999999999 1.66363636363636 \\
0.0363636363636366 1.6 \\
0.0363636363636366 1.5 \\
0.0363636363636366 1.4 \\
0.0363636363636366 1.3 \\
0.0363636363636366 1.2 \\
0.0363636363636366 1.1 \\
0.0363636363636366 1 \\
0.0363636363636366 0.899999999999999 \\
0.0363636363636366 0.8 \\
0.0363636363636366 0.699999999999999 \\
0.0363636363636366 0.6 \\
0.0363636363636366 0.5 \\
0.0363636363636366 0.399999999999999 \\
0.0363636363636366 0.3 \\
0.0363636363636366 0.199999999999999 \\
0.0363636363636366 0.0999999999999996 \\
0 0.0636363636363631 \\
-0.0636363636363634 0 \\
0 -0.0636363636363631 \\
0.0363636363636366 -0.0999999999999996 \\
0.0363636363636366 -0.199999999999999 \\
0.0363636363636366 -0.3 \\
0.0363636363636366 -0.399999999999999 \\
0.0363636363636366 -0.5 \\
0.0363636363636366 -0.6 \\
0.0363636363636366 -0.699999999999999 \\
0.0363636363636366 -0.8 \\
0.0363636363636366 -0.899999999999999 \\
0.0363636363636366 -1 \\
0.0363636363636366 -1.1 \\
0.0363636363636366 -1.2 \\
0.0363636363636366 -1.3 \\
0.0363636363636366 -1.4 \\
0.0363636363636366 -1.5 \\
0.0363636363636366 -1.6 \\
0.0999999999999999 -1.66363636363636 \\
0.136363636363637 -1.7 \\
0.136363636363637 -1.8 \\
0.136363636363637 -1.9 \\
0.136363636363637 -2 \\
0.2 -2.06363636363636 \\
0.236363636363637 -2.1 \\
0.236363636363637 -2.2 \\
0.236363636363637 -2.3 \\
0.236363636363637 -2.4 \\
0.236363636363637 -2.5 \\
0.2 -2.53636363636364 \\
0.136363636363637 -2.6 \\
0.136363636363637 -2.7 \\
0.0999999999999999 -2.73636363636364 \\
0.0363636363636366 -2.8 \\
0 -2.83636363636364 \\
-0.1 -2.83636363636364 \\
-0.163636363636363 -2.9 \\
-0.2 -2.93636363636364 \\
-0.3 -2.93636363636364 \\
-0.4 -2.93636363636364 \\
-0.5 -2.93636363636364 \\
-0.536363636363637 -2.9 \\
-0.6 -2.83636363636364 \\
-0.7 -2.83636363636364 \\
-0.736363636363637 -2.8 \\
-0.8 -2.73636363636364 \\
-0.836363636363637 -2.7 \\
-0.9 -2.63636363636364 \\
-0.936363636363637 -2.6 \\
-1 -2.53636363636364 \\
-1.03636363636364 -2.5 \\
-1.1 -2.43636363636364 \\
-1.13636363636364 -2.4 \\
-1.2 -2.33636363636364 \\
-1.23636363636364 -2.3 \\
-1.3 -2.23636363636364 \\
-1.4 -2.23636363636364 \\
-1.43636363636364 -2.2 \\
-1.5 -2.13636363636364 \\
-1.53636363636364 -2.1 \\
-1.6 -2.03636363636364 \\
-1.7 -2.03636363636364 \\
-1.73636363636364 -2 \\
-1.8 -1.93636363636364 \\
-1.9 -1.93636363636364 \\
-1.93636363636364 -1.9 \\
-2 -1.83636363636364 \\
-2.03636363636364 -1.8 \\
-2.1 -1.73636363636364 \\
-2.2 -1.73636363636364 \\
-2.23636363636364 -1.7 \\
-2.3 -1.63636363636364 \\
-2.33636363636364 -1.6 \\
-2.33636363636364 -1.5 \\
-2.4 -1.43636363636364 \\
-2.43636363636364 -1.4 \\
-2.5 -1.33636363636364 \\
-2.53636363636364 -1.3 \\
-2.53636363636364 -1.2 \\
-2.6 -1.13636363636364 \\
-2.63636363636364 -1.1 \\
-2.63636363636364 -1 \\
-2.63636363636364 -0.899999999999999 \\
-2.7 -0.836363636363636 \\
-2.73636363636364 -0.8 \\
-2.73636363636364 -0.699999999999999 \\
-2.73636363636364 -0.6 \\
-2.73636363636364 -0.5 \\
-2.73636363636364 -0.399999999999999 \\
-2.73636363636364 -0.3 \\
-2.73636363636364 -0.199999999999999 \\
-2.73636363636364 -0.0999999999999996 \\
-2.73636363636364 0 \\
-2.73636363636364 0.0999999999999996 \\
-2.73636363636364 0.199999999999999 \\
-2.73636363636364 0.3 \\
-2.73636363636364 0.399999999999999 \\
-2.73636363636364 0.5 \\
-2.73636363636364 0.6 \\
-2.73636363636364 0.699999999999999 \\
-2.73636363636364 0.8 \\
-2.7 0.836363636363636 \\
NaN NaN \\
};


\addplot[dashed,line width=2.0pt,draw=mycolor1,forget plot]
table[row sep=crcr]{
x y\\
-2.7 0.827272727272726 \\
-2.62727272727273 0.899999999999999 \\
-2.62727272727273 1 \\
-2.62727272727273 1.1 \\
-2.6 1.12727272727273 \\
-2.52727272727273 1.2 \\
-2.52727272727273 1.3 \\
-2.5 1.32727272727273 \\
-2.42727272727273 1.4 \\
-2.4 1.42727272727273 \\
-2.32727272727273 1.5 \\
-2.32727272727273 1.6 \\
-2.3 1.62727272727273 \\
-2.22727272727273 1.7 \\
-2.2 1.72727272727273 \\
-2.1 1.72727272727273 \\
-2.02727272727273 1.8 \\
-2 1.82727272727273 \\
-1.92727272727273 1.9 \\
-1.9 1.92727272727273 \\
-1.8 1.92727272727273 \\
-1.72727272727273 2 \\
-1.7 2.02727272727273 \\
-1.6 2.02727272727273 \\
-1.52727272727273 2.1 \\
-1.5 2.12727272727273 \\
-1.42727272727273 2.2 \\
-1.4 2.22727272727273 \\
-1.3 2.22727272727273 \\
-1.22727272727273 2.3 \\
-1.2 2.32727272727273 \\
-1.12727272727273 2.4 \\
-1.1 2.42727272727273 \\
-1.02727272727273 2.5 \\
-1 2.52727272727273 \\
-0.927272727272728 2.6 \\
-0.9 2.62727272727273 \\
-0.827272727272728 2.7 \\
-0.8 2.72727272727273 \\
-0.727272727272728 2.8 \\
-0.7 2.82727272727273 \\
-0.6 2.82727272727273 \\
-0.527272727272727 2.9 \\
-0.5 2.92727272727273 \\
-0.4 2.92727272727273 \\
-0.3 2.92727272727273 \\
-0.2 2.92727272727273 \\
-0.172727272727273 2.9 \\
-0.1 2.82727272727273 \\
0 2.82727272727273 \\
0.0272727272727273 2.8 \\
0.0999999999999999 2.72727272727273 \\
0.127272727272727 2.7 \\
0.127272727272727 2.6 \\
0.2 2.52727272727273 \\
0.227272727272727 2.5 \\
0.227272727272727 2.4 \\
0.227272727272727 2.3 \\
0.227272727272727 2.2 \\
0.227272727272727 2.1 \\
0.2 2.07272727272727 \\
0.127272727272727 2 \\
0.127272727272727 1.9 \\
0.127272727272727 1.8 \\
0.127272727272727 1.7 \\
0.0999999999999999 1.67272727272727 \\
0.0272727272727273 1.6 \\
0.0272727272727273 1.5 \\
0.0272727272727273 1.4 \\
0.0272727272727273 1.3 \\
0.0272727272727273 1.2 \\
0.0272727272727273 1.1 \\
0.0272727272727273 1 \\
0.0272727272727273 0.899999999999999 \\
0.0272727272727273 0.8 \\
0.0272727272727273 0.699999999999999 \\
0.0272727272727273 0.6 \\
0.0272727272727273 0.5 \\
0.0272727272727273 0.399999999999999 \\
0.0272727272727273 0.3 \\
0.0272727272727273 0.199999999999999 \\
0.0272727272727273 0.0999999999999996 \\
0 0.0727272727272724 \\
-0.0727272727272727 0 \\
0 -0.0727272727272724 \\
0.0272727272727273 -0.0999999999999996 \\
0.0272727272727273 -0.199999999999999 \\
0.0272727272727273 -0.3 \\
0.0272727272727273 -0.399999999999999 \\
0.0272727272727273 -0.5 \\
0.0272727272727273 -0.6 \\
0.0272727272727273 -0.699999999999999 \\
0.0272727272727273 -0.8 \\
0.0272727272727273 -0.899999999999999 \\
0.0272727272727273 -1 \\
0.0272727272727273 -1.1 \\
0.0272727272727273 -1.2 \\
0.0272727272727273 -1.3 \\
0.0272727272727273 -1.4 \\
0.0272727272727273 -1.5 \\
0.0272727272727273 -1.6 \\
0.0999999999999999 -1.67272727272727 \\
0.127272727272727 -1.7 \\
0.127272727272727 -1.8 \\
0.127272727272727 -1.9 \\
0.127272727272727 -2 \\
0.2 -2.07272727272727 \\
0.227272727272727 -2.1 \\
0.227272727272727 -2.2 \\
0.227272727272727 -2.3 \\
0.227272727272727 -2.4 \\
0.227272727272727 -2.5 \\
0.2 -2.52727272727273 \\
0.127272727272727 -2.6 \\
0.127272727272727 -2.7 \\
0.0999999999999999 -2.72727272727273 \\
0.0272727272727273 -2.8 \\
0 -2.82727272727273 \\
-0.1 -2.82727272727273 \\
-0.172727272727273 -2.9 \\
-0.2 -2.92727272727273 \\
-0.3 -2.92727272727273 \\
-0.4 -2.92727272727273 \\
-0.5 -2.92727272727273 \\
-0.527272727272727 -2.9 \\
-0.6 -2.82727272727273 \\
-0.7 -2.82727272727273 \\
-0.727272727272728 -2.8 \\
-0.8 -2.72727272727273 \\
-0.827272727272728 -2.7 \\
-0.9 -2.62727272727273 \\
-0.927272727272728 -2.6 \\
-1 -2.52727272727273 \\
-1.02727272727273 -2.5 \\
-1.1 -2.42727272727273 \\
-1.12727272727273 -2.4 \\
-1.2 -2.32727272727273 \\
-1.22727272727273 -2.3 \\
-1.3 -2.22727272727273 \\
-1.4 -2.22727272727273 \\
-1.42727272727273 -2.2 \\
-1.5 -2.12727272727273 \\
-1.52727272727273 -2.1 \\
-1.6 -2.02727272727273 \\
-1.7 -2.02727272727273 \\
-1.72727272727273 -2 \\
-1.8 -1.92727272727273 \\
-1.9 -1.92727272727273 \\
-1.92727272727273 -1.9 \\
-2 -1.82727272727273 \\
-2.02727272727273 -1.8 \\
-2.1 -1.72727272727273 \\
-2.2 -1.72727272727273 \\
-2.22727272727273 -1.7 \\
-2.3 -1.62727272727273 \\
-2.32727272727273 -1.6 \\
-2.32727272727273 -1.5 \\
-2.4 -1.42727272727273 \\
-2.42727272727273 -1.4 \\
-2.5 -1.32727272727273 \\
-2.52727272727273 -1.3 \\
-2.52727272727273 -1.2 \\
-2.6 -1.12727272727273 \\
-2.62727272727273 -1.1 \\
-2.62727272727273 -1 \\
-2.62727272727273 -0.899999999999999 \\
-2.7 -0.827272727272727 \\
-2.72727272727273 -0.8 \\
-2.72727272727273 -0.699999999999999 \\
-2.72727272727273 -0.6 \\
-2.72727272727273 -0.5 \\
-2.72727272727273 -0.399999999999999 \\
-2.72727272727273 -0.3 \\
-2.72727272727273 -0.199999999999999 \\
-2.72727272727273 -0.0999999999999996 \\
-2.72727272727273 0 \\
-2.72727272727273 0.0999999999999996 \\
-2.72727272727273 0.199999999999999 \\
-2.72727272727273 0.3 \\
-2.72727272727273 0.399999999999999 \\
-2.72727272727273 0.5 \\
-2.72727272727273 0.6 \\
-2.72727272727273 0.699999999999999 \\
-2.72727272727273 0.8 \\
-2.7 0.827272727272726 \\
NaN NaN \\
};


\addplot[dashed,line width=2.0pt,draw=mycolor1,forget plot]
table[row sep=crcr]{
x y\\
-2.7 0.818181818181818 \\
-2.61818181818182 0.899999999999999 \\
-2.61818181818182 1 \\
-2.61818181818182 1.1 \\
-2.6 1.11818181818182 \\
-2.51818181818182 1.2 \\
-2.51818181818182 1.3 \\
-2.5 1.31818181818182 \\
-2.41818181818182 1.4 \\
-2.4 1.41818181818182 \\
-2.31818181818182 1.5 \\
-2.31818181818182 1.6 \\
-2.3 1.61818181818182 \\
-2.21818181818182 1.7 \\
-2.2 1.71818181818182 \\
-2.1 1.71818181818182 \\
-2.01818181818182 1.8 \\
-2 1.81818181818182 \\
-1.91818181818182 1.9 \\
-1.9 1.91818181818182 \\
-1.8 1.91818181818182 \\
-1.71818181818182 2 \\
-1.7 2.01818181818182 \\
-1.6 2.01818181818182 \\
-1.51818181818182 2.1 \\
-1.5 2.11818181818182 \\
-1.41818181818182 2.2 \\
-1.4 2.21818181818182 \\
-1.3 2.21818181818182 \\
-1.21818181818182 2.3 \\
-1.2 2.31818181818182 \\
-1.11818181818182 2.4 \\
-1.1 2.41818181818182 \\
-1.01818181818182 2.5 \\
-1 2.51818181818182 \\
-0.918181818181819 2.6 \\
-0.9 2.61818181818182 \\
-0.818181818181818 2.7 \\
-0.8 2.71818181818182 \\
-0.718181818181818 2.8 \\
-0.7 2.81818181818182 \\
-0.6 2.81818181818182 \\
-0.518181818181818 2.9 \\
-0.5 2.91818181818182 \\
-0.4 2.91818181818182 \\
-0.3 2.91818181818182 \\
-0.2 2.91818181818182 \\
-0.181818181818182 2.9 \\
-0.1 2.81818181818182 \\
0 2.81818181818182 \\
0.018181818181818 2.8 \\
0.0999999999999999 2.71818181818182 \\
0.118181818181818 2.7 \\
0.118181818181818 2.6 \\
0.2 2.51818181818182 \\
0.218181818181818 2.5 \\
0.218181818181818 2.4 \\
0.218181818181818 2.3 \\
0.218181818181818 2.2 \\
0.218181818181818 2.1 \\
0.2 2.08181818181818 \\
0.118181818181818 2 \\
0.118181818181818 1.9 \\
0.118181818181818 1.8 \\
0.118181818181818 1.7 \\
0.0999999999999999 1.68181818181818 \\
0.018181818181818 1.6 \\
0.018181818181818 1.5 \\
0.018181818181818 1.4 \\
0.018181818181818 1.3 \\
0.018181818181818 1.2 \\
0.018181818181818 1.1 \\
0.018181818181818 1 \\
0.018181818181818 0.899999999999999 \\
0.018181818181818 0.8 \\
0.018181818181818 0.699999999999999 \\
0.018181818181818 0.6 \\
0.018181818181818 0.5 \\
0.018181818181818 0.399999999999999 \\
0.018181818181818 0.3 \\
0.018181818181818 0.199999999999999 \\
0.018181818181818 0.0999999999999996 \\
0 0.0818181818181817 \\
-0.0818181818181821 0 \\
0 -0.0818181818181817 \\
0.018181818181818 -0.0999999999999996 \\
0.018181818181818 -0.199999999999999 \\
0.018181818181818 -0.3 \\
0.018181818181818 -0.399999999999999 \\
0.018181818181818 -0.5 \\
0.018181818181818 -0.6 \\
0.018181818181818 -0.699999999999999 \\
0.018181818181818 -0.8 \\
0.018181818181818 -0.899999999999999 \\
0.018181818181818 -1 \\
0.018181818181818 -1.1 \\
0.018181818181818 -1.2 \\
0.018181818181818 -1.3 \\
0.018181818181818 -1.4 \\
0.018181818181818 -1.5 \\
0.018181818181818 -1.6 \\
0.0999999999999999 -1.68181818181818 \\
0.118181818181818 -1.7 \\
0.118181818181818 -1.8 \\
0.118181818181818 -1.9 \\
0.118181818181818 -2 \\
0.2 -2.08181818181818 \\
0.218181818181818 -2.1 \\
0.218181818181818 -2.2 \\
0.218181818181818 -2.3 \\
0.218181818181818 -2.4 \\
0.218181818181818 -2.5 \\
0.2 -2.51818181818182 \\
0.118181818181818 -2.6 \\
0.118181818181818 -2.7 \\
0.0999999999999999 -2.71818181818182 \\
0.018181818181818 -2.8 \\
0 -2.81818181818182 \\
-0.1 -2.81818181818182 \\
-0.181818181818182 -2.9 \\
-0.2 -2.91818181818182 \\
-0.3 -2.91818181818182 \\
-0.4 -2.91818181818182 \\
-0.5 -2.91818181818182 \\
-0.518181818181818 -2.9 \\
-0.6 -2.81818181818182 \\
-0.7 -2.81818181818182 \\
-0.718181818181818 -2.8 \\
-0.8 -2.71818181818182 \\
-0.818181818181818 -2.7 \\
-0.9 -2.61818181818182 \\
-0.918181818181819 -2.6 \\
-1 -2.51818181818182 \\
-1.01818181818182 -2.5 \\
-1.1 -2.41818181818182 \\
-1.11818181818182 -2.4 \\
-1.2 -2.31818181818182 \\
-1.21818181818182 -2.3 \\
-1.3 -2.21818181818182 \\
-1.4 -2.21818181818182 \\
-1.41818181818182 -2.2 \\
-1.5 -2.11818181818182 \\
-1.51818181818182 -2.1 \\
-1.6 -2.01818181818182 \\
-1.7 -2.01818181818182 \\
-1.71818181818182 -2 \\
-1.8 -1.91818181818182 \\
-1.9 -1.91818181818182 \\
-1.91818181818182 -1.9 \\
-2 -1.81818181818182 \\
-2.01818181818182 -1.8 \\
-2.1 -1.71818181818182 \\
-2.2 -1.71818181818182 \\
-2.21818181818182 -1.7 \\
-2.3 -1.61818181818182 \\
-2.31818181818182 -1.6 \\
-2.31818181818182 -1.5 \\
-2.4 -1.41818181818182 \\
-2.41818181818182 -1.4 \\
-2.5 -1.31818181818182 \\
-2.51818181818182 -1.3 \\
-2.51818181818182 -1.2 \\
-2.6 -1.11818181818182 \\
-2.61818181818182 -1.1 \\
-2.61818181818182 -1 \\
-2.61818181818182 -0.899999999999999 \\
-2.7 -0.818181818181818 \\
-2.71818181818182 -0.8 \\
-2.71818181818182 -0.699999999999999 \\
-2.71818181818182 -0.6 \\
-2.71818181818182 -0.5 \\
-2.71818181818182 -0.399999999999999 \\
-2.71818181818182 -0.3 \\
-2.71818181818182 -0.199999999999999 \\
-2.71818181818182 -0.0999999999999996 \\
-2.71818181818182 0 \\
-2.71818181818182 0.0999999999999996 \\
-2.71818181818182 0.199999999999999 \\
-2.71818181818182 0.3 \\
-2.71818181818182 0.399999999999999 \\
-2.71818181818182 0.5 \\
-2.71818181818182 0.6 \\
-2.71818181818182 0.699999999999999 \\
-2.71818181818182 0.8 \\
-2.7 0.818181818181818 \\
NaN NaN \\
};


\addplot[dashed,line width=2.0pt,draw=mycolor1,forget plot]
table[row sep=crcr]{
x y\\
-2.7 0.809090909090909 \\
-2.60909090909091 0.899999999999999 \\
-2.60909090909091 1 \\
-2.60909090909091 1.1 \\
-2.6 1.10909090909091 \\
-2.50909090909091 1.2 \\
-2.50909090909091 1.3 \\
-2.5 1.30909090909091 \\
-2.40909090909091 1.4 \\
-2.4 1.40909090909091 \\
-2.30909090909091 1.5 \\
-2.30909090909091 1.6 \\
-2.3 1.60909090909091 \\
-2.20909090909091 1.7 \\
-2.2 1.70909090909091 \\
-2.1 1.70909090909091 \\
-2.00909090909091 1.8 \\
-2 1.80909090909091 \\
-1.90909090909091 1.9 \\
-1.9 1.90909090909091 \\
-1.8 1.90909090909091 \\
-1.70909090909091 2 \\
-1.7 2.00909090909091 \\
-1.6 2.00909090909091 \\
-1.50909090909091 2.1 \\
-1.5 2.10909090909091 \\
-1.40909090909091 2.2 \\
-1.4 2.20909090909091 \\
-1.3 2.20909090909091 \\
-1.20909090909091 2.3 \\
-1.2 2.30909090909091 \\
-1.10909090909091 2.4 \\
-1.1 2.40909090909091 \\
-1.00909090909091 2.5 \\
-1 2.50909090909091 \\
-0.909090909090909 2.6 \\
-0.9 2.60909090909091 \\
-0.80909090909091 2.7 \\
-0.8 2.70909090909091 \\
-0.70909090909091 2.8 \\
-0.7 2.80909090909091 \\
-0.6 2.80909090909091 \\
-0.509090909090909 2.9 \\
-0.5 2.90909090909091 \\
-0.4 2.90909090909091 \\
-0.3 2.90909090909091 \\
-0.2 2.90909090909091 \\
-0.190909090909091 2.9 \\
-0.1 2.80909090909091 \\
0 2.80909090909091 \\
0.00909090909090934 2.8 \\
0.0999999999999999 2.70909090909091 \\
0.109090909090909 2.7 \\
0.109090909090909 2.6 \\
0.2 2.50909090909091 \\
0.209090909090909 2.5 \\
0.209090909090909 2.4 \\
0.209090909090909 2.3 \\
0.209090909090909 2.2 \\
0.209090909090909 2.1 \\
0.2 2.09090909090909 \\
0.109090909090909 2 \\
0.109090909090909 1.9 \\
0.109090909090909 1.8 \\
0.109090909090909 1.7 \\
0.0999999999999999 1.69090909090909 \\
0.00909090909090934 1.6 \\
0.00909090909090934 1.5 \\
0.00909090909090934 1.4 \\
0.00909090909090934 1.3 \\
0.00909090909090934 1.2 \\
0.00909090909090934 1.1 \\
0.00909090909090934 1 \\
0.00909090909090934 0.899999999999999 \\
0.00909090909090934 0.8 \\
0.00909090909090934 0.699999999999999 \\
0.00909090909090934 0.6 \\
0.00909090909090934 0.5 \\
0.00909090909090934 0.399999999999999 \\
0.00909090909090934 0.3 \\
0.00909090909090934 0.199999999999999 \\
0.00909090909090934 0.0999999999999996 \\
0 0.0909090909090903 \\
-0.0909090909090907 0 \\
0 -0.0909090909090903 \\
0.00909090909090934 -0.0999999999999996 \\
0.00909090909090934 -0.199999999999999 \\
0.00909090909090934 -0.3 \\
0.00909090909090934 -0.399999999999999 \\
0.00909090909090934 -0.5 \\
0.00909090909090934 -0.6 \\
0.00909090909090934 -0.699999999999999 \\
0.00909090909090934 -0.8 \\
0.00909090909090934 -0.899999999999999 \\
0.00909090909090934 -1 \\
0.00909090909090934 -1.1 \\
0.00909090909090934 -1.2 \\
0.00909090909090934 -1.3 \\
0.00909090909090934 -1.4 \\
0.00909090909090934 -1.5 \\
0.00909090909090934 -1.6 \\
0.0999999999999999 -1.69090909090909 \\
0.109090909090909 -1.7 \\
0.109090909090909 -1.8 \\
0.109090909090909 -1.9 \\
0.109090909090909 -2 \\
0.2 -2.09090909090909 \\
0.209090909090909 -2.1 \\
0.209090909090909 -2.2 \\
0.209090909090909 -2.3 \\
0.209090909090909 -2.4 \\
0.209090909090909 -2.5 \\
0.2 -2.50909090909091 \\
0.109090909090909 -2.6 \\
0.109090909090909 -2.7 \\
0.0999999999999999 -2.70909090909091 \\
0.00909090909090934 -2.8 \\
0 -2.80909090909091 \\
-0.1 -2.80909090909091 \\
-0.190909090909091 -2.9 \\
-0.2 -2.90909090909091 \\
-0.3 -2.90909090909091 \\
-0.4 -2.90909090909091 \\
-0.5 -2.90909090909091 \\
-0.509090909090909 -2.9 \\
-0.6 -2.80909090909091 \\
-0.7 -2.80909090909091 \\
-0.70909090909091 -2.8 \\
-0.8 -2.70909090909091 \\
-0.80909090909091 -2.7 \\
-0.9 -2.60909090909091 \\
-0.909090909090909 -2.6 \\
-1 -2.50909090909091 \\
-1.00909090909091 -2.5 \\
-1.1 -2.40909090909091 \\
-1.10909090909091 -2.4 \\
-1.2 -2.30909090909091 \\
-1.20909090909091 -2.3 \\
-1.3 -2.20909090909091 \\
-1.4 -2.20909090909091 \\
-1.40909090909091 -2.2 \\
-1.5 -2.10909090909091 \\
-1.50909090909091 -2.1 \\
-1.6 -2.00909090909091 \\
-1.7 -2.00909090909091 \\
-1.70909090909091 -2 \\
-1.8 -1.90909090909091 \\
-1.9 -1.90909090909091 \\
-1.90909090909091 -1.9 \\
-2 -1.80909090909091 \\
-2.00909090909091 -1.8 \\
-2.1 -1.70909090909091 \\
-2.2 -1.70909090909091 \\
-2.20909090909091 -1.7 \\
-2.3 -1.60909090909091 \\
-2.30909090909091 -1.6 \\
-2.30909090909091 -1.5 \\
-2.4 -1.40909090909091 \\
-2.40909090909091 -1.4 \\
-2.5 -1.30909090909091 \\
-2.50909090909091 -1.3 \\
-2.50909090909091 -1.2 \\
-2.6 -1.10909090909091 \\
-2.60909090909091 -1.1 \\
-2.60909090909091 -1 \\
-2.60909090909091 -0.899999999999999 \\
-2.7 -0.809090909090909 \\
-2.70909090909091 -0.8 \\
-2.70909090909091 -0.699999999999999 \\
-2.70909090909091 -0.6 \\
-2.70909090909091 -0.5 \\
-2.70909090909091 -0.399999999999999 \\
-2.70909090909091 -0.3 \\
-2.70909090909091 -0.199999999999999 \\
-2.70909090909091 -0.0999999999999996 \\
-2.70909090909091 0 \\
-2.70909090909091 0.0999999999999996 \\
-2.70909090909091 0.199999999999999 \\
-2.70909090909091 0.3 \\
-2.70909090909091 0.399999999999999 \\
-2.70909090909091 0.5 \\
-2.70909090909091 0.6 \\
-2.70909090909091 0.699999999999999 \\
-2.70909090909091 0.8 \\
-2.7 0.809090909090909 \\
NaN NaN \\
};


\addplot[dashed,line width=2.0pt,draw=mycolor1,forget plot]
table[row sep=crcr]{
x y\\
-2.7 0.8 \\
-2.6 0.899999999999999 \\
-2.6 1 \\
-2.6 1.1 \\
-2.6 1.1 \\
-2.5 1.2 \\
-2.5 1.3 \\
-2.5 1.3 \\
-2.4 1.4 \\
-2.4 1.4 \\
-2.3 1.5 \\
-2.3 1.6 \\
-2.3 1.6 \\
-2.2 1.7 \\
-2.2 1.7 \\
-2.1 1.7 \\
-2 1.8 \\
-2 1.8 \\
-1.9 1.9 \\
-1.9 1.9 \\
-1.8 1.9 \\
-1.7 2 \\
-1.7 2 \\
-1.6 2 \\
-1.5 2.1 \\
-1.5 2.1 \\
-1.4 2.2 \\
-1.4 2.2 \\
-1.3 2.2 \\
-1.2 2.3 \\
-1.2 2.3 \\
-1.1 2.4 \\
-1.1 2.4 \\
-1 2.5 \\
-1 2.5 \\
-0.9 2.6 \\
-0.9 2.6 \\
-0.8 2.7 \\
-0.8 2.7 \\
-0.7 2.8 \\
-0.7 2.8 \\
-0.6 2.8 \\
-0.5 2.9 \\
-0.5 2.9 \\
-0.4 2.9 \\
-0.3 2.9 \\
-0.2 2.9 \\
-0.2 2.9 \\
-0.1 2.8 \\
0 2.8 \\
0 2.8 \\
0.0999999999999999 2.7 \\
0.0999999999999999 2.7 \\
0.0999999999999999 2.6 \\
0.2 2.5 \\
0.2 2.5 \\
0.2 2.4 \\
0.2 2.3 \\
0.2 2.2 \\
0.2 2.1 \\
0.2 2.1 \\
0.0999999999999999 2 \\
0.0999999999999999 1.9 \\
0.0999999999999999 1.8 \\
0.0999999999999999 1.7 \\
0.0999999999999999 1.7 \\
0 1.6 \\
0 1.5 \\
0 1.4 \\
0 1.3 \\
0 1.2 \\
0 1.1 \\
0 1 \\
0 0.899999999999999 \\
0 0.8 \\
0 0.699999999999999 \\
0 0.6 \\
0 0.5 \\
0 0.399999999999999 \\
0 0.3 \\
0 0.199999999999999 \\
0 0.0999999999999996 \\
0 0.0999999999999996 \\
-0.1 0 \\
0 -0.0999999999999996 \\
0 -0.0999999999999996 \\
0 -0.199999999999999 \\
0 -0.3 \\
0 -0.399999999999999 \\
0 -0.5 \\
0 -0.6 \\
0 -0.699999999999999 \\
0 -0.8 \\
0 -0.899999999999999 \\
0 -1 \\
0 -1.1 \\
0 -1.2 \\
0 -1.3 \\
0 -1.4 \\
0 -1.5 \\
0 -1.6 \\
0.0999999999999999 -1.7 \\
0.0999999999999999 -1.7 \\
0.0999999999999999 -1.8 \\
0.0999999999999999 -1.9 \\
0.0999999999999999 -2 \\
0.2 -2.1 \\
0.2 -2.1 \\
0.2 -2.2 \\
0.2 -2.3 \\
0.2 -2.4 \\
0.2 -2.5 \\
0.2 -2.5 \\
0.0999999999999999 -2.6 \\
0.0999999999999999 -2.7 \\
0.0999999999999999 -2.7 \\
0 -2.8 \\
0 -2.8 \\
-0.1 -2.8 \\
-0.2 -2.9 \\
-0.2 -2.9 \\
-0.3 -2.9 \\
-0.4 -2.9 \\
-0.5 -2.9 \\
-0.5 -2.9 \\
-0.6 -2.8 \\
-0.7 -2.8 \\
-0.7 -2.8 \\
-0.8 -2.7 \\
-0.8 -2.7 \\
-0.9 -2.6 \\
-0.9 -2.6 \\
-1 -2.5 \\
-1 -2.5 \\
-1.1 -2.4 \\
-1.1 -2.4 \\
-1.2 -2.3 \\
-1.2 -2.3 \\
-1.3 -2.2 \\
-1.4 -2.2 \\
-1.4 -2.2 \\
-1.5 -2.1 \\
-1.5 -2.1 \\
-1.6 -2 \\
-1.7 -2 \\
-1.7 -2 \\
-1.8 -1.9 \\
-1.9 -1.9 \\
-1.9 -1.9 \\
-2 -1.8 \\
-2 -1.8 \\
-2.1 -1.7 \\
-2.2 -1.7 \\
-2.2 -1.7 \\
-2.3 -1.6 \\
-2.3 -1.6 \\
-2.3 -1.5 \\
-2.4 -1.4 \\
-2.4 -1.4 \\
-2.5 -1.3 \\
-2.5 -1.3 \\
-2.5 -1.2 \\
-2.6 -1.1 \\
-2.6 -1.1 \\
-2.6 -1 \\
-2.6 -0.899999999999999 \\
-2.7 -0.8 \\
-2.7 -0.8 \\
-2.7 -0.699999999999999 \\
-2.7 -0.6 \\
-2.7 -0.5 \\
-2.7 -0.399999999999999 \\
-2.7 -0.3 \\
-2.7 -0.199999999999999 \\
-2.7 -0.0999999999999996 \\
-2.7 0 \\
-2.7 0.0999999999999996 \\
-2.7 0.199999999999999 \\
-2.7 0.3 \\
-2.7 0.399999999999999 \\
-2.7 0.5 \\
-2.7 0.6 \\
-2.7 0.699999999999999 \\
-2.7 0.8 \\
-2.7 0.8 \\
NaN NaN \\
};

\end{axis}
\end{tikzpicture}%
\caption{Stabilit�tsgebiete f�r einige Einschrittverfahren}
\label{fig:stabis}
\end{figure}

\section{Mehrschrittverfahren}
Die Stabilit�t von Mehrschrittverfahren ist nicht so einfach allgemein nachzuweisen, wie bei den allgemeinen Runge-Kutta Verfahren. In \cite{gruene} wird gezeigt, dass die \glqq Stabilit�t\grqq \ des Polynoms $P_a(z)$ eine Voraussetzung f�r die \textit{Konvergenz} des Fehlers eines Mehrschrittverfahrens darstellt. Dabei gilt das Polynom als stabil, wenn f�r alle Nullstellen $\lambda_i$ gilt:
\begin{align}
	&|\lambda_i| \leq 1, \\
	&|\lambda_i|=1 \ \Rightarrow \text{$\lambda_i$ ist einfache Nullstelle.}
\end{align}
Diese \glqq Stabilit�tsbedingung\grqq \ darf jedoch nicht mit den f�r Einschrittverfahren gezeigten Stabilit�tskriterien verwechselt werden, da sie noch keine Aussage �ber die tats�chliche Stabilit�t des Verfahrens trifft, sondern lediglich die Konvergenz des Fehlers erm�glicht, welche bei den Einschrittverfahren ohnehin vorausgesetzt wurde.
Da eine allgemeine Aussage �ber die Stabilit�t eines Mehrschrittverfahrens schwierig ist, wird in \cite{friedrich} gezeigt, dass die Stabilit�t eines Verfahrens im konkreten Fall �berp�ft werden kann, indem es folgenderma�en umgeformt wird:
\begin{align}
	x_{i+1}+a_{k-1}x_{i}+a_{k-2}x_{i-1}+\ldots+a_0 x_{i+1-k}=0.
\end{align}
Das hei�t, die Koeffizienten $a_l (l=0,\ldots,s-1)$ beinhalten bereits die Koeffizienten des gew�hlten Verfahrens, sowie die gew�hlte Schrittweite. �hnlich wie bei der Herleitung der Polynome $P_a(z)$ und $P_b(z)$ erh�lt man damit ein \textit{charakteristisches Polynom}
\begin{align}
	P_k(z) = z^k + a_{k-1}z^{k-1}+ a_{k-2}z^{k-2}+\ldots+a_0 = \sum\limits_{l=0}^{k}a_l z^l. 
\end{align}
Das gew�hlte Verfahren ist mit der verwendeten Schrittweite genau dann stabil, wenn f�r alle Nullstellen des charakteristischen Polynoms $z_l, (l=1,\ldots,k)$ gilt:
\begin{align}
	|z_l| < 1.
\end{align}

%%%%%%%%%%%%%%%%%%%%%%%%%%%%%%%%%%%%%%%%%%
\chapter{Anwendung der Integrationsmethoden in Simulationssoftware}
% Vielleicht keine Unterpunkte, sondern nur ein Kapitel mit unter�berschriften oder so, evtl als subsubsection machen oder so... Weil ist nicht so viel.
\label{sec:anwend}
Um die Bedeutung der vorgestellten Integrationsmethoden zu verdeutlichen, soll in diesem Kapitel kurz dargelegt werden, welche Verfahren von einigen ausgew�hlten Programmen, die in der Energietechnik zum Einsatz kommen, verwendet werden.

\subsubsection{PSS\textregistered Netomac, PSS\textregistered E}
Die Firma \siemens \ liefert mit den Programmen \textsc{PSS\textregistered Netomac} und \textsc{PSS\textregistered E} gleich zwei Softwarepakete, die in der Energietechnik Anwendung finden. 
Ersteres verwendet die Trapezregel als Integrationsverfahren. Das hei�t, es wird Modellwissen verwendet, um die implizite Gleichung wirklich aufzul�sen.
Da die Trapezregel Isometrieerhaltend ist, ergeben sich keine Stabilit�tsprobleme und die Wahl der Schrittweite beeinflusst lediglich die Genauigkeit.\\
In letzterem Softwarepaket wird nach \cite{PSSE} das \glqq second-order Euler\grqq-Verfahren, oder \glqq modified Euler\grqq-Verfahren verwendet. Nach der Schreibweise in dieser Arbeit entspricht dies dem Heun-Verfahren (\ref{eq:Heun}). In diesem Fall muss die Schrittweite klein genug gew�hlt werden, um Stabilit�t zu gew�hrleisten. In \cite{PSSE} wird deshalb empfohlen, die Schrittweite zwischen $\frac{1}{5}$ und $\frac{1}{4}$ der kleinsten im Prozess auftretenden Zeitkonstanten zu w�hlen.
 
\subsubsection{PowerFactory}
Ein weitere relevante Software in der Energietechnik ist \textsc{PowerFactory} von \mbox{\textsc{DIgSILENT}}. In \cite{powerfac} wird nicht genau beschrieben, welches Integrationsverfahren verwendet wird. Es wird angedeutet, dass das Verfahren der RMS-Simulation A-stabil ist, was auf ein implizites Verfahren hindeutet. Auch hier sind dynamische Modelle hinterlegt, die das Auswerten der impliziten Gleichung erm�glichen. Dar�ber hinaus wird angegeben, dass eine automatische Schrittweitenanpassung erfolgt beziehungsweise bei EMT-Simulationen mit variabler oder fester Schrittweite gerechnet werden kann.


\subsubsection{Eurostag}
Die Software \textsc{Eurostag} des franz�sischen Energieversorgers \textsc{RTE} greift nach \cite{eurostag} bei der Simulation auf das in \cite{mixed} beschriebenen \glqq mixed adams - BDF\grqq-Vefahren zur�ck. Dabei handelt es sich um ein Pr�diktor-Korrektor Verfahren, bei dem je nach Problemstellung ein implizites Adams-Verfahren oder ein BDF-Verfahren als Korrektor eingesetzt wird. Dar�ber hinaus wird die Schrittweite, �hnlich wie bei Einschrittverfahren gezeigt, gesteuert. Wie bei allen Verfahren mit Schrittweitensteuerung kommt es hierbei f�r die Genauigkeit und f�r die Stabilit�t auf die vom Anwender eingestellte Toleranz des lokalen Fehlers an. 

























%\section{Einbindung von Bildern}
%\label{kap:einbindungbilder}
%Abbildungen sind mit Hilfe des Pakets \textit{graphicx} einzuf�gen. Sie k�nnen im PDF-Format durch die Nutzung des folgenden Codes implementiert werden.
%\begin{verbatim}
%\begin{figure}[!htb]\centering
% \includegraphics*[width = \textwidth]{beispiel}
% \caption{Beispiel f�r die Einbindung eines Bildes}
% \label{abb:beispiel}
%\end{figure}
%\end{verbatim}
%Das Ergebnis ist die Anzeige des Bildes, mittig, wie der Text breit mit der angegebenen Unterschrift. Alternativ kann bei der Breite eine absolute Angabe in mm erfolgen. �ber das label \textit{abb:beispiel} kann das Bild referenziert werden.
%\begin{figure}[!htb]\centering
% \includegraphics*[width = \textwidth]{beispiel}
% \caption{Beispiel f�r die Einbindung eines Bildes}
% \label{abb:beispiel}
%\end{figure}
%Um auf das Bild \ref{abb:beispiel} zu verweisen, bedient man sich der folgenden Funktion:
%\begin{verbatim}
%\ref{abb:beispiel}
%\end{verbatim}
%Die referenzierte Nummerierung erfolgt Kapitelweise. Will man weiterhin eine Quelle in der Bildunterschrift angeben, so ist darauf zu achten, dass die Einbindung der Bildunterschriften durch ein optionales Element (eingeschlossen in eckigen Klammern) erweitert wird, welches die Beschriftung f�r das Abbildungsverzeichnis enth�lt:
%\begin{verbatim}
% \caption[Beispiel...]{Beispiel... , aus \cite{schwab}}
%\end{verbatim}
%Diese Variante verhindert, dass LaTex die Quellen bereits im Abbildungsverzeichnis zu z�hlen anf�ngt.
%
%\section{Einbindung von Tabellen}
%\label{kap:einbindungtabellen}
%Die Tabellen sollen mit Hilfe des Pakets \textit{tabularx} eingebunden werden. Im Folgenden ist ein Beispiel f�r die Einbindung von Tabellen aufgef�hrt. Mit
%\begin{verbatim}
%\begin{table}[!htb]
%\centering
%\caption{Beispiel einer Tabelle}
%\label{tab:tabelle1}
%\begin{tabularx}{\textwidth}{|X|c|c|c|c|c|c|c|c|}
%\hline
%        & Spalte 1 & Spalte 2 & Spalte 3 & Spalte 4 & Spalte 5 \\
%\hline
%Zeile 1 &          &          &          &          &          \\
%\hline
%Zeile 2 &          &          &          &          &          \\
%\hline
%Zeile 3 &          &          &          &          &          \\
%\hline
%\end{tabularx}
%\end{table}
%\end{verbatim}
%ergibt sich die folgende Tabellenausgabe \ref{tab:tabelle1}.
%\begin{table}[!htb]
%\centering
%\caption{Beispiel einer Tabelle}
%\label{tab:tabelle1}
%\begin{tabularx}{\textwidth}{|X|c|c|c|c|c|c|c|c|}
%\hline
% 				& Spalte 1 	& Spalte 2 	& Spalte 3 	& Spalte 4 	& Spalte 5 \\
%\hline
%Zeile 1 & 					& 					& 					& 					& \\
%\hline
%Zeile 2 & 					& 					& 					& 					& \\
%\hline
%Zeile 3 & 					& 					& 					& 					& \\
%\hline
%\end{tabularx}
%\end{table}
%Bei Angabe von Quellen in der Tabellen�berschrift ist �hnlich wie im Kapitel \ref{kap:einbindungbilder} zu verfahren.
%
%\section{Eingabe von Gleichungen}
%\label{kap:einbindunggleichungen}
%Die Eingabe von Gleichungen erfolgt nach dem folgenden Beispiel:
%\begin{verbatim}
%\begin{align}
%U=R \cdot I
%\end{align}
%\end{verbatim}
%Das Ergebnis ist die folgende Darstellung mit der automatischen Nummerierung:
%\begin{align}
%U=R \cdot I
%\end{align}
%Durch die Einbindung des eigens definierten Pakets \textit{befehle.sty} stehen zudem die folgenden Funktionen f�r die Gleichungseingabe zur Verf�gung:
%\begin{tabbing}
%XXXXXX 	\= XXXXXXXXXXXXXXXXXXXXXXXXXXXXXXXXXXXXXXXXXXXXXXXXXXXXXXXX\kill
%einh{} 	\> Einheiten\\
%iz{}		\> Nicht-kursiver Index, unten\\
%izo{}		\> Nicht-kursiver Index, oben\\
%izf{}		\> Nicht-kursiver Index, unten, kleinere Schrift\\
%izof{}	\> Nicht-kursiver Index, oben, kleinere Schrift
%\end{tabbing}
%Die Anwendung der Funktionen kann anhand des folgenden Beispiels nachvollzogen werden. Mit dem Code
%\begin{verbatim}
%U\iz{R}=I \cdot R = 10\einh{A} \cdot 1\einh{$\Omega$}=10\einh{V}
%\end{verbatim}
%ergibt sich die folgende Ausgabe:
%\begin{align}
%U\iz{R}=I \cdot R = 10\einh{A} \cdot 1\einh{$\Omega$}=10\einh{V}
%\end{align}
%Weitere Funktionen k�nnen der Datei entnommen und individuell erweitert werden.
%
%\section{Quellenangabe}
%\label{kap:quellenangabe}
%Die Angabe von Quellen ist mit folgendem Code als Beispiel m�glich:
%\begin{verbatim}
%\cite{schwab}
%\end{verbatim}
%Die Ausgabe wird an der Stelle eingef�gt, an der man es einsetzt. Zum Beispiel hier: \cite{schwab}. Die Zahlen werden nach dem Vorkommen im Text vergeben und durchnummeriert. D.h. die n�chste Angabe erh�lt die Nummer zwei und zwar hier: \cite{oedingoswald}. Die Eintr�ge in der Literaturdatenbank k�nnen manuell in der Datei \textit{literatur.bib} oder mit der Software \textit{JabRef} angepasst werden.														% Hauptteil (Bezeichnung passend zur Arbeit!)

\chapter{Zusammenfassung und Ausblick}
\label{kap:zusammenfassung}

In der Zusammenfassung werden die Ergebnisse der Arbeit kurz zusammengefasst. Der
Umfang betr�gt ca. eine Seite.											% Zusammenfassung (Bezeichnung passend zur Arbeit!)

\appendix																							% Anhang

\chapter{�berblick einiger Verfahren}
\section{Einschrittverfahren}
\begin{sideways}\centering
\begin{tabular}{|lcccc|}\hline
\textbf{Name}&\textbf{Ordnung}&\textbf{Butcher-Schema}&\textbf{Art}&\textbf{Stabilit�tsfunktion}\\ \hline
Euler &1 & $\begin{array}[c]{c|c}
 0 &  \\ \hline
 & 1
\end{array}$ & explizit & ...\\ \hdashline[1pt/2pt]
%%
Heun & 2 & $\begin{array}[c]{c|cc}
 0   \\
1 & 1 \\ \hline
 & \eh & \eh \end{array}$ & explizit & ... \\\hdashline[1pt/2pt]
%%%
klassisches Runge-Kutta & 4 & $\begin{array}[c]{c|cccc}
		0 &\\ 
		\eh&\eh \\ 
		\eh & 0 & \eh \\
		1&0&0&1 \\ \hline
		& \frac{1}{6}&\frac{2}{6}&\frac{2}{6}&\frac{1}{6}
	\end{array}$ & explizit & ...\\\hdashline[1pt/2pt]
%%%%%%%%%%%
Euler & 1 & $\begin{array}[c]{c|c}
		0 &\\ \hline
		&1
	\end{array}$ & implizit & ...\\\hdashline[1pt/2pt]
Mittelpunktsregel & 2 & $\begin{array}[c]{c|cc}
		0 &\\ 
		1&1 \\ \hline
		& \eh & \eh
	\end{array}$ & implizit & ... \\\hdashline[1pt/2pt]
Trapezregel & 2 &$\begin{array}[c]{c|cccc}
		0 &\\ 
		\eh&\eh \\ 
		\eh & 0 & \eh \\
		1&0&0&1 \\ \hline
		& \frac{1}{6}&\frac{2}{6}&\frac{2}{6}&\frac{1}{6}
	\end{array}$ & implizit & ... \\ \lasthline
\end{tabular}
\end{sideways}
%%%%%%%%%%%%%%%%%%%%%%%%%%%%%%%%%%%%%%%%%%%%%%
\newpage
\section{Eingebettete Runge-Kutta-Verfahren (Schrittweitensteuerung)}
\label{sec:anheingeb}
\begin{longtable}{|lcc|}\hline
\tikzmark[2pt]{1}{\parbox{2cm}{\textbf{Name/\\Ordnung}}}&\textbf{Stufen}&\textbf{Butcher-Schema}\\ \hline
\parbox{2cm}{Fehlberg-RK4(3)} & 5 & $
\begin{array}[c]{c|ccccc}
		0\\
		\eh & \eh \\
		\eh & 0 & \eh \\
		1 & 0 & 0 & 1 \\%
		1 & \frac{1}{6} & \frac{2}{6} & \frac{2}{6} & \frac{1}{6}\\\hline
		& \frac{1}{6} & \frac{2}{6} & \frac{2}{6} & \frac{1}{6} & 0 \\ \hline
		& \frac{1}{6} & \frac{2}{6} & \frac{2}{6} & 0 & \frac{1}{6} 
\end{array}
$ \\ \hdashline[1pt/2pt]
\parbox{2cm}{Dormand-Prince-RK5(4)} & 6 & $
\begin{array}[c]{c|ccccccc}
		0\\
		\frac{1}{5} & \frac{1}{5} \\
		\frac{3}{10} & \frac{3}{40} & \frac{9}{40} \\
		\frac{4}{5} & \frac{44}{45} & -\frac{56}{15} & \frac{32}{9} \\%
		\frac{8}{9} & \frac{19372}{6561} & -\frac{25360}{2187} & \frac{64448}{6561} & -\frac{5103}{18656}\\
		1 & \frac{9017}{3168} & -\frac{355}{33} & \frac{46732}{5247} & \frac{49}{176} & -\frac{5103}{18656}\\
		1 & \frac{35}{384} & 0 & \frac{500}{1113} & \frac{125}{92} & -\frac{2187}{6487} & \frac{11}{84}\\\hline
		& \frac{35}{384} & 0 & \frac{500}{1113} & \frac{125}{92} & -\frac{2187}{6487} & \frac{11}{84}&0 \\ \hline
		& \frac{5179}{57600} & 0 & \frac{7571}{16695} & \frac{393}{640} & -\frac{92097}{339200} & \frac{187}{2100} & \frac{1}{40} 
\end{array}
$ \\ \lasthline
\end{longtable}
%%%%%%%%%%%%%%%%%%%%%%%%%%%%%%%%%%%%%%%%%
\begin{landscape}
\section{Mehrschrittverfahren}
\label{sec:anhmehr}
\small
\arraycolsep0pt
\begin{longtable}{|llcccc|}\hline
$k$ &\textbf{Name}&\textbf{$P_a(z)$}&\textbf{$P_b(z)$}&\textbf{Art}&\textbf{Stabilit�t}\\ \hline
$\begin{array}{c}2\end{array}$ & \parbox{1.4cm}{Milne-Simpson} &  $z^2-1$ & $\frac{1}{3}z^2 + \frac{4}{3}z + \frac{1}{3}$ & implizit & ...? (Sinn?) \\[4pt] \hdashline[1pt/2pt]
$\begin{array}[c]{l}
1\\2\\3\\4
\end{array}$&
\parbox{1.4cm}{Adams-Bashforth} & $z^{k}-z^{k-1}$ &$
\begin{array}[c]{l}
	1 \ \text{\footnotesize($\hat{=}$ Euler, explizit)}\\
	(3z-1)/2 \\
	(23z^2-16z+5)/12\\
	(55z^3-59z^2+37z-9)/24
\end{array}$ & explizit & ... \\\hdashline[1pt/2pt]
$\begin{array}[c]{l}
2\\3\\4
\end{array}$&
\parbox{1.4cm}{Adams-Moulton} & $z^{k}-z^{k-1}$ &$
\begin{array}[c]{l}
	(5z^2+8z-1)/12\\
	(9z^3+19z^2-5z+1)/24\\
	(251z^4+646z^3+-246z^2+106z-19)
\end{array}$ & implizit & ... \\[4pt]\hdashline[1pt/2pt]
$\begin{array}[c]{l}
1\\2\\3\\4
\end{array}$&
BDF &%\makebox[2.5cm][l]{$\begin{array}[c]{l}
	$\begin{array}[c]{l}
	z-1 \ \text{\footnotesize($\hat{=}$ Euler, implizit)}\\
	\frac{3}{2}z^2-2z+\eh\\
	\frac{11}{6}z^3-3z^2+\frac{3}{2}z-\frac{1}{3}\\
	\frac{25}{12}z^4-4z^3+3z^2-\frac{4}{3}z+\frac{1}{4}\\
\end{array}$ & $z^k$ & implizit & \parbox{4cm}{unendlich gro�es Stabilit�tsgebiet $\Rightarrow$ gut f�r steife DGL} \\\lasthline
\end{longtable}
\end{landscape}
\chapter{Anhang Teil 2}
\chapter{Anhang Teil 3}























																% Anhangsinhalt

\addchap{Symbol- und Abk�rzungsverzeichnis}
\label{kap:symbole}

Sollten in einer Ausarbeitung viele Abk�rzungen und Formelzeichen auftreten, so
empfiehlt es sich, diese gesondert in einem Kapitel aufzufuhren. Dieses kann auch nach
dem Inhaltsverzeichnis (Abbildungs- und Tabellenverzeichnis) folgen.
\begin{tabbing}
XXXXXX \= XXXX \= XXXXXXXXXXXXXXXXXXXXXXXXXXXXXXXXXXXXXXXXXXXXXXXXX \kill
A											\>																	\> Abk�rzung f�r ... \\
$\cos\varphi$					\>																	\> Leistungsfaktor \\
$U\iz{s}$							\> V																\> Betrag der Statorspannung \\
$\varphi$							\> rad															\> Winkel zwischen Spannung und Strom \\
\end{tabbing}															% Symbol- und Abk�rzungsverzeichnis

\begin{spacing}{1.5}          												% evtl. kleinerer Zeilenabstand im LV
%\nocite{*}                              							% alle Literaturquellen einbinden, sonst werden nur die zitierten
                                            					% Quellen im Literaturverzeichnis angezeigt.
\bibliographystyle{unsrtdin}		         						  % Darstellung numerisch nach Vorkommen im Text
\bibliography{literatur}
\end{spacing}

%\newpage
%\renewcommand{\indexname}{Stichwortverzeichnis}			% Stichwortverzeichnisbezeichnung
%\addcontentsline{toc}{chapter}{Stichwortverzeichnis} % Stichwortverzeichnis ins IV einf?gen
%\begin{spacing}{1.0} 
%\printindex																					% Stichwortverzeichnis anzeigen
%\end{spacing}

\end{document}
%------------------------- Ende des Dokuments -----------------------